
\documentclass[11pt]{amsart}

\usepackage{amsmath}
\usepackage{amssymb}
\usepackage{amsthm}
\usepackage{mathtools}
\usepackage{enumitem}
%\usepackage{thmtools}
\usepackage{hyperref}
\hypersetup{hypertexnames=false}
\usepackage[nameinlink,noabbrev]{cleveref}
\usepackage{booktabs}
\usepackage{graphicx}
\usepackage{xcolor}


% ---- Theorem environments ----
\theoremstyle{plain}
\newtheorem{theorem}{Theorem}
\newtheorem{proposition}[theorem]{Proposition}
\newtheorem{lemma}[theorem]{Lemma}
\newtheorem{corollary}[theorem]{Corollary}
\newtheorem{claim}[theorem]{Claim}

\theoremstyle{definition}
\newtheorem{definition}[theorem]{Definition}
\newtheorem{assumption}[theorem]{Assumption}
\newtheorem{example}[theorem]{Example}

\theoremstyle{remark}
\newtheorem{remark}[theorem]{Remark}

% ---- Macros ----
\newcommand{\cond}{\operatorname{cond}}
\newcommand{\KL}{D_{\operatorname{KL}}}
\newcommand{\E}{\mathbb{E}}
\newcommand{\Var}{\operatorname{Var}}
\newcommand{\Cov}{\operatorname{Cov}}
\newcommand{\tr}{\operatorname{Tr}}
\newcommand{\diag}{\operatorname{diag}}

\newcommand{\Range}{\operatorname{Range}}
\newcommand{\Span}{\operatorname{span}}
\newcommand{\id}{\mathrm{id}}
% \newcommand{\TODO}[1]{\textcolor{red}{[TODO: #1]}} % removed for submission

% ---- Cleveref names ----
\crefname{theorem}{theorem}{theorems}
\Crefname{theorem}{Theorem}{Theorems}
\crefname{proposition}{proposition}{propositions}
\Crefname{proposition}{Proposition}{Propositions}
\crefname{lemma}{lemma}{lemmas}
\Crefname{lemma}{Lemma}{Lemmas}
\crefname{corollary}{corollary}{corollaries}
\Crefname{corollary}{Corollary}{Corollaries}
\crefname{claim}{claim}{claims}
\Crefname{claim}{Claim}{Claims}
\crefname{definition}{definition}{definitions}
\Crefname{definition}{Definition}{Definitions}
\crefname{assumption}{assumption}{assumptions}
\Crefname{assumption}{Assumption}{Assumptions}
\crefname{example}{example}{examples}
\Crefname{example}{Example}{Examples}
\crefname{remark}{remark}{remarks}
\Crefname{remark}{Remark}{Remarks}
\crefname{equation}{Eq.}{Eqs.}
\Crefname{equation}{Equation}{Equations}
\crefname{figure}{fig.}{figs.}
\Crefname{figure}{Fig.}{Figs.}
\crefname{table}{table}{tables}
\Crefname{table}{Table}{Tables}
\crefname{section}{Sec.}{Secs.}
\Crefname{section}{Section}{Sections}

\title[Singular Gaussian Conditioning in Hilbert Space]{Singular Gaussian Conditioning in Hilbert Space: Inverse-Free Schur Projection and Variance Collapse}

\author{Lucia Shang}
\thanks{Manuscript received May 2026. Repository information omitted for submission.}

\begin{document}

% ============================================================
%  ABSTRACT
% ============================================================
\begin{abstract}
We study singular coordinate conditioning for Gaussian measures on separable Hilbert spaces. Since trace-class covariance blocks are compact, the usual Schur-complement formula cannot be implemented through a bounded inverse. We introduce an inverse-free variational Schur projection on the Cameron--Martin range and prove canonical normalization together with an $O(\lambda^{-6})$ residual estimate. The construction is stable under Galerkin truncation for exponentially stable analytic semigroup dynamics and applies to sectorial elliptic Gaussian fields.
\end{abstract}

\keywords{Singular conditioning; Gaussian measures; Cameron--Martin space; Schur complement; Hilbert-space covariance operators; trace-class covariance; Radon Gaussian conditioning; Galerkin approximation; stochastic evolution equations}

\maketitle

% ============================================================
%  \S1  INTRODUCTION
% ============================================================
\section{Introduction}
\label{sec:intro}

Singular coordinate conditioning arises when an exact constraint $X_k=x_k^*$ is reached as a zero-noise or infinite-precision limit of Gaussian conditioning problems.  Finite-dimensional Gaussian conditioning is governed by the ordinary Schur complement of a covariance matrix \cite{horn2012matrix,rasmussen2006gaussian}.  Infinite-dimensional Gaussian conditioning is subtler: Radon Gaussian measures on separable Hilbert or Banach spaces have compact covariance operators, and their Cameron--Martin spaces are covariance-norm domains rather than ambient statistical parameter spaces \cite{bogachev1998gaussian,daprato2014stochastic,steinwart2024conditioning}.  In Hilbert-space Gaussian conditioning, shorted operators give the covariance of conditioning on closed subspaces, while recent inverse-problem work emphasizes that positive-noise posterior formulas must be interpreted through approximation, disintegration, or Hilbert-space posterior equivalence rather than through a global covariance inverse \cite{owhadi2015conditioning,stuart2010inverse,chen2024gaussian,carerelie2024optimal,carerelie2025posterior}.

The obstruction studied here is the covariance-block obstruction behind that phenomenon.  In a finite-dimensional model, the conditional variance of a coordinate can be written as $A-BD^{-1}B^*$ after a block decomposition.  In a separable Hilbert space the corresponding block $D$ is typically compact and trace class, hence $D^{-1}$ is not a bounded operator on the ambient space \cite{bogachev1998gaussian,simon2005trace,daprato2014stochastic}.  Classical Schur-complement and shorted-operator theory, originating in Krein's theory of non-negative extensions and developed systematically by Anderson and Trapp, provides variational tools for positive block operators \cite{krein1947selfadjoint,anderson1975shorted,douglas1966majorization}.  However, this theory does not by itself give the scaling law for a singular Gaussian conditioning residual \cite{tretter2008spectral,contino2018schur,friedrich2017generalized,hassi2023complementation}.  This paper formulates the residual directly as a Cameron--Martin norm projection and never invokes a pseudoinverse or an ambient precision operator.

The hard-conditioning limit also reaches the Feldman--H\'{a}jek measure-class boundary.  With positive observational noise, Gaussian Bayesian inverse problems may remain in the equivalent-measure regime \cite{stuart2010inverse,carerelie2024optimal,carerelie2025posterior}.  Under an exact coordinate constraint, the limiting conditional law is supported on a hyperplane of prior measure zero, so the equivalent-measure regime is left \cite{feldman1958equivalence,hajek1958property,bogachev1998gaussian}.  A useful theory must therefore separate the diverging joint divergence caused by exact evidence from the finite covariance residual that remains visible in the block structure.

We prove that this finite residual is controlled by an inverse-free variational Schur projection.  Under the Lyapunov block assumptions stated below, the target residual has the exact form
\[
  S_{k,\lambda}=\frac{\varepsilon_0}{2\lambda^2}-R_\lambda,
  \qquad
  R_\lambda=O(\lambda^{-6}).
\]
Consequently,
\[
  S_{k,\lambda}^{-1}=\frac{2}{\varepsilon_0}\lambda^2+O(\lambda^{-2}).
\]
The proof uses covariance-norm forms, shorting, Douglas range inclusion, Lyapunov block identities, and trace-ideal convergence. These are standard inputs in operator and SPDE analysis \cite{krein1947selfadjoint,anderson1975shorted,douglas1966majorization} \cite{simon2005trace,daprato2014stochastic}. The Cameron--Martin space is used only as a Hilbert subspace equipped with its covariance norm; no differentiable statistical or geometric structure is assumed \cite{bogachev1998gaussian}.
The operator base is an exponentially stable analytic semigroup with bounded drift perturbations, so dissipative sectorial generators such as $-(-\Delta+\kappa^2)$ fall inside the abstract setup.

The paper has four main contributions:
\begin{enumerate}
  \item \textbf{Canonical normalization.} We identify the block-Lyapunov mechanism that fixes the observed coordinate variance at order $\lambda^{-2}$ with explicit coefficient $\varepsilon_0/2$.
  \item \textbf{Inverse-free Schur projection.} We identify the residual subtraction with the target-coordinate quadratic form of the shorted-operator variational formula, then prove the $O(\lambda^{-6})$ residual estimate in the same operator-theoretic regime where compact covariance blocks preclude a bounded Schur inverse.
  \item \textbf{Analytic-semigroup and Galerkin stability.} We work with an analytic semigroup drift base plus bounded perturbations, and prove that finite-dimensional truncations converge to the Hilbert-space residual in trace ideal norm, matching the discretization-independent perspective used in Hilbert-space Gaussian inverse problems \cite{stuart2010inverse,carerelie2024optimal,carerelie2025posterior}.
  \item \textbf{Radon Gaussian conditioning.} For Radon Gaussian laws, we construct the conditional predictor as a measurable Cameron--Martin projection and show that the associated renormalized conditioning functional remains bounded, using regular conditional probabilities and continuous Gaussian disintegration \cite{bogachev1998gaussian,lagatta2013continuous,steinwart2024conditioning}.
\end{enumerate}

\subsection{Positioning in the literature}
\label{sec:related}

The paper is closest to four bodies of work.  First, it uses the classical measure-class structure of Gaussian measures, especially the Cameron--Martin theorem and the Feldman--H\'{a}jek dichotomy, as presented in standard references on Gaussian measures \cite{feldman1958equivalence,hajek1958property,bogachev1998gaussian}.  Second, it is adjacent to Bayesian inverse problems and Hilbert-space Gaussian conditioning, including the shorted-operator description of Gaussian conditioning on closed subspaces, positive-noise posterior formulas, low-rank posterior approximations, and nonlinear observation maps \cite{owhadi2015conditioning,stuart2010inverse,chen2024gaussian,steinwart2024conditioning,carerelie2024optimal,carerelie2025posterior}.  Third, the variational residual is informed by Schur complements, shorted operators, Douglas range inclusion, and modern complementability theory for block operators \cite{krein1947selfadjoint,anderson1975shorted,douglas1966majorization,tretter2008spectral,contino2018schur,friedrich2017generalized,hassi2023complementation}.  Fourth, the motivating examples come from stochastic evolution equations and elliptic Gaussian fields, where dissipative sectorial generators produce compact trace-class stationary covariances that are naturally approximated by Galerkin truncation \cite{daprato2014stochastic,lord2014introduction,simon2005trace}.

\subsection{Organization of the paper}
\label{sec:organization}

\Cref{sec:setup} defines the operator model and establishes the canonical normalization. \Cref{sec:infdim} establishes the Hilbert-space variational Schur residual. \Cref{sec:gaussian_zero} gives the Radon Gaussian conditional extension and the renormalized conditioning functional. \Cref{sec:spde} gives an elliptic Gaussian-field example. \Cref{sec:discussion} records limitations and extensions, and \Cref{sec:conclusion} concludes. Proofs are given in \Cref{app:block_cumulant_analysis,app:gaussian_lift}; numerical diagnostics and Galerkin verification are given in \Cref{app:experiments}.

\section{Operator Model and Canonical Normalization}
\label{sec:setup}

\subsection{Operator Dynamics and Regularized Conditioning}
\label{sec:setup:dynamics}

We model the continuous-time dynamics as stochastic evolution on a separable Hilbert space $\mathcal{H}$ \cite{daprato2014stochastic}:
\begin{equation}
  \label{eq:stochastic_evolution}
  dX_t = \mathcal{A} X_t \, dt + \mathcal{D}^{1/2} dW_t
\end{equation}
where $\mathcal{A}$ is the drift operator and $\mathcal{D}$ is the trace-class diffusion operator. We decompose the Hilbert space as $\mathcal{H} = \mathcal{H}_k \oplus \mathcal{H}_I$, where $\mathcal{H}_k \coloneqq \mathbb{R}\phi_k$ is the one-dimensional target coordinate subspace, and $\mathcal{H}_I \coloneqq \mathcal{H}_k^\perp$ is the free bulk subspace. Let $\Pi_k \coloneqq \phi_k \otimes \phi_k$ and $\Pi_I \coloneqq I - \Pi_k$ be the corresponding orthogonal projections.

To approximate the exact point conditioning $X_k = x_k^*$, we set the incoming target-coordinate blocks to zero and apply a regularization parameter $\lambda > 0$.

\begin{definition}[Calibrated Approximating Dynamics Family]
  \label{def:soft_do}
  The regularized conditioning operator of strength $\lambda > 0$ at coordinate $X_k$ acts on the operator dynamics \eqref{eq:stochastic_evolution} by modifying both the drift and the diffusion:
  \begin{equation}
    \label{eq:soft_do_sde}
    \mathcal{A}^{(\lambda)} \coloneqq \mathcal{A}_{\cond} - \lambda \Pi_k, \qquad
    \mathcal{D}^{(\lambda)} \coloneqq \mathcal{D}_{\cond} + d_\lambda \Pi_k
  \end{equation}
  where $\mathcal{A}_{\cond}$ is the block drift operator with incoming target-coordinate blocks set to zero ($\Pi_k \mathcal{A}_{\cond} = 0$), $\mathcal{D}_{\cond}$ has target-coordinate cross-noise set to zero, and $d_\lambda > 0$ is the local diffusion variance of the target coordinate at regularization strength $\lambda$. In the SDE, this corresponds to:
  \begin{equation}
    \label{eq:soft_do_sde_general}
    dX_{k,t} = -\lambda (X_{k,t} - x_k^*) dt + \sqrt{d_\lambda}\, dW_{k,t}.
  \end{equation}
\end{definition}

\subsection{Canonical Normalization}
\label{sec:setup:calibration}

The block-decoupled structure under the regularized dynamics yields:
\begin{equation}
  \mathcal{A}^{(\lambda)}=
  \begin{pmatrix}-\lambda&0\\a_{Ik}&\mathcal{A}_{II}\end{pmatrix},\qquad
  \mathcal{D}^{(\lambda)}=
  \begin{pmatrix}d_\lambda&0\\0&D_{II}\end{pmatrix},
\end{equation}
where $a_{Ik} \coloneqq \mathcal{A}_{Ik} \phi_k \in \mathcal{H}_I$ represents the off-diagonal drift coupling.
Let $\Sigma_\lambda$ denote the stationary covariance operator under regularization strength $\lambda$, solving the Lyapunov equation:
\begin{equation}
  \mathcal{A}^{(\lambda)}\Sigma_\lambda + \Sigma_\lambda(\mathcal{A}^{(\lambda)})^* + \mathcal{D}^{(\lambda)} = 0.
\end{equation}
The target block $kk$ decouples completely from the free coordinates, yielding the exact scalar equation:
\begin{equation}
  -2\lambda\Sigma_{\lambda,kk} + d_\lambda = 0 \quad \Longrightarrow \quad \Sigma_{\lambda,kk} = \frac{d_\lambda}{2\lambda}.
\end{equation}

\begin{proposition}[Canonical Scaling]
  \label{prop:exact_calibration}
  To maintain a target noise scale $\varepsilon_0 > 0$ independent of the auxiliary parameter $\lambda$ under the singular limit, we define the canonical scaling condition:
  \begin{equation}
    \lambda^2 \Sigma_{\lambda,kk} \equiv \frac{\varepsilon_0}{2} \quad \text{for all } \lambda > 0,
  \end{equation}
  which uniquely requires:
  \begin{equation}
    d_\lambda = \frac{\varepsilon_0}{\lambda} \quad \text{for every } \lambda > 0.
  \end{equation}
  Substitution of this target diffusion scaling yields the calibrated approximating dynamics:
  \begin{equation}
    dX_{k,t} = -\lambda (X_{k,t} - x_k^*) dt + \sqrt{\varepsilon_0 / \lambda}\, dW_{k,t}.
  \end{equation}
\end{proposition}

All subsequent results are stated for this calibrated family.

\begin{definition}[Schur LMMSE Invariant]
  \label{def:schur_invariant}
  Let $S_{k,\lambda} \coloneqq \|X_k - P_{\mathcal{M}_I} X_k\|_{L^2}^2$ represent the linear minimum mean squared error (LMMSE) residual of the target coordinate $X_k$ on the closed linear span $\mathcal{M}_I$ of the free coordinates $X_I$. The \emph{Schur LMMSE invariant} $m_2^{\mathrm{Schur}}$ is defined as:
  \begin{equation}
    \label{eq:m2_schur_definition}
    m_2^{\mathrm{Schur}} \coloneqq \lim_{\lambda\to\infty} \lambda^2 S_{k,\lambda}.
  \end{equation}
\end{definition}

\subsection{Operator Setup and Assumptions}
\label{sec:setup:assumptions}

We state the operator-theoretic assumptions that make the limit passage independent of the Galerkin dimension.

\begin{definition}[Drift Operator with Analytic Base]
  \label{def:drift_operator}
  A drift operator on a separable Hilbert space $\mathcal{H}$ is an operator of the form $\mathcal{A} = \mathcal{A}_0 + \mathcal{K}$, where:
  \begin{enumerate}
    \item[(i)] $\mathcal{A}_0: \operatorname{dom}(\mathcal{A}_0) \subset \mathcal{H} \to \mathcal{H}$ is closed, densely defined, and dissipative, and it generates a strongly continuous analytic contraction semigroup $(e^{t\mathcal{A}_0})_{t \ge 0}$ on $\mathcal{H}$ satisfying $\|e^{t\mathcal{A}_0}\| \le e^{-\alpha t}$ for some spectral gap $\alpha > 0$; equivalently, the base operator $-\mathcal{A}_0$ is sectorial of angle $<\pi/2$ in this exponentially stable contraction setting.
    \item[(ii)] $\mathcal{K}: \mathcal{H} \to \mathcal{H}$ is a bounded linear operator.
    \item[(iii)] The full operator $\mathcal{A} = \mathcal{A}_0 + \mathcal{K}$, with $\operatorname{dom}(\mathcal{A})=\operatorname{dom}(\mathcal{A}_0)$, generates a strongly continuous analytic semigroup $(e^{t\mathcal{A}})_{t \ge 0}$ by the bounded perturbation theorem, and this semigroup remains exponentially stable: $\|e^{t\mathcal{A}}\| \le M e^{-\beta t}$ for some $M \ge 1$, $\beta > 0$.
  \end{enumerate}
  In finite dimensions ($\mathcal{H} = \mathbb{R}^n$), this reduces to a decomposition of a Hurwitz drift matrix into an exponentially stable analytic base plus a bounded perturbation.
\end{definition}

\begin{definition}[Trace-Class Noise and Covariance]
  \label{def:trace_class}
  Let $\mathcal{S}_1(\mathcal{H})$ denote the Banach space of trace-class (nuclear) operators on $\mathcal{H}$ equipped with the trace norm $\|T\|_{\mathcal{S}_1} \coloneqq \tr(|T|) = \sum_{i=1}^\infty \langle \phi_i, (T^*T)^{1/2} \phi_i \rangle < \infty$.
  \begin{enumerate}
    \item[(i)] The diffusion operator $\mathcal{D}: \mathcal{H} \to \mathcal{H}$ is positive, self-adjoint, and belongs to $\mathcal{S}_1(\mathcal{H})$.
    \item[(ii)] Assuming the pair $(\mathcal{A}, \mathcal{D}^{1/2})$ is controllable, the stationary covariance operator $\Sigma: \mathcal{H} \to \mathcal{H}$ is the unique strictly positive, self-adjoint, trace-class solution to the operator Lyapunov equation:
    \begin{equation}
      \label{eq:op_lyapunov}
      \mathcal{A}\Sigma + \Sigma\mathcal{A}^* + \mathcal{D} = 0
    \end{equation}
    which is given by the convergent Bochner integral
    \[
      \Sigma = \int_0^\infty e^{t\mathcal{A}} \mathcal{D} e^{t\mathcal{A}^*} dt.
    \]
  \end{enumerate}
\end{definition}

\begin{assumption}[Core Analytical Assumptions for Singular Hyperplane Conditioning]
  This work relies on two core analytical assumptions regarding the target-coordinate block decoupling and the outgoing coupling:
  \begin{enumerate}
    \item[(i)] \textbf{Linear Second-Moment Channel:} \label{assum:linear_second_moment} A centered linear stationary process on $\mathcal{H}=\mathcal{H}_k\oplus \mathcal{H}_I$ has finite second moments and covariance solving \Cref{eq:op_lyapunov}. The block-decoupling constraint zeroes all incoming target drift and cross-noise, so that $\mathcal{A}_{kI}=\mathcal{D}_{kI}=\mathcal{D}_{Ik}=0$ and $\mathcal{A}_{kk}=-\lambda$. The target diffusion $\mathcal{D}_{kk}=\varepsilon_0/\lambda$ is the unique scaling determined by the canonical normalization (Proposition~\ref{prop:exact_calibration}). The free noise is independent of the local target noise. (Gaussianity is not assumed).
    \item[(ii)] \textbf{Cameron--Martin Stability of Off-Diagonal Drift Coupling:} \label{assum:cm_stability} Let $Q_0\coloneqq\Sigma_{II}^{(0)}$ denote the free-block covariance produced by the free noise under block decoupling, and let $H_{Q_0}\coloneqq\operatorname{Range}(Q_0^{1/2})$ with the covariance norm $|u|_{Q_0}\coloneqq\|Q_0^{-1/2}u\|_{\mathcal{H}_I}$. The outgoing vector $a_{Ik}\coloneqq\mathcal{A}_{Ik}\phi_k$ belongs to $H_{Q_0}$, and $e^{t\mathcal{A}_{II}}$ restricts to a strongly continuous exponentially stable semigroup on $H_{Q_0}$:
    \begin{equation}
      \|e^{t\mathcal{A}_{II}}\|_{\mathcal{L}(H_{Q_0})} \le M e^{-\omega t},\qquad t\ge0,
    \end{equation}
    for constants $M\ge1$ and $\omega>0$ independent of $\lambda$.
  \end{enumerate}
  \vspace{1mm}
  \noindent\textit{Intuitive Explanation:} Part (i) states that the regularized dynamics act as a strong localized drift penalization on the target coordinate $X_k$ while isolating it from incoming influences, giving an operator implementation of the block-decoupling constraint. Part (ii) requires that the off-diagonal drift coupling $a_{Ik}$ from the target coordinate has finite covariance norm relative to the background fluctuations of the free coordinates. This ensures that target fluctuations propagating downstream do not overload the variance of the infinite-dimensional system.
\end{assumption}

\begin{remark}[Finite-Norm Outgoing Covariance Coupling and Finite-Dimensional Calibration]
  \label{rem:cm_information_flow}
  The membership $a_{Ik}\in H_{Q_0}$ is an admissibility condition,
  not a consequence of the free-block Lyapunov equation alone: an arbitrary
  outgoing vector need not lie in
  $\operatorname{Range}((\Sigma_{II}^{(0)})^{1/2})$.  It says that the
  deterministic channel carrying target fluctuations into the free-block has
  finite covariance norm relative to the free-noise background.
  
  To understand this condition intuitively, consider a finite-dimensional setting where the free-block state space is $\mathcal{H}_I = \mathbb{R}^{d-1}$ and the free stationary covariance $Q_0 \coloneqq \Sigma_{II}^{(0)}$ is strictly positive-definite. Since $Q_0 \succ 0$, its square root $Q_0^{1/2}$ is invertible, and its range is the entire space $\mathbb{R}^{d-1}$. In this case, the Cameron--Martin space is the full space $H_{Q_0} = \operatorname{Range}(Q_0^{1/2}) = \mathbb{R}^{d-1}$, and the Cameron--Martin norm $|u|_{Q_0} = \|Q_0^{-1/2} u\|_2$ is equivalent to the standard Euclidean norm. Consequently, the admissibility condition $a_{Ik} \in H_{Q_0}$ is trivially satisfied for any outgoing vector $a_{Ik} \in \mathbb{R}^{d-1}$.
  
  However, if some direction in the free coordinates is noise-free under block decoupling (meaning $Q_0$ is positive semidefinite but singular), $H_{Q_0}$ becomes a proper subspace of $\mathbb{R}^{d-1}$. In this case, the condition $a_{Ik} \in H_{Q_0}$ requires that the leakage vector $a_{Ik}$ has no component in the null space of $Q_0$. Equivalently, it prevents the target coordinate from leaking fluctuations into a downstream coordinate that has zero background noise, which would otherwise result in an infinite relative covariance norm.
  
  In infinite dimensions, $\Sigma_{II}^{(0)}$ is trace-class and compact, so its eigenvalues decay to zero, making $\operatorname{Range}((\Sigma_{II}^{(0)})^{1/2})$ a dense but strictly proper subspace of $\mathcal{H}_I$. The condition $a_{Ik} \in H_{Q_0}$ requires the off-diagonal drift coupling coefficients to decay sufficiently fast relative to the background noise eigenvalues to keep the relative covariance norm finite.
  
  Under this assumption, invariance and exponential stability on $H_{Q_0}$ imply:
  \begin{equation}
  \begin{aligned}
    (\lambda I-\mathcal{A}_{II})^{-1}a_{Ik}
       &\in H_{Q_0},\\
    |(\lambda I-\mathcal{A}_{II})^{-1}a_{Ik}|_{Q_0}
       &\le \frac{M}{\lambda+\omega}|a_{Ik}|_{Q_0}.
  \end{aligned}
  \end{equation}
\end{remark}

\begin{remark}[Constant Dependency in LMMSE Expansion]
  Consequently, the exact leakage cross-covariance
  $\Sigma_{\lambda,Ik}$ has finite Cameron--Martin norm and is
  $O(\lambda^{-3})$ in that norm.  This is a condition on the deterministic
  response in the covariance norm; it does not assert that infinite-dimensional
  Gaussian sample paths lie in $H_{Q_0}$ almost surely.
\end{remark}

% ============================================================
%  §3  HILBERT-SPACE VARIATIONAL SCHUR RESIDUAL
% ============================================================
\section{Hilbert-Space Variational Schur Residual}
\label{sec:infdim}

\subsection{Variational Schur Subtraction and Norm Bounds}
\label{sec:infdim:variational}

In finite dimensions, the Schur complement of the bulk covariance block $\Sigma_{II}$ is defined as $S_k \coloneqq \Sigma_{kk} - \Sigma_{kI} \Sigma_{II}^{-1} \Sigma_{Ik}$. Under stable and controllable dynamics, the bulk covariance $\Sigma_{II}$ is positive-definite and invertible, so the subtraction is well-defined. However, in infinite dimensions, the stationary covariance operator is trace-class and compact, meaning its restriction $\Sigma_{\lambda,II}$ lacks a bounded global inverse. The appropriate operator-theoretic replacement is the shorted operator: for a non-negative block operator, the Schur complement is recovered as the quadratic form of the short to the target subspace, with the inverse term replaced by a variational supremum \cite{krein1947selfadjoint,anderson1975shorted}. To lift the Schur complement to infinite dimensions in the present singular-conditioning problem, we formulate the projection residual directly using covariance-norm forms.

Let $Q_0 \coloneqq \Sigma_{II}^{(0)}$ denote the fixed background bulk covariance under full decoupling ($\lambda = \infty$). We define the core Cameron--Martin space $H_{Q_0} \coloneqq \operatorname{Range}(Q_0^{1/2})$ and the core covariance norm:
\begin{equation}
  |u|_{Q_0} \coloneqq \|Q_0^{-1/2} u\| \quad \text{for all } u \in H_{Q_0}.
\end{equation}
By Assumption~\ref{assum:cm_stability}, the off-diagonal drift coupling vector satisfies $a_{Ik} \in H_{Q_0}$, and the restricted drift generator $\mathcal{A}_{II}$ generates an exponentially stable semigroup on $H_{Q_0}$.

Under this assumption, the exact cross-covariance resolvent formula (derived in \Cref{app:block_cumulant_analysis}) satisfies:
\begin{equation}
  \Sigma_{\lambda,Ik} = \frac{\varepsilon_0}{2\lambda^2} (\lambda I - \mathcal{A}_{II})^{-1} a_{Ik}.
\end{equation}
This resolvent response yields the Cameron--Martin norm bound on the cross-covariance block:
\begin{equation}
  \label{eq:cross_norm_bound}
  |\Sigma_{\lambda,Ik}|_{Q_0} = O(\lambda^{-3}) \quad \text{as } \lambda \to \infty.
\end{equation}

To define the infinite-dimensional projection residual, we introduce the variational Schur subtraction $R_\lambda$:
\begin{equation}
  \label{eq:variational_schur_subtraction}
  R_\lambda \coloneqq \sup_{\langle u, \Sigma_{\lambda,II} u \rangle > 0} \frac{|\langle u, \Sigma_{\lambda,Ik} \rangle|^2}{\langle u, \Sigma_{\lambda,II} u \rangle}.
\end{equation}
Equivalently, if $\Sigma_\lambda$ is viewed as a non-negative block operator on $\mathbb{C}\phi_k\oplus\mathcal{H}_I$, then \eqref{eq:variational_schur_subtraction} is the one-dimensional Anderson--Trapp shorted-operator subtraction, namely
\begin{equation}
  \label{eq:shorted_operator_scalar}
  \big\langle \phi_k,\mathcal{S}_{\mathbb{C}\phi_k}(\Sigma_\lambda)\phi_k\big\rangle
  =\Sigma_{\lambda,kk}-R_\lambda.
\end{equation}
The restriction to vectors with $\langle u,\Sigma_{\lambda,II}u\rangle>0$ simply removes null directions of the denominator; for a non-negative block operator, such directions do not change the supremum. Thus $R_\lambda$ is not an ambient inverse computation but the covariance-norm contribution removed by shorting to the target coordinate \cite{anderson1975shorted,owhadi2015conditioning}.

Since the target noise channel and bulk noise channel are independent, the bulk state decomposes into independent background and target-driven components, meaning the bulk covariance dominates the background covariance ($\Sigma_{\lambda,II} = Q_0 + \operatorname{Cov}(Y_I) \succeq Q_0$). Douglas range inclusion then guarantees that the variational subtraction is bounded:
\begin{equation}
  0 \le R_\lambda \le |\Sigma_{\lambda,Ik}|_{Q_0}^2 = O(\lambda^{-6}).
\end{equation}

\subsection{Hilbert-Space Schur Residual Scaling}
\label{sec:infdim:scaling}

We now establish the main scaling theorem for the $L^2$-projection residual.

\begin{theorem}[Hilbert Schur Residual Invariant]
  \label{thm:inf_cancellation}
  Under Assumptions~\ref{assum:linear_second_moment} and \ref{assum:cm_stability}, the $L^2$-projection residual satisfies:
  \begin{equation}
    \label{eq:main_schur_sharp}
    S_{k,\lambda} \coloneqq \|X_k - P_{\mathcal{M}_I} X_k\|_{L^2}^2 = \frac{\varepsilon_0}{2\lambda^2} - R_\lambda,
  \end{equation}
  where the remainder satisfies $0 \le R_\lambda = O(\lambda^{-6})$. In particular, the Schur LMMSE invariant is:
  \begin{equation}
    \label{eq:m2_schur_value}
    m_2^{\mathrm{Schur}} = \lim_{\lambda\to\infty} \lambda^2 S_{k,\lambda} = \frac{\varepsilon_0}{2},
  \end{equation}
  and the inverse projection residual scaling satisfies:
  \begin{equation}
    \label{eq:schur_scaling_inverse}
    S_{k,\lambda}^{-1} = \frac{2}{\varepsilon_0}\lambda^2 + O(\lambda^{-2}).
  \end{equation}
  This result depends only on covariance and orthogonal projection; it holds for any linear stationary model satisfying these assumptions, without Gaussianity.
\end{theorem}
\begin{proof}
  See \Cref{app:block_cumulant_analysis} for the full derivation.
\end{proof}

\subsection{Convergence of Finite-Dimensional Truncations}
\label{sec:proofs:convergence}

Let $\Pi_n: \mathcal{H} \to \mathcal{H}_n \coloneqq \operatorname{span}\{\phi_1, \ldots, \phi_n\}$ be the $n$-dimensional Galerkin projection. Write $\mathcal{A}_n^{(\lambda)} \coloneqq \Pi_n \mathcal{A}^{(\lambda)} \Pi_n$, $\mathcal{D}_n^{(\lambda)} \coloneqq \Pi_n \mathcal{D}^{(\lambda)} \Pi_n$, and let $\Sigma_{n,\lambda}$ solve the truncated Lyapunov equation $\mathcal{A}_n^{(\lambda)} \Sigma_{n,\lambda} + \Sigma_{n,\lambda} (\mathcal{A}_n^{(\lambda)})^* + \mathcal{D}_n^{(\lambda)} = 0$.

\begin{proposition}[Galerkin Consistency of the Schur Invariant]
  \label{prop:galerkin_consistency}
  Under Assumption~\ref{assum:linear_second_moment}, let $\Pi_n:\mathcal{H}\to\mathcal{H}_n$ be finite-rank orthogonal Galerkin projections preserving the target-bulk block decomposition:
  \begin{equation}
    \Pi_n\phi_k=\phi_k,\qquad \Pi_n\mathcal{H}_I\subset\mathcal{H}_I,\qquad \Pi_n\to I\ \text{strongly on }\mathcal{H}.
  \end{equation}
  Assume the truncated regularized dynamics satisfy the uniform exponential stability and trace-class conditions stated in Proposition~\ref{prop:dynamic_galerkin} of \Cref{app:block_cumulant_analysis}.  Then the truncated covariance operators converge in trace norm:
  \begin{equation}
    \|\Sigma_{n,\lambda}-\Sigma_\lambda\|_{\mathcal{S}_1}\to0\quad\text{as }n\to\infty,
  \end{equation}
  and the truncated Schur invariant is exact at every level:
  \begin{equation}
    m_2^{\mathrm{Schur},(n)} \coloneqq \lim_{\lambda\to\infty} \lambda^2 S_{k,\lambda}^{(n)} = \frac{\varepsilon_0}{2} \quad \text{for all } n \ge k.
  \end{equation}
\end{proposition}
\begin{proof}
  See \Cref{app:block_cumulant_analysis} for the dynamic and spectral support.
\end{proof}

% ============================================================
%  §4  GAUSSIAN CONDITIONAL EXTENSION AND FINITE-PART FUNCTIONAL
% ============================================================
\section{Gaussian Conditional Extension and Finite-Part Functional}
\label{sec:gaussian_zero}

This section gives the Gaussian extension of the base theorem, using Radon Gaussian laws and disintegration to formulate conditioning under exact evidence and define the renormalized finite-part functional.

\subsection{Radon Gaussian Setup and Disintegration}
\label{sec:gaussian:setup}

\begin{assumption}[Hilbert Gaussian Conditioning and Affine Mean Channel]
  \label{assum:gaussian_clamp}
  In addition to Assumptions~\ref{assum:linear_second_moment} and \ref{assum:cm_stability}, let $\mu_\lambda = \mathcal{N}(m_\lambda, \Sigma_\lambda)$ be a Gaussian Radon law on $\mathcal{H}$ with injective trace-class covariance. Let $\mu_{\mathrm{obs}}$ be the canonical Gaussian disintegration member on the evidence hyperplane $X_k = x_k^*$. Suppose that a fixed deterministic forcing $b \in \mathcal{H}$ gives the stationary mean equation:
  \begin{equation}
    \mathcal{A}_{\cond} m_\lambda + b - \lambda \Pi_k (m_\lambda - x_k^* \phi_k) = 0, \quad \Pi_k \mathcal{A}_{\cond} = 0.
  \end{equation}
  We define the target deterministic offset constant:
  \begin{equation}
    c_\lambda \coloneqq \lambda(m_{\lambda,k} - x_k^*) = b_k,
  \end{equation}
  where $m_{\lambda,k} \coloneqq \langle m_\lambda, \phi_k \rangle$ and $b_k \coloneqq \langle b, \phi_k \rangle$.
\end{assumption}

Under the Gaussian Radon law, the coordinate projection $\pi_k: \mathcal{H} \to \mathbb{R}$ is continuous, so the target variable $X_k$ is a well-defined Borel random variable. The Polish structure of the Hilbert space guarantees the existence of a regular conditional probability distribution. 
By selecting the weakly continuous version of the conditional law at the single point $t = x_k^*$ (as detailed in \Cref{app:gaussian_lift}), we obtain the unique pointwise disintegration member $\mu_{\mathrm{obs}}$. Under this measure, the target coordinate satisfies the exact target zeros:
\begin{equation}
  \label{eq:exact_target_zeros}
  m_{\text{obs},k} - x_k^* = 0, \quad \Var_{\text{obs}}(X_k) = 0 \quad (\text{exact}).
\end{equation}

\subsection{Global-Inverse-Free Measurable Predictor}
\label{sec:gaussian:predictor}

Let $C_\lambda \coloneqq \Sigma_{\lambda,II}$ denote the Gaussian bulk covariance operator on $\mathcal{H}_I$. Because $C_\lambda$ is trace-class and compact, it lacks a bounded global inverse. We define the bulk Cameron--Martin space $H_{C_\lambda} \coloneqq \operatorname{Range}(C_\lambda^{1/2})$ and the corresponding covariance norm.

The leakage cross-covariance vector satisfies $\Sigma_{\lambda,Ik} \in H_{C_\lambda}$ with norm $\|\Sigma_{\lambda,Ik}\|_{H_{C_\lambda}} = O(\lambda^{-3})$. This range inclusion allows us to define the conditional predictor $\mu_{k|I}(X_I) \coloneqq \E[X_k \mid X_I]$ as a measurable Wiener functional:
\begin{equation}
  \mu_{k|I}(X_I) = m_{\lambda,k} + W_{\Sigma_{\lambda,Ik}}(X_I - m_{\lambda,I}),
\end{equation}
which is the unique LMMSE predictor in $L^2$. The estimates in \Cref{app:gaussian_lift} imply that, for all sufficiently large $\lambda$, the bulk marginals $\mu_{\text{obs},I}$ and $\mu_{\lambda,I}$ are equivalent and have finite bulk relative entropy:
\begin{equation}
  \KL(\mu_{\text{obs},I} \parallel \mu_{\lambda,I}) = O(\lambda^{-4}) \to 0 \quad \text{as } \lambda \to \infty.
\end{equation}

\subsection{Gaussian Renormalized Finite-Part Functional}
\label{sec:gaussian:functional}

Although the joint KL divergence diverges under the exact conditioning, we can define a renormalized finite-part action.

\begin{definition}[Gaussian Renormalized Finite-Part Functional]
  \label{def:unified_action}
  Whenever $S_{k,\lambda}>0$ and $\mu_{\text{obs},I}$ is equivalent to $\mu_{\lambda,I}$, we define the Gaussian renormalized finite-part functional of the conditioned law:
  \begin{equation}
    \label{eq:unified_action}
    \begin{aligned}
      \mathfrak{J}_\lambda(\mu_{\mathrm{obs}}; \mu_\lambda) \coloneqq &\KL(\mu_{\mathrm{obs},I} \parallel \mu_{\lambda,I}) \\
      &+ \frac{1}{2} S_{k,\lambda}^{-1} \E_{\mu_{\mathrm{obs}}}\left[ \left( X_k - \mu_{k|I}(X_I) \right)^2 \right].
    \end{aligned}
  \end{equation}
  This represents the finite part of the joint KL divergence under the conditional normalization scheme, where the Dirac and Gaussian normalization singularities are subtracted.
\end{definition}

\paragraph{Weak scheme-independence.}
For an evidence level $t\in\mathbb{R}$, write $\mu_\lambda^t$ for the weakly continuous conditional law from \Cref{app:gaussian_lift} and set
\[
  \mathfrak{J}_\lambda(t)\coloneqq
  \mathfrak{J}_\lambda(\mu_\lambda^t;\mu_\lambda).
\]
Then, for any two evidence levels $t_1,t_2$ for which the finite part is defined,
\begin{equation}
  \label{eq:weak_scheme_independence}
  \mathfrak{J}_\lambda(t_2)-\mathfrak{J}_\lambda(t_1)
  =
  \frac{(t_2-m_{\lambda,k})^2-(t_1-m_{\lambda,k})^2}
       {2\Sigma_{\lambda,kk}}.
\end{equation}
In the calibrated scaling $\Sigma_{\lambda,kk}=\varepsilon_0/(2\lambda^2)$, if
$b_i=\lambda(m_{\lambda,k}-t_i)$, this becomes
\[
  \mathfrak{J}_\lambda(t_2)-\mathfrak{J}_\lambda(t_1)
  =
  \frac{b_2^2-b_1^2}{\varepsilon_0}.
\]
\begin{proof}
Let $\Delta_t\coloneqq t-m_{\lambda,k}$ and
\[
  \rho_\lambda
  \coloneqq
  \frac{\|\Sigma_{\lambda,Ik}\|_{H_{C_\lambda}}^2}{\Sigma_{\lambda,kk}}
  =
  1-\frac{S_{k,\lambda}}{\Sigma_{\lambda,kk}}.
\]
The Gaussian entropy computation in \Cref{app:gaussian_lift} gives
\[
  \KL(\mu_{\lambda,I}^t\parallel\mu_{\lambda,I})
  =
  \frac{1}{2}\frac{\Delta_t^2}{\Sigma_{\lambda,kk}}\rho_\lambda
  -\frac{1}{2}\log(1-\rho_\lambda)-\frac{1}{2}\rho_\lambda .
\]
The conditional-residual term in \eqref{eq:unified_action} is
\[
  \frac{1}{2S_{k,\lambda}}
  \E_{\mu_\lambda^t}\!\left[
    (X_k-\mu_{k|I}(X_I))^2
  \right]
  =
  \frac{1}{2}\rho_\lambda
  +
  \frac{1}{2}\frac{\Delta_t^2}{\Sigma_{\lambda,kk}}(1-\rho_\lambda).
\]
Adding these two identities cancels the $\rho_\lambda/2$ terms and yields
\[
  \mathfrak{J}_\lambda(t)
  =
  \frac{(t-m_{\lambda,k})^2}{2\Sigma_{\lambda,kk}}
  -\frac{1}{2}\log(1-\rho_\lambda).
\]
The logarithmic term depends on the reference covariance and the conditioning
scale, but not on the evidence value $t$. Hence it cancels in the difference
between $t_2$ and $t_1$, proving \eqref{eq:weak_scheme_independence}. The same
argument shows weak independence from the mollifying kernel: replacing the
Gaussian Dirac mollifier in \Cref{prop:mollification_limit} by any centered
approximate identity with finite second moment changes the extracted finite
part only by an additive term depending on the kernel, $\lambda$, and
$S_{k,\lambda}$, not on $t$. Therefore all evidence differences
$\Delta\mathfrak{J}_\lambda$ are kernel-independent.
\end{proof}

\subsection{Gaussian Main Theorem}
\label{sec:gaussian:theorem}

We now establish the main regularization and cancellation theorems for the finite-part functional.

\begin{theorem}[Gaussian Finite-Part Regularization]
  \label{thm:finite_part_regularization}
  Under Assumption~\ref{assum:gaussian_clamp}, there exists $\lambda_0>0$ such that, for $\lambda\ge\lambda_0$, the renormalized finite-part functional under the conditioned law $\mu_{\mathrm{obs}}$ is well-defined and satisfies:
  \begin{equation}
    \label{eq:unified_regularization}
    \mathfrak{J}_\lambda(\mu_{\mathrm{obs}}; \mu_\lambda) = \frac{b_k^2}{\varepsilon_0} + O(\lambda^{-4}).
  \end{equation}
  In particular, the functional remains bounded at $O(1)$ and converges to:
  \begin{equation}
    \lim_{\lambda\to\infty} \mathfrak{J}_\lambda(\mu_{\mathrm{obs}}; \mu_\lambda) = \frac{b_k^2}{\varepsilon_0}.
  \end{equation}
\end{theorem}
\begin{proof}
  See \Cref{app:gaussian_lift} for the full derivation.
\end{proof}

\begin{corollary}[Exact Decoupled-Centered Cancellation]
  \label{cor:exact_pole_zero}
  Under Assumption~\ref{assum:gaussian_clamp}, if the target coordinate is dynamically decoupled from the bulk ($a_{Ik} = 0$) and the evidence is centered at the stationary mean ($x_k^* = m_{\lambda,k}$), then:
  \begin{equation}
    \mathfrak{J}_\lambda(\mu_{\mathrm{obs}}; \mu_\lambda) = 0 \quad \text{for every } \lambda > 0.
  \end{equation}
\end{corollary}
\begin{proof}
  The decoupling implies $\Sigma_{\lambda,Ik} = 0$, so $\KL(\mu_{\text{obs},I} \parallel \mu_{\lambda,I}) = 0$. The centering makes the conditional residual zero under the conditioning, so both terms in \eqref{eq:unified_action} vanish.
\end{proof}

% ============================================================
%  §5  ELLIPTIC GAUSSIAN FIELD EXAMPLE
% ============================================================
\section{Elliptic Gaussian Field Example}
\label{sec:spde}

We now record a standard elliptic Gaussian field example to show that the abstract Hilbert-space assumptions are nonempty and natural for covariance operators arising from stochastic evolution equations and elliptic Gaussian priors \cite{daprato2014stochastic,lord2014introduction,stuart2010inverse}.

\subsection{Elliptic Gaussian field on $[0,1]$}

Let
\[
  L=-\frac{d^2}{dx^2}+\kappa^2
\]
on $L^2(0,1)$ with Dirichlet boundary conditions and $\kappa>0$. Then $\mathcal{A}_0=-L$ with domain $H^2(0,1)\cap H_0^1(0,1)$ is dissipative, $L$ is positive self-adjoint and sectorial, and $e^{-tL}$ is an exponentially stable analytic contraction semigroup. The inverse covariance
\[
  Q_0=L^{-1}
\]
is compact, positive, self-adjoint, and trace class. With eigenfunctions $e_j(x)=\sqrt{2}\sin(\pi jx)$, its eigenvalues are
\[
  q_j=(\pi^2j^2+\kappa^2)^{-1}.
\]
Thus $Q_0$ is a typical infinite-dimensional covariance operator: it is compact and trace class, has no bounded inverse on the ambient space, and carries its inverse only as a Cameron--Martin quadratic form \cite{bogachev1998gaussian,simon2005trace}.

\subsection{Cameron--Martin range and covariance norm}

The Cameron--Martin space associated with $Q_0$ is
\[
  H_{Q_0}=\operatorname{Range}(Q_0^{1/2})=H_0^1(0,1),
\]
with equivalent covariance norm
\[
  \|h\|_{Q_0}^2
  =\langle h,Lh\rangle_{L^2}
  =\int_0^1 \bigl(|h'(x)|^2+\kappa^2|h(x)|^2\bigr)\,dx.
\]
This identifies the variational Schur projection in a Sobolev quadratic form, as in the standard covariance-norm representation of Gaussian Hilbert priors \cite{bogachev1998gaussian,stuart2010inverse}. The covariance inverse does not act as a bounded operator on $L^2(0,1)$; it acts only through this quadratic form on its Cameron--Martin domain.

\subsection{Spectral sensor conditioning}

Fix a target eigenmode $e_k$ and define the observed coordinate
\[
  X_k=\langle X,e_k\rangle_{L^2}.
\]
The unperturbed target variance is
\[
  \Sigma_{0,kk}=q_k=(\pi^2k^2+\kappa^2)^{-1}.
\]
A regularized restoring drift on this coordinate produces the exact normalization
\[
  \Sigma_{\lambda,kk}=\frac{\varepsilon_0}{2\lambda^2}-R_\lambda,
  \qquad R_\lambda=O(\lambda^{-6}),
\]
under the block assumptions of \Cref{sec:infdim}. Off-diagonal bounded perturbations of the drift operator generate the cross-covariance vector entering the variational Schur residual. The resulting subtraction is evaluated by the Cameron--Martin quadratic form above, not by a global inverse on $L^2(0,1)$.

This example shows that the abstract theorem applies to an infinite-dimensional Gaussian field: the residual is controlled by a Sobolev norm on $H_0^1(0,1)$ while the ambient covariance remains compact and non-invertible.

\section{Discussion}
\label{sec:discussion}

We have shown that singular Gaussian coordinate conditioning in a separable Hilbert space admits an inverse-free Schur normal form.  The main obstruction is the absence of a bounded inverse for compact covariance blocks.  The variational Cameron--Martin projection is the target-coordinate shorted-operator subtraction evaluated in the covariance-energy scale, and it yields the residual estimate needed for the hard-conditioning limit.  The scalar nature of the target coordinate is essential to the present proof.  Three extensions mark the boundary between the theorem proved here and a genuinely finite-rank conditioning theory. We formalize these extensions as exact theorems below.

\subsection{Matrix algebraic extension}
\label{sec:discussion:matrix-extension}

For an $m$-dimensional observed subspace $E=\operatorname{span}\{\phi_1,\ldots,\phi_m\}$, the scalar target block $\Sigma_{\lambda,kk}$ is replaced by the covariance matrix $\Sigma_{\lambda,EE}\in\mathbb{R}^{m\times m}$. The scalar calibration becomes a matrix Lyapunov balancing problem.

\begin{theorem}[Matrix Lyapunov Calibration]
  \label{thm:matrix_lyapunov}
  Let $E$ be an $m$-dimensional target subspace. Suppose the regularized drift and diffusion on $E$ are given by $A_{\lambda,E} = -\lambda \Lambda_E$ and $D_{\lambda,E} = \frac{1}{\lambda} \Delta_E$, where $\Lambda_E, \Delta_E \in \mathbb{R}^{m \times m}$ are strictly positive definite matrices. Let $\Sigma_{\lambda,EE}$ be the unique positive definite solution to the matrix Lyapunov equation:
  \begin{equation}
    \label{eq:discussion_matrix_lyapunov}
    A_{\lambda,E}\Sigma_{\lambda,EE}
    +\Sigma_{\lambda,EE}A_{\lambda,E}^{*}
    +D_{\lambda,E}=0.
  \end{equation}
  Then $\Sigma_{\lambda,EE}$ scales exactly as $\lambda^{-2}$. Specifically, the normalized matrix $\widetilde{\Sigma}_{EE} \coloneqq \lambda^2 \Sigma_{\lambda,EE}$ is independent of $\lambda$ and is the unique positive definite solution to the scale-free Lyapunov equation:
  \begin{equation}
    \label{eq:scale_free_lyapunov}
    \Lambda_E \widetilde{\Sigma}_{EE} + \widetilde{\Sigma}_{EE} \Lambda_E^* = \Delta_E.
  \end{equation}
\end{theorem}
\begin{proof}
  Substituting $A_{\lambda,E} = -\lambda \Lambda_E$ and $D_{\lambda,E} = \frac{1}{\lambda} \Delta_E$ into \eqref{eq:discussion_matrix_lyapunov} yields:
  \[
    -\lambda \Lambda_E \Sigma_{\lambda,EE} - \lambda \Sigma_{\lambda,EE} \Lambda_E^* + \frac{1}{\lambda} \Delta_E = 0.
  \]
  Multiplying the entire equation by $\lambda$ gives:
  \[
    \Lambda_E (\lambda^2 \Sigma_{\lambda,EE}) + (\lambda^2 \Sigma_{\lambda,EE}) \Lambda_E^* = \Delta_E.
  \]
  Defining $\widetilde{\Sigma}_{EE} \coloneqq \lambda^2 \Sigma_{\lambda,EE}$, we recover \eqref{eq:scale_free_lyapunov}. Since $\Lambda_E$ is positive definite, its spectrum lies in the open right half-plane, making $-\Lambda_E$ Hurwitz. Because $\Delta_E$ is positive definite, standard continuous Lyapunov theory guarantees that \eqref{eq:scale_free_lyapunov} admits a unique positive definite solution $\widetilde{\Sigma}_{EE}$. Therefore, the exact scaling $\Sigma_{\lambda,EE} = \lambda^{-2} \widetilde{\Sigma}_{EE}$ holds for all $\lambda > 0$, properly generalizing the scalar calibration condition $\Sigma_{\lambda,kk} = \varepsilon_0 / (2\lambda^2)$.
\end{proof}

\subsection{Operator-valued residuals}
\label{sec:discussion:operator-valued-residuals}

For an $m$-dimensional target space $E$, the Schur residual must be a positive-semidefinite operator on $E$. Because the free-block covariance $\Sigma_{\lambda,II}$ is compact and lacks a bounded global inverse, the formal finite-dimensional matrix expression $R_{\lambda,E} = \Sigma_{\lambda,EI}\Sigma_{\lambda,II}^{-1}\Sigma_{\lambda,IE}$ must be rigorously replaced by the shorted-operator construction evaluated in the Cameron--Martin energy norm.

\begin{theorem}[Operator-Valued Schur Residual]
  \label{thm:operator_residual}
  Let $\mathcal{H} = E \oplus \mathcal{H}_I$, with $E$ finite-dimensional. Let $\Sigma_\lambda$ be the non-negative trace-class covariance operator on $\mathcal{H}$ with blocks $\Sigma_{\lambda,EE}$, $\Sigma_{\lambda,EI}$, $\Sigma_{\lambda,IE}$, and $\Sigma_{\lambda,II}$. Define the shorted-operator residual $R_{\lambda,E}$ on $E$ by its quadratic forms:
  \begin{equation}
    \label{eq:discussion_operator_residual_form}
    \langle v,R_{\lambda,E}v\rangle_E
    =
    \sup_{\langle u,\Sigma_{\lambda,II}u\rangle>0}
    \frac{|\langle u,\Sigma_{\lambda,IE}v\rangle|^2}
         {\langle u,\Sigma_{\lambda,II}u\rangle},
    \qquad v\in E.
  \end{equation}
  Assume the structural range inclusion $\operatorname{Range}(\Sigma_{\lambda,IE}) \subset \operatorname{Range}(\Sigma_{\lambda,II}^{1/2}) \eqqcolon H_{C_\lambda}$. Then $R_{\lambda,E}$ is a well-defined bounded positive-semidefinite operator on $E$, uniquely characterized by:
  \begin{equation}
    R_{\lambda,E} = (\Sigma_{\lambda,II}^{-1/2}\Sigma_{\lambda,IE})^* (\Sigma_{\lambda,II}^{-1/2}\Sigma_{\lambda,IE}).
  \end{equation}
  Furthermore, if the cross-covariance satisfies the energy bound $\|\Sigma_{\lambda,IE} v\|_{H_{C_\lambda}} \le M \lambda^{-3} \|v\|_E$ for all $v \in E$ and some constant $M > 0$, then the operator norm of the residual satisfies $\|R_{\lambda,E}\|_{\mathcal{L}(E)} = O(\lambda^{-6})$.
\end{theorem}
\begin{proof}
  By the assumed range inclusion, for any $v \in E$, the vector $y_v \coloneqq \Sigma_{\lambda,IE} v$ lies in $H_{C_\lambda} = \operatorname{Range}(\Sigma_{\lambda,II}^{1/2})$. Thus, there exists a unique element $w_v \in \overline{\operatorname{Range}(\Sigma_{\lambda,II}^{1/2})}$ such that $\Sigma_{\lambda,II}^{1/2} w_v = y_v$. We denote $w_v \coloneqq \Sigma_{\lambda,II}^{-1/2} \Sigma_{\lambda,IE} v$.
  
  For any $u \in \mathcal{H}_I$ with $\langle u, \Sigma_{\lambda,II} u \rangle > 0$, the ratio in \eqref{eq:discussion_operator_residual_form} can be rewritten as:
  \[
    \frac{|\langle u, \Sigma_{\lambda,II}^{1/2} w_v \rangle|^2}{\|\Sigma_{\lambda,II}^{1/2} u\|^2} = \frac{|\langle \Sigma_{\lambda,II}^{1/2} u, w_v \rangle|^2}{\|\Sigma_{\lambda,II}^{1/2} u\|^2}.
  \]
  By the Cauchy--Schwarz inequality, this ratio is bounded above by $\|w_v\|^2$. The supremum is achieved by taking a sequence of vectors $\Sigma_{\lambda,II}^{1/2} u$ converging to $w_v$, yielding:
  \[
    \langle v, R_{\lambda,E} v \rangle_E = \|w_v\|^2 = \|\Sigma_{\lambda,II}^{-1/2} \Sigma_{\lambda,IE} v\|^2 = \langle v, (\Sigma_{\lambda,II}^{-1/2}\Sigma_{\lambda,IE})^* (\Sigma_{\lambda,II}^{-1/2}\Sigma_{\lambda,IE}) v \rangle_E.
  \]
  Since this holds for all $v \in E$, it uniquely defines $R_{\lambda,E}$ as a positive-semidefinite matrix on $E$. Finally, the assumption $\|\Sigma_{\lambda,IE} v\|_{H_{C_\lambda}} \le M \lambda^{-3} \|v\|_E$ means exactly that $\|w_v\| \le M \lambda^{-3} \|v\|_E$. Therefore, $\langle v, R_{\lambda,E} v \rangle_E \le M^2 \lambda^{-6} \|v\|_E^2$, which implies the operator norm bound $\|R_{\lambda,E}\|_{\mathcal{L}(E)} \le M^2 \lambda^{-6} = O(\lambda^{-6})$.
\end{proof}

\subsection{General observation operators and misalignment}
\label{sec:discussion:misalignment}

The scalar proof in this paper uses a block decomposition aligned with a single observed coordinate. However, typical finite-rank observation operators $L:\mathcal H\to\mathbb R^m$ select a subspace $E$ that does not commute with the free prior covariance, meaning the target-free split is not an eigenbasis split. This geometric misalignment generates cross-coupling that requires a precise resolvent analysis.

\begin{theorem}[Misaligned Covariance Coupling]
  \label{thm:misaligned_coupling}
  Let $P_E$ be an orthogonal projection onto an $m$-dimensional observed subspace $E$, and let $\mathcal{H}_I = E^\perp$. Let the regularized dynamics be decomposed such that the target block matches \Cref{thm:matrix_lyapunov} with $A_{\lambda,E} = -\lambda \Lambda_E$, the cross-noise is zero ($D_{\lambda,IE} = 0$), and the off-diagonal drift coupling is denoted by $A_{\lambda,IE}$. 
  If $A_{\lambda,IE}$ maps $E$ into the Cameron--Martin space $H_{Q_0}$ of the free background covariance $Q_0$ on $\mathcal{H}_I$, and the restricted drift $A_{II}$ generates an exponentially stable semigroup on $H_{Q_0}$, then the cross-covariance satisfies the exact identity:
  \begin{equation}
    \label{eq:misaligned_resolvent_identity}
    \Sigma_{\lambda,IE} = \int_0^\infty e^{t A_{II}} A_{\lambda,IE} \Sigma_{\lambda,EE} e^{t A_{\lambda,E}^*} dt.
  \end{equation}
  Furthermore, the misaligned cross-covariance energy scales as $\|\Sigma_{\lambda,IE} v\|_{H_{Q_0}} = O(\lambda^{-3})$ for any $v \in E$.
\end{theorem}
\begin{proof}
  The off-diagonal block of the global Lyapunov equation $\mathcal{A}^{(\lambda)} \Sigma_\lambda + \Sigma_\lambda (\mathcal{A}^{(\lambda)})^* + \mathcal{D}^{(\lambda)} = 0$ yields the Sylvester equation:
  \[
    A_{II} \Sigma_{\lambda,IE} + \Sigma_{\lambda,IE} A_{\lambda,E}^* + A_{\lambda,IE} \Sigma_{\lambda,EE} = 0.
  \]
  Since $A_{II}$ generates an exponentially stable semigroup $e^{t A_{II}}$ and $A_{\lambda,E} = -\lambda \Lambda_E$ is Hurwitz, the unique solution is given by the integral representation \eqref{eq:misaligned_resolvent_identity}.
  
  To bound the Cameron--Martin norm, fix $v \in E$. We have $e^{t A_{\lambda,E}^*} v = e^{-\lambda t \Lambda_E^*} v$. By \Cref{thm:matrix_lyapunov}, $\Sigma_{\lambda,EE} = O(\lambda^{-2})$. Since $A_{\lambda,IE}$ maps into $H_{Q_0}$ and $E$ is finite-dimensional, there exists a constant $C$ such that $\|A_{\lambda,IE} w\|_{H_{Q_0}} \le C \|w\|_E$ for all $w \in E$. Using the uniform exponential stability of $e^{t A_{II}}$ on $H_{Q_0}$, namely $\|e^{t A_{II}}\|_{\mathcal{L}(H_{Q_0})} \le M e^{-\omega t}$ for some $M \ge 1, \omega > 0$, we estimate:
  \begin{align*}
    \|\Sigma_{\lambda,IE} v\|_{H_{Q_0}} 
    &\le \int_0^\infty \|e^{t A_{II}}\|_{\mathcal{L}(H_{Q_0})} \|A_{\lambda,IE} \Sigma_{\lambda,EE} e^{-\lambda t \Lambda_E^*} v\|_{H_{Q_0}} dt \\
    &\le \int_0^\infty M e^{-\omega t} C \|\Sigma_{\lambda,EE}\|_{\mathcal{L}(E)} \|e^{-\lambda t \Lambda_E^*}\|_{\mathcal{L}(E)} \|v\|_E dt.
  \end{align*}
  Let $\gamma > 0$ be the smallest eigenvalue of the positive definite matrix $\Lambda_E^*$, so $\|e^{-\lambda t \Lambda_E^*}\|_{\mathcal{L}(E)} \le K e^{-\lambda \gamma t}$. Then:
  \begin{align*}
    \|\Sigma_{\lambda,IE} v\|_{H_{Q_0}} 
    &\le M C K \|\Sigma_{\lambda,EE}\|_{\mathcal{L}(E)} \|v\|_E \int_0^\infty e^{-(\omega + \lambda \gamma) t} dt \\
    &= M C K \|\Sigma_{\lambda,EE}\|_{\mathcal{L}(E)} \|v\|_E \frac{1}{\omega + \lambda \gamma}.
  \end{align*}
  Since $\|\Sigma_{\lambda,EE}\|_{\mathcal{L}(E)} = O(\lambda^{-2})$ and $(\omega + \lambda \gamma)^{-1} = O(\lambda^{-1})$, we conclude that $\|\Sigma_{\lambda,IE} v\|_{H_{Q_0}} = O(\lambda^{-3})$. Combined with \Cref{thm:operator_residual}, this verifies that the shorted-operator residual maintains the $O(\lambda^{-6})$ scaling even in the presence of observation misalignment.
\end{proof}

The result is therefore deliberately limited to covariance and Gaussian-conditioning questions in the rank-one aligned setting.  It does not assert a general non-Gaussian conditioning theory.  For non-Gaussian stationary processes, the projection residual $S_{k,\lambda}=\varepsilon_0/(2\lambda^2)-R_\lambda$ still has a second-moment interpretation, but Radon disintegration, finite-part functionals, and measure-class behavior require separate hypotheses \cite{bogachev2007measure,kallenberg2021foundations}.  Dynamical relaxation after conditioning, or transport distances between post-conditioning laws, belongs to a different problem.  The present contribution is the operator-theoretic construction of the singular Gaussian conditioning residual itself.

% ============================================================
%  \S6  CONCLUSION
% ============================================================
\section{Conclusion}
\label{sec:conclusion}

We established a Hilbert-space framework for singular coordinate conditioning of Gaussian measures. The central object is an inverse-free variational Schur projection on the Cameron--Martin range, equivalently the one-dimensional target-coordinate form of the shorted-operator subtraction. It gives canonical normalization, an $O(\lambda^{-6})$ residual estimate, Galerkin stability, and a Radon Gaussian conditioning statement. The elliptic field example shows that the assumptions are natural for trace-class covariance operators arising from stochastic evolution equations and elliptic Gaussian priors \cite{daprato2014stochastic,lord2014introduction,stuart2010inverse}.

\section*{Declaration of Generative AI and AI-Assisted Technologies in the Manuscript Preparation Process}

During the preparation of this work, the author(s) used Google DeepMind's Antigravity assistant in order to merge the LaTeX manuscripts, clean up cross-referencing tags, format mathematical notation, and edit the draft for style consistency. After using this tool/service, the author(s) reviewed and edited the content as needed and take(s) full responsibility for the content of the published article.

% ============================================================
%  REFERENCES
% ============================================================
\bibliographystyle{amsplain}

\documentclass[11pt]{amsart}

\usepackage{amsmath}
\usepackage{amssymb}
\usepackage{amsthm}
\usepackage{mathtools}
\usepackage{enumitem}
%\usepackage{thmtools}
\usepackage{hyperref}
\hypersetup{hypertexnames=false}
\usepackage[nameinlink,noabbrev]{cleveref}
\usepackage{booktabs}
\usepackage{graphicx}
\usepackage{xcolor}


% ---- Theorem environments ----
\theoremstyle{plain}
\newtheorem{theorem}{Theorem}
\newtheorem{proposition}[theorem]{Proposition}
\newtheorem{lemma}[theorem]{Lemma}
\newtheorem{corollary}[theorem]{Corollary}
\newtheorem{claim}[theorem]{Claim}

\theoremstyle{definition}
\newtheorem{definition}[theorem]{Definition}
\newtheorem{assumption}[theorem]{Assumption}
\newtheorem{example}[theorem]{Example}

\theoremstyle{remark}
\newtheorem{remark}[theorem]{Remark}

% ---- Macros ----
\newcommand{\cond}{\operatorname{cond}}
\newcommand{\KL}{D_{\operatorname{KL}}}
\newcommand{\E}{\mathbb{E}}
\newcommand{\Var}{\operatorname{Var}}
\newcommand{\Cov}{\operatorname{Cov}}
\newcommand{\tr}{\operatorname{Tr}}
\newcommand{\diag}{\operatorname{diag}}

\newcommand{\Range}{\operatorname{Range}}
\newcommand{\Span}{\operatorname{span}}
\newcommand{\id}{\mathrm{id}}
% \newcommand{\TODO}[1]{\textcolor{red}{[TODO: #1]}} % removed for submission

% ---- Cleveref names ----
\crefname{theorem}{theorem}{theorems}
\Crefname{theorem}{Theorem}{Theorems}
\crefname{proposition}{proposition}{propositions}
\Crefname{proposition}{Proposition}{Propositions}
\crefname{lemma}{lemma}{lemmas}
\Crefname{lemma}{Lemma}{Lemmas}
\crefname{corollary}{corollary}{corollaries}
\Crefname{corollary}{Corollary}{Corollaries}
\crefname{claim}{claim}{claims}
\Crefname{claim}{Claim}{Claims}
\crefname{definition}{definition}{definitions}
\Crefname{definition}{Definition}{Definitions}
\crefname{assumption}{assumption}{assumptions}
\Crefname{assumption}{Assumption}{Assumptions}
\crefname{example}{example}{examples}
\Crefname{example}{Example}{Examples}
\crefname{remark}{remark}{remarks}
\Crefname{remark}{Remark}{Remarks}
\crefname{equation}{Eq.}{Eqs.}
\Crefname{equation}{Equation}{Equations}
\crefname{figure}{fig.}{figs.}
\Crefname{figure}{Fig.}{Figs.}
\crefname{table}{table}{tables}
\Crefname{table}{Table}{Tables}
\crefname{section}{Sec.}{Secs.}
\Crefname{section}{Section}{Sections}

\title[Singular Gaussian Conditioning in Hilbert Space]{Singular Gaussian Conditioning in Hilbert Space: Inverse-Free Schur Projection and Variance Collapse}

\author{Lucia Shang}
\thanks{Manuscript received May 2026. Repository information omitted for submission.}

\begin{document}

% ============================================================
%  ABSTRACT
% ============================================================
\begin{abstract}
We study singular coordinate conditioning for Gaussian measures on separable Hilbert spaces. Since trace-class covariance blocks are compact, the usual Schur-complement formula cannot be implemented through a bounded inverse. We introduce an inverse-free variational Schur projection on the Cameron--Martin range and prove canonical normalization together with an $O(\lambda^{-6})$ residual estimate. The construction is stable under Galerkin truncation for exponentially stable analytic semigroup dynamics and applies to sectorial elliptic Gaussian fields.
\end{abstract}

\keywords{Singular conditioning; Gaussian measures; Cameron--Martin space; Schur complement; Hilbert-space covariance operators; trace-class covariance; Radon Gaussian conditioning; Galerkin approximation; stochastic evolution equations}

\maketitle

% ============================================================
%  \S1  INTRODUCTION
% ============================================================
\section{Introduction}
\label{sec:intro}

Singular coordinate conditioning arises when an exact constraint $X_k=x_k^*$ is reached as a zero-noise or infinite-precision limit of Gaussian conditioning problems.  Finite-dimensional Gaussian conditioning is governed by the ordinary Schur complement of a covariance matrix \cite{horn2012matrix,rasmussen2006gaussian}.  Infinite-dimensional Gaussian conditioning is subtler: Radon Gaussian measures on separable Hilbert or Banach spaces have compact covariance operators, and their Cameron--Martin spaces are covariance-norm domains rather than ambient statistical parameter spaces \cite{bogachev1998gaussian,daprato2014stochastic,steinwart2024conditioning}.  In Hilbert-space Gaussian conditioning, shorted operators give the covariance of conditioning on closed subspaces, while recent inverse-problem work emphasizes that positive-noise posterior formulas must be interpreted through approximation, disintegration, or Hilbert-space posterior equivalence rather than through a global covariance inverse \cite{owhadi2015conditioning,stuart2010inverse,chen2024gaussian,carerelie2024optimal,carerelie2025posterior}.

The obstruction studied here is the covariance-block obstruction behind that phenomenon.  In a finite-dimensional model, the conditional variance of a coordinate can be written as $A-BD^{-1}B^*$ after a block decomposition.  In a separable Hilbert space the corresponding block $D$ is typically compact and trace class, hence $D^{-1}$ is not a bounded operator on the ambient space \cite{bogachev1998gaussian,simon2005trace,daprato2014stochastic}.  Classical Schur-complement and shorted-operator theory, originating in Krein's theory of non-negative extensions and developed systematically by Anderson and Trapp, provides variational tools for positive block operators \cite{krein1947selfadjoint,anderson1975shorted,douglas1966majorization}.  However, this theory does not by itself give the scaling law for a singular Gaussian conditioning residual \cite{tretter2008spectral,contino2018schur,friedrich2017generalized,hassi2023complementation}.  This paper formulates the residual directly as a Cameron--Martin norm projection and never invokes a pseudoinverse or an ambient precision operator.

The hard-conditioning limit also reaches the Feldman--H\'{a}jek measure-class boundary.  With positive observational noise, Gaussian Bayesian inverse problems may remain in the equivalent-measure regime \cite{stuart2010inverse,carerelie2024optimal,carerelie2025posterior}.  Under an exact coordinate constraint, the limiting conditional law is supported on a hyperplane of prior measure zero, so the equivalent-measure regime is left \cite{feldman1958equivalence,hajek1958property,bogachev1998gaussian}.  A useful theory must therefore separate the diverging joint divergence caused by exact evidence from the finite covariance residual that remains visible in the block structure.

We prove that this finite residual is controlled by an inverse-free variational Schur projection.  Under the Lyapunov block assumptions stated below, the target residual has the exact form
\[
  S_{k,\lambda}=\frac{\varepsilon_0}{2\lambda^2}-R_\lambda,
  \qquad
  R_\lambda=O(\lambda^{-6}).
\]
Consequently,
\[
  S_{k,\lambda}^{-1}=\frac{2}{\varepsilon_0}\lambda^2+O(\lambda^{-2}).
\]
The proof uses covariance-norm forms, shorting, Douglas range inclusion, Lyapunov block identities, and trace-ideal convergence. These are standard inputs in operator and SPDE analysis \cite{krein1947selfadjoint,anderson1975shorted,douglas1966majorization} \cite{simon2005trace,daprato2014stochastic}. The Cameron--Martin space is used only as a Hilbert subspace equipped with its covariance norm; no differentiable statistical or geometric structure is assumed \cite{bogachev1998gaussian}.
The operator base is an exponentially stable analytic semigroup with bounded drift perturbations, so dissipative sectorial generators such as $-(-\Delta+\kappa^2)$ fall inside the abstract setup.

The paper has four main contributions:
\begin{enumerate}
  \item \textbf{Canonical normalization.} We identify the block-Lyapunov mechanism that fixes the observed coordinate variance at order $\lambda^{-2}$ with explicit coefficient $\varepsilon_0/2$.
  \item \textbf{Inverse-free Schur projection.} We identify the residual subtraction with the target-coordinate quadratic form of the shorted-operator variational formula, then prove the $O(\lambda^{-6})$ residual estimate in the same operator-theoretic regime where compact covariance blocks preclude a bounded Schur inverse.
  \item \textbf{Analytic-semigroup and Galerkin stability.} We work with an analytic semigroup drift base plus bounded perturbations, and prove that finite-dimensional truncations converge to the Hilbert-space residual in trace ideal norm, matching the discretization-independent perspective used in Hilbert-space Gaussian inverse problems \cite{stuart2010inverse,carerelie2024optimal,carerelie2025posterior}.
  \item \textbf{Radon Gaussian conditioning.} For Radon Gaussian laws, we construct the conditional predictor as a measurable Cameron--Martin projection and show that the associated renormalized conditioning functional remains bounded, using regular conditional probabilities and continuous Gaussian disintegration \cite{bogachev1998gaussian,lagatta2013continuous,steinwart2024conditioning}.
\end{enumerate}

\subsection{Positioning in the literature}
\label{sec:related}

The paper is closest to four bodies of work.  First, it uses the classical measure-class structure of Gaussian measures, especially the Cameron--Martin theorem and the Feldman--H\'{a}jek dichotomy, as presented in standard references on Gaussian measures \cite{feldman1958equivalence,hajek1958property,bogachev1998gaussian}.  Second, it is adjacent to Bayesian inverse problems and Hilbert-space Gaussian conditioning, including the shorted-operator description of Gaussian conditioning on closed subspaces, positive-noise posterior formulas, low-rank posterior approximations, and nonlinear observation maps \cite{owhadi2015conditioning,stuart2010inverse,chen2024gaussian,steinwart2024conditioning,carerelie2024optimal,carerelie2025posterior}.  Third, the variational residual is informed by Schur complements, shorted operators, Douglas range inclusion, and modern complementability theory for block operators \cite{krein1947selfadjoint,anderson1975shorted,douglas1966majorization,tretter2008spectral,contino2018schur,friedrich2017generalized,hassi2023complementation}.  Fourth, the motivating examples come from stochastic evolution equations and elliptic Gaussian fields, where dissipative sectorial generators produce compact trace-class stationary covariances that are naturally approximated by Galerkin truncation \cite{daprato2014stochastic,lord2014introduction,simon2005trace}.

\subsection{Organization of the paper}
\label{sec:organization}

\Cref{sec:setup} defines the operator model and establishes the canonical normalization. \Cref{sec:infdim} establishes the Hilbert-space variational Schur residual. \Cref{sec:gaussian_zero} gives the Radon Gaussian conditional extension and the renormalized conditioning functional. \Cref{sec:spde} gives an elliptic Gaussian-field example. \Cref{sec:discussion} records limitations and extensions, and \Cref{sec:conclusion} concludes. Proofs are given in \Cref{app:block_cumulant_analysis,app:gaussian_lift}; numerical diagnostics and Galerkin verification are given in \Cref{app:experiments}.

\section{Operator Model and Canonical Normalization}
\label{sec:setup}

\subsection{Operator Dynamics and Regularized Conditioning}
\label{sec:setup:dynamics}

We model the continuous-time dynamics as stochastic evolution on a separable Hilbert space $\mathcal{H}$ \cite{daprato2014stochastic}:
\begin{equation}
  \label{eq:stochastic_evolution}
  dX_t = \mathcal{A} X_t \, dt + \mathcal{D}^{1/2} dW_t
\end{equation}
where $\mathcal{A}$ is the drift operator and $\mathcal{D}$ is the trace-class diffusion operator. We decompose the Hilbert space as $\mathcal{H} = \mathcal{H}_k \oplus \mathcal{H}_I$, where $\mathcal{H}_k \coloneqq \mathbb{R}\phi_k$ is the one-dimensional target coordinate subspace, and $\mathcal{H}_I \coloneqq \mathcal{H}_k^\perp$ is the free bulk subspace. Let $\Pi_k \coloneqq \phi_k \otimes \phi_k$ and $\Pi_I \coloneqq I - \Pi_k$ be the corresponding orthogonal projections.

To approximate the exact point conditioning $X_k = x_k^*$, we set the incoming target-coordinate blocks to zero and apply a regularization parameter $\lambda > 0$.

\begin{definition}[Calibrated Approximating Dynamics Family]
  \label{def:soft_do}
  The regularized conditioning operator of strength $\lambda > 0$ at coordinate $X_k$ acts on the operator dynamics \eqref{eq:stochastic_evolution} by modifying both the drift and the diffusion:
  \begin{equation}
    \label{eq:soft_do_sde}
    \mathcal{A}^{(\lambda)} \coloneqq \mathcal{A}_{\cond} - \lambda \Pi_k, \qquad
    \mathcal{D}^{(\lambda)} \coloneqq \mathcal{D}_{\cond} + d_\lambda \Pi_k
  \end{equation}
  where $\mathcal{A}_{\cond}$ is the block drift operator with incoming target-coordinate blocks set to zero ($\Pi_k \mathcal{A}_{\cond} = 0$), $\mathcal{D}_{\cond}$ has target-coordinate cross-noise set to zero, and $d_\lambda > 0$ is the local diffusion variance of the target coordinate at regularization strength $\lambda$. In the SDE, this corresponds to:
  \begin{equation}
    \label{eq:soft_do_sde_general}
    dX_{k,t} = -\lambda (X_{k,t} - x_k^*) dt + \sqrt{d_\lambda}\, dW_{k,t}.
  \end{equation}
\end{definition}

\subsection{Canonical Normalization}
\label{sec:setup:calibration}

The block-decoupled structure under the regularized dynamics yields:
\begin{equation}
  \mathcal{A}^{(\lambda)}=
  \begin{pmatrix}-\lambda&0\\a_{Ik}&\mathcal{A}_{II}\end{pmatrix},\qquad
  \mathcal{D}^{(\lambda)}=
  \begin{pmatrix}d_\lambda&0\\0&D_{II}\end{pmatrix},
\end{equation}
where $a_{Ik} \coloneqq \mathcal{A}_{Ik} \phi_k \in \mathcal{H}_I$ represents the off-diagonal drift coupling.
Let $\Sigma_\lambda$ denote the stationary covariance operator under regularization strength $\lambda$, solving the Lyapunov equation:
\begin{equation}
  \mathcal{A}^{(\lambda)}\Sigma_\lambda + \Sigma_\lambda(\mathcal{A}^{(\lambda)})^* + \mathcal{D}^{(\lambda)} = 0.
\end{equation}
The target block $kk$ decouples completely from the free coordinates, yielding the exact scalar equation:
\begin{equation}
  -2\lambda\Sigma_{\lambda,kk} + d_\lambda = 0 \quad \Longrightarrow \quad \Sigma_{\lambda,kk} = \frac{d_\lambda}{2\lambda}.
\end{equation}

\begin{proposition}[Canonical Scaling]
  \label{prop:exact_calibration}
  To maintain a target noise scale $\varepsilon_0 > 0$ independent of the auxiliary parameter $\lambda$ under the singular limit, we define the canonical scaling condition:
  \begin{equation}
    \lambda^2 \Sigma_{\lambda,kk} \equiv \frac{\varepsilon_0}{2} \quad \text{for all } \lambda > 0,
  \end{equation}
  which uniquely requires:
  \begin{equation}
    d_\lambda = \frac{\varepsilon_0}{\lambda} \quad \text{for every } \lambda > 0.
  \end{equation}
  Substitution of this target diffusion scaling yields the calibrated approximating dynamics:
  \begin{equation}
    dX_{k,t} = -\lambda (X_{k,t} - x_k^*) dt + \sqrt{\varepsilon_0 / \lambda}\, dW_{k,t}.
  \end{equation}
\end{proposition}

All subsequent results are stated for this calibrated family.

\begin{definition}[Schur LMMSE Invariant]
  \label{def:schur_invariant}
  Let $S_{k,\lambda} \coloneqq \|X_k - P_{\mathcal{M}_I} X_k\|_{L^2}^2$ represent the linear minimum mean squared error (LMMSE) residual of the target coordinate $X_k$ on the closed linear span $\mathcal{M}_I$ of the free coordinates $X_I$. The \emph{Schur LMMSE invariant} $m_2^{\mathrm{Schur}}$ is defined as:
  \begin{equation}
    \label{eq:m2_schur_definition}
    m_2^{\mathrm{Schur}} \coloneqq \lim_{\lambda\to\infty} \lambda^2 S_{k,\lambda}.
  \end{equation}
\end{definition}

\subsection{Operator Setup and Assumptions}
\label{sec:setup:assumptions}

We state the operator-theoretic assumptions that make the limit passage independent of the Galerkin dimension.

\begin{definition}[Drift Operator with Analytic Base]
  \label{def:drift_operator}
  A drift operator on a separable Hilbert space $\mathcal{H}$ is an operator of the form $\mathcal{A} = \mathcal{A}_0 + \mathcal{K}$, where:
  \begin{enumerate}
    \item[(i)] $\mathcal{A}_0: \operatorname{dom}(\mathcal{A}_0) \subset \mathcal{H} \to \mathcal{H}$ is closed, densely defined, and dissipative, and it generates a strongly continuous analytic contraction semigroup $(e^{t\mathcal{A}_0})_{t \ge 0}$ on $\mathcal{H}$ satisfying $\|e^{t\mathcal{A}_0}\| \le e^{-\alpha t}$ for some spectral gap $\alpha > 0$; equivalently, the base operator $-\mathcal{A}_0$ is sectorial of angle $<\pi/2$ in this exponentially stable contraction setting.
    \item[(ii)] $\mathcal{K}: \mathcal{H} \to \mathcal{H}$ is a bounded linear operator.
    \item[(iii)] The full operator $\mathcal{A} = \mathcal{A}_0 + \mathcal{K}$, with $\operatorname{dom}(\mathcal{A})=\operatorname{dom}(\mathcal{A}_0)$, generates a strongly continuous analytic semigroup $(e^{t\mathcal{A}})_{t \ge 0}$ by the bounded perturbation theorem, and this semigroup remains exponentially stable: $\|e^{t\mathcal{A}}\| \le M e^{-\beta t}$ for some $M \ge 1$, $\beta > 0$.
  \end{enumerate}
  In finite dimensions ($\mathcal{H} = \mathbb{R}^n$), this reduces to a decomposition of a Hurwitz drift matrix into an exponentially stable analytic base plus a bounded perturbation.
\end{definition}

\begin{definition}[Trace-Class Noise and Covariance]
  \label{def:trace_class}
  Let $\mathcal{S}_1(\mathcal{H})$ denote the Banach space of trace-class (nuclear) operators on $\mathcal{H}$ equipped with the trace norm $\|T\|_{\mathcal{S}_1} \coloneqq \tr(|T|) = \sum_{i=1}^\infty \langle \phi_i, (T^*T)^{1/2} \phi_i \rangle < \infty$.
  \begin{enumerate}
    \item[(i)] The diffusion operator $\mathcal{D}: \mathcal{H} \to \mathcal{H}$ is positive, self-adjoint, and belongs to $\mathcal{S}_1(\mathcal{H})$.
    \item[(ii)] Assuming the pair $(\mathcal{A}, \mathcal{D}^{1/2})$ is controllable, the stationary covariance operator $\Sigma: \mathcal{H} \to \mathcal{H}$ is the unique strictly positive, self-adjoint, trace-class solution to the operator Lyapunov equation:
    \begin{equation}
      \label{eq:op_lyapunov}
      \mathcal{A}\Sigma + \Sigma\mathcal{A}^* + \mathcal{D} = 0
    \end{equation}
    which is given by the convergent Bochner integral
    \[
      \Sigma = \int_0^\infty e^{t\mathcal{A}} \mathcal{D} e^{t\mathcal{A}^*} dt.
    \]
  \end{enumerate}
\end{definition}

\begin{assumption}[Core Analytical Assumptions for Singular Hyperplane Conditioning]
  This work relies on two core analytical assumptions regarding the target-coordinate block decoupling and the outgoing coupling:
  \begin{enumerate}
    \item[(i)] \textbf{Linear Second-Moment Channel:} \label{assum:linear_second_moment} A centered linear stationary process on $\mathcal{H}=\mathcal{H}_k\oplus \mathcal{H}_I$ has finite second moments and covariance solving \Cref{eq:op_lyapunov}. The block-decoupling constraint zeroes all incoming target drift and cross-noise, so that $\mathcal{A}_{kI}=\mathcal{D}_{kI}=\mathcal{D}_{Ik}=0$ and $\mathcal{A}_{kk}=-\lambda$. The target diffusion $\mathcal{D}_{kk}=\varepsilon_0/\lambda$ is the unique scaling determined by the canonical normalization (Proposition~\ref{prop:exact_calibration}). The free noise is independent of the local target noise. (Gaussianity is not assumed).
    \item[(ii)] \textbf{Cameron--Martin Stability of Off-Diagonal Drift Coupling:} \label{assum:cm_stability} Let $Q_0\coloneqq\Sigma_{II}^{(0)}$ denote the free-block covariance produced by the free noise under block decoupling, and let $H_{Q_0}\coloneqq\operatorname{Range}(Q_0^{1/2})$ with the covariance norm $|u|_{Q_0}\coloneqq\|Q_0^{-1/2}u\|_{\mathcal{H}_I}$. The outgoing vector $a_{Ik}\coloneqq\mathcal{A}_{Ik}\phi_k$ belongs to $H_{Q_0}$, and $e^{t\mathcal{A}_{II}}$ restricts to a strongly continuous exponentially stable semigroup on $H_{Q_0}$:
    \begin{equation}
      \|e^{t\mathcal{A}_{II}}\|_{\mathcal{L}(H_{Q_0})} \le M e^{-\omega t},\qquad t\ge0,
    \end{equation}
    for constants $M\ge1$ and $\omega>0$ independent of $\lambda$.
  \end{enumerate}
  \vspace{1mm}
  \noindent\textit{Intuitive Explanation:} Part (i) states that the regularized dynamics act as a strong localized drift penalization on the target coordinate $X_k$ while isolating it from incoming influences, giving an operator implementation of the block-decoupling constraint. Part (ii) requires that the off-diagonal drift coupling $a_{Ik}$ from the target coordinate has finite covariance norm relative to the background fluctuations of the free coordinates. This ensures that target fluctuations propagating downstream do not overload the variance of the infinite-dimensional system.
\end{assumption}

\begin{remark}[Finite-Norm Outgoing Covariance Coupling and Finite-Dimensional Calibration]
  \label{rem:cm_information_flow}
  The membership $a_{Ik}\in H_{Q_0}$ is an admissibility condition,
  not a consequence of the free-block Lyapunov equation alone: an arbitrary
  outgoing vector need not lie in
  $\operatorname{Range}((\Sigma_{II}^{(0)})^{1/2})$.  It says that the
  deterministic channel carrying target fluctuations into the free-block has
  finite covariance norm relative to the free-noise background.
  
  To understand this condition intuitively, consider a finite-dimensional setting where the free-block state space is $\mathcal{H}_I = \mathbb{R}^{d-1}$ and the free stationary covariance $Q_0 \coloneqq \Sigma_{II}^{(0)}$ is strictly positive-definite. Since $Q_0 \succ 0$, its square root $Q_0^{1/2}$ is invertible, and its range is the entire space $\mathbb{R}^{d-1}$. In this case, the Cameron--Martin space is the full space $H_{Q_0} = \operatorname{Range}(Q_0^{1/2}) = \mathbb{R}^{d-1}$, and the Cameron--Martin norm $|u|_{Q_0} = \|Q_0^{-1/2} u\|_2$ is equivalent to the standard Euclidean norm. Consequently, the admissibility condition $a_{Ik} \in H_{Q_0}$ is trivially satisfied for any outgoing vector $a_{Ik} \in \mathbb{R}^{d-1}$.
  
  However, if some direction in the free coordinates is noise-free under block decoupling (meaning $Q_0$ is positive semidefinite but singular), $H_{Q_0}$ becomes a proper subspace of $\mathbb{R}^{d-1}$. In this case, the condition $a_{Ik} \in H_{Q_0}$ requires that the leakage vector $a_{Ik}$ has no component in the null space of $Q_0$. Equivalently, it prevents the target coordinate from leaking fluctuations into a downstream coordinate that has zero background noise, which would otherwise result in an infinite relative covariance norm.
  
  In infinite dimensions, $\Sigma_{II}^{(0)}$ is trace-class and compact, so its eigenvalues decay to zero, making $\operatorname{Range}((\Sigma_{II}^{(0)})^{1/2})$ a dense but strictly proper subspace of $\mathcal{H}_I$. The condition $a_{Ik} \in H_{Q_0}$ requires the off-diagonal drift coupling coefficients to decay sufficiently fast relative to the background noise eigenvalues to keep the relative covariance norm finite.
  
  Under this assumption, invariance and exponential stability on $H_{Q_0}$ imply:
  \begin{equation}
  \begin{aligned}
    (\lambda I-\mathcal{A}_{II})^{-1}a_{Ik}
       &\in H_{Q_0},\\
    |(\lambda I-\mathcal{A}_{II})^{-1}a_{Ik}|_{Q_0}
       &\le \frac{M}{\lambda+\omega}|a_{Ik}|_{Q_0}.
  \end{aligned}
  \end{equation}
\end{remark}

\begin{remark}[Constant Dependency in LMMSE Expansion]
  Consequently, the exact leakage cross-covariance
  $\Sigma_{\lambda,Ik}$ has finite Cameron--Martin norm and is
  $O(\lambda^{-3})$ in that norm.  This is a condition on the deterministic
  response in the covariance norm; it does not assert that infinite-dimensional
  Gaussian sample paths lie in $H_{Q_0}$ almost surely.
\end{remark}

% ============================================================
%  §3  HILBERT-SPACE VARIATIONAL SCHUR RESIDUAL
% ============================================================
\section{Hilbert-Space Variational Schur Residual}
\label{sec:infdim}

\subsection{Variational Schur Subtraction and Norm Bounds}
\label{sec:infdim:variational}

In finite dimensions, the Schur complement of the bulk covariance block $\Sigma_{II}$ is defined as $S_k \coloneqq \Sigma_{kk} - \Sigma_{kI} \Sigma_{II}^{-1} \Sigma_{Ik}$. Under stable and controllable dynamics, the bulk covariance $\Sigma_{II}$ is positive-definite and invertible, so the subtraction is well-defined. However, in infinite dimensions, the stationary covariance operator is trace-class and compact, meaning its restriction $\Sigma_{\lambda,II}$ lacks a bounded global inverse. The appropriate operator-theoretic replacement is the shorted operator: for a non-negative block operator, the Schur complement is recovered as the quadratic form of the short to the target subspace, with the inverse term replaced by a variational supremum \cite{krein1947selfadjoint,anderson1975shorted}. To lift the Schur complement to infinite dimensions in the present singular-conditioning problem, we formulate the projection residual directly using covariance-norm forms.

Let $Q_0 \coloneqq \Sigma_{II}^{(0)}$ denote the fixed background bulk covariance under full decoupling ($\lambda = \infty$). We define the core Cameron--Martin space $H_{Q_0} \coloneqq \operatorname{Range}(Q_0^{1/2})$ and the core covariance norm:
\begin{equation}
  |u|_{Q_0} \coloneqq \|Q_0^{-1/2} u\| \quad \text{for all } u \in H_{Q_0}.
\end{equation}
By Assumption~\ref{assum:cm_stability}, the off-diagonal drift coupling vector satisfies $a_{Ik} \in H_{Q_0}$, and the restricted drift generator $\mathcal{A}_{II}$ generates an exponentially stable semigroup on $H_{Q_0}$.

Under this assumption, the exact cross-covariance resolvent formula (derived in \Cref{app:block_cumulant_analysis}) satisfies:
\begin{equation}
  \Sigma_{\lambda,Ik} = \frac{\varepsilon_0}{2\lambda^2} (\lambda I - \mathcal{A}_{II})^{-1} a_{Ik}.
\end{equation}
This resolvent response yields the Cameron--Martin norm bound on the cross-covariance block:
\begin{equation}
  \label{eq:cross_norm_bound}
  |\Sigma_{\lambda,Ik}|_{Q_0} = O(\lambda^{-3}) \quad \text{as } \lambda \to \infty.
\end{equation}

To define the infinite-dimensional projection residual, we introduce the variational Schur subtraction $R_\lambda$:
\begin{equation}
  \label{eq:variational_schur_subtraction}
  R_\lambda \coloneqq \sup_{\langle u, \Sigma_{\lambda,II} u \rangle > 0} \frac{|\langle u, \Sigma_{\lambda,Ik} \rangle|^2}{\langle u, \Sigma_{\lambda,II} u \rangle}.
\end{equation}
Equivalently, if $\Sigma_\lambda$ is viewed as a non-negative block operator on $\mathbb{C}\phi_k\oplus\mathcal{H}_I$, then \eqref{eq:variational_schur_subtraction} is the one-dimensional Anderson--Trapp shorted-operator subtraction, namely
\begin{equation}
  \label{eq:shorted_operator_scalar}
  \big\langle \phi_k,\mathcal{S}_{\mathbb{C}\phi_k}(\Sigma_\lambda)\phi_k\big\rangle
  =\Sigma_{\lambda,kk}-R_\lambda.
\end{equation}
The restriction to vectors with $\langle u,\Sigma_{\lambda,II}u\rangle>0$ simply removes null directions of the denominator; for a non-negative block operator, such directions do not change the supremum. Thus $R_\lambda$ is not an ambient inverse computation but the covariance-norm contribution removed by shorting to the target coordinate \cite{anderson1975shorted,owhadi2015conditioning}.

Since the target noise channel and bulk noise channel are independent, the bulk state decomposes into independent background and target-driven components, meaning the bulk covariance dominates the background covariance ($\Sigma_{\lambda,II} = Q_0 + \operatorname{Cov}(Y_I) \succeq Q_0$). Douglas range inclusion then guarantees that the variational subtraction is bounded:
\begin{equation}
  0 \le R_\lambda \le |\Sigma_{\lambda,Ik}|_{Q_0}^2 = O(\lambda^{-6}).
\end{equation}

\subsection{Hilbert-Space Schur Residual Scaling}
\label{sec:infdim:scaling}

We now establish the main scaling theorem for the $L^2$-projection residual.

\begin{theorem}[Hilbert Schur Residual Invariant]
  \label{thm:inf_cancellation}
  Under Assumptions~\ref{assum:linear_second_moment} and \ref{assum:cm_stability}, the $L^2$-projection residual satisfies:
  \begin{equation}
    \label{eq:main_schur_sharp}
    S_{k,\lambda} \coloneqq \|X_k - P_{\mathcal{M}_I} X_k\|_{L^2}^2 = \frac{\varepsilon_0}{2\lambda^2} - R_\lambda,
  \end{equation}
  where the remainder satisfies $0 \le R_\lambda = O(\lambda^{-6})$. In particular, the Schur LMMSE invariant is:
  \begin{equation}
    \label{eq:m2_schur_value}
    m_2^{\mathrm{Schur}} = \lim_{\lambda\to\infty} \lambda^2 S_{k,\lambda} = \frac{\varepsilon_0}{2},
  \end{equation}
  and the inverse projection residual scaling satisfies:
  \begin{equation}
    \label{eq:schur_scaling_inverse}
    S_{k,\lambda}^{-1} = \frac{2}{\varepsilon_0}\lambda^2 + O(\lambda^{-2}).
  \end{equation}
  This result depends only on covariance and orthogonal projection; it holds for any linear stationary model satisfying these assumptions, without Gaussianity.
\end{theorem}
\begin{proof}
  See \Cref{app:block_cumulant_analysis} for the full derivation.
\end{proof}

\subsection{Convergence of Finite-Dimensional Truncations}
\label{sec:proofs:convergence}

Let $\Pi_n: \mathcal{H} \to \mathcal{H}_n \coloneqq \operatorname{span}\{\phi_1, \ldots, \phi_n\}$ be the $n$-dimensional Galerkin projection. Write $\mathcal{A}_n^{(\lambda)} \coloneqq \Pi_n \mathcal{A}^{(\lambda)} \Pi_n$, $\mathcal{D}_n^{(\lambda)} \coloneqq \Pi_n \mathcal{D}^{(\lambda)} \Pi_n$, and let $\Sigma_{n,\lambda}$ solve the truncated Lyapunov equation $\mathcal{A}_n^{(\lambda)} \Sigma_{n,\lambda} + \Sigma_{n,\lambda} (\mathcal{A}_n^{(\lambda)})^* + \mathcal{D}_n^{(\lambda)} = 0$.

\begin{proposition}[Galerkin Consistency of the Schur Invariant]
  \label{prop:galerkin_consistency}
  Under Assumption~\ref{assum:linear_second_moment}, let $\Pi_n:\mathcal{H}\to\mathcal{H}_n$ be finite-rank orthogonal Galerkin projections preserving the target-bulk block decomposition:
  \begin{equation}
    \Pi_n\phi_k=\phi_k,\qquad \Pi_n\mathcal{H}_I\subset\mathcal{H}_I,\qquad \Pi_n\to I\ \text{strongly on }\mathcal{H}.
  \end{equation}
  Assume the truncated regularized dynamics satisfy the uniform exponential stability and trace-class conditions stated in Proposition~\ref{prop:dynamic_galerkin} of \Cref{app:block_cumulant_analysis}.  Then the truncated covariance operators converge in trace norm:
  \begin{equation}
    \|\Sigma_{n,\lambda}-\Sigma_\lambda\|_{\mathcal{S}_1}\to0\quad\text{as }n\to\infty,
  \end{equation}
  and the truncated Schur invariant is exact at every level:
  \begin{equation}
    m_2^{\mathrm{Schur},(n)} \coloneqq \lim_{\lambda\to\infty} \lambda^2 S_{k,\lambda}^{(n)} = \frac{\varepsilon_0}{2} \quad \text{for all } n \ge k.
  \end{equation}
\end{proposition}
\begin{proof}
  See \Cref{app:block_cumulant_analysis} for the dynamic and spectral support.
\end{proof}

% ============================================================
%  §4  GAUSSIAN CONDITIONAL EXTENSION AND FINITE-PART FUNCTIONAL
% ============================================================
\section{Gaussian Conditional Extension and Finite-Part Functional}
\label{sec:gaussian_zero}

This section gives the Gaussian extension of the base theorem, using Radon Gaussian laws and disintegration to formulate conditioning under exact evidence and define the renormalized finite-part functional.

\subsection{Radon Gaussian Setup and Disintegration}
\label{sec:gaussian:setup}

\begin{assumption}[Hilbert Gaussian Conditioning and Affine Mean Channel]
  \label{assum:gaussian_clamp}
  In addition to Assumptions~\ref{assum:linear_second_moment} and \ref{assum:cm_stability}, let $\mu_\lambda = \mathcal{N}(m_\lambda, \Sigma_\lambda)$ be a Gaussian Radon law on $\mathcal{H}$ with injective trace-class covariance. Let $\mu_{\mathrm{obs}}$ be the canonical Gaussian disintegration member on the evidence hyperplane $X_k = x_k^*$. Suppose that a fixed deterministic forcing $b \in \mathcal{H}$ gives the stationary mean equation:
  \begin{equation}
    \mathcal{A}_{\cond} m_\lambda + b - \lambda \Pi_k (m_\lambda - x_k^* \phi_k) = 0, \quad \Pi_k \mathcal{A}_{\cond} = 0.
  \end{equation}
  We define the target deterministic offset constant:
  \begin{equation}
    c_\lambda \coloneqq \lambda(m_{\lambda,k} - x_k^*) = b_k,
  \end{equation}
  where $m_{\lambda,k} \coloneqq \langle m_\lambda, \phi_k \rangle$ and $b_k \coloneqq \langle b, \phi_k \rangle$.
\end{assumption}

Under the Gaussian Radon law, the coordinate projection $\pi_k: \mathcal{H} \to \mathbb{R}$ is continuous, so the target variable $X_k$ is a well-defined Borel random variable. The Polish structure of the Hilbert space guarantees the existence of a regular conditional probability distribution. 
By selecting the weakly continuous version of the conditional law at the single point $t = x_k^*$ (as detailed in \Cref{app:gaussian_lift}), we obtain the unique pointwise disintegration member $\mu_{\mathrm{obs}}$. Under this measure, the target coordinate satisfies the exact target zeros:
\begin{equation}
  \label{eq:exact_target_zeros}
  m_{\text{obs},k} - x_k^* = 0, \quad \Var_{\text{obs}}(X_k) = 0 \quad (\text{exact}).
\end{equation}

\subsection{Global-Inverse-Free Measurable Predictor}
\label{sec:gaussian:predictor}

Let $C_\lambda \coloneqq \Sigma_{\lambda,II}$ denote the Gaussian bulk covariance operator on $\mathcal{H}_I$. Because $C_\lambda$ is trace-class and compact, it lacks a bounded global inverse. We define the bulk Cameron--Martin space $H_{C_\lambda} \coloneqq \operatorname{Range}(C_\lambda^{1/2})$ and the corresponding covariance norm.

The leakage cross-covariance vector satisfies $\Sigma_{\lambda,Ik} \in H_{C_\lambda}$ with norm $\|\Sigma_{\lambda,Ik}\|_{H_{C_\lambda}} = O(\lambda^{-3})$. This range inclusion allows us to define the conditional predictor $\mu_{k|I}(X_I) \coloneqq \E[X_k \mid X_I]$ as a measurable Wiener functional:
\begin{equation}
  \mu_{k|I}(X_I) = m_{\lambda,k} + W_{\Sigma_{\lambda,Ik}}(X_I - m_{\lambda,I}),
\end{equation}
which is the unique LMMSE predictor in $L^2$. The estimates in \Cref{app:gaussian_lift} imply that, for all sufficiently large $\lambda$, the bulk marginals $\mu_{\text{obs},I}$ and $\mu_{\lambda,I}$ are equivalent and have finite bulk relative entropy:
\begin{equation}
  \KL(\mu_{\text{obs},I} \parallel \mu_{\lambda,I}) = O(\lambda^{-4}) \to 0 \quad \text{as } \lambda \to \infty.
\end{equation}

\subsection{Gaussian Renormalized Finite-Part Functional}
\label{sec:gaussian:functional}

Although the joint KL divergence diverges under the exact conditioning, we can define a renormalized finite-part action.

\begin{definition}[Gaussian Renormalized Finite-Part Functional]
  \label{def:unified_action}
  Whenever $S_{k,\lambda}>0$ and $\mu_{\text{obs},I}$ is equivalent to $\mu_{\lambda,I}$, we define the Gaussian renormalized finite-part functional of the conditioned law:
  \begin{equation}
    \label{eq:unified_action}
    \begin{aligned}
      \mathfrak{J}_\lambda(\mu_{\mathrm{obs}}; \mu_\lambda) \coloneqq &\KL(\mu_{\mathrm{obs},I} \parallel \mu_{\lambda,I}) \\
      &+ \frac{1}{2} S_{k,\lambda}^{-1} \E_{\mu_{\mathrm{obs}}}\left[ \left( X_k - \mu_{k|I}(X_I) \right)^2 \right].
    \end{aligned}
  \end{equation}
  This represents the finite part of the joint KL divergence under the conditional normalization scheme, where the Dirac and Gaussian normalization singularities are subtracted.
\end{definition}

\paragraph{Weak scheme-independence.}
For an evidence level $t\in\mathbb{R}$, write $\mu_\lambda^t$ for the weakly continuous conditional law from \Cref{app:gaussian_lift} and set
\[
  \mathfrak{J}_\lambda(t)\coloneqq
  \mathfrak{J}_\lambda(\mu_\lambda^t;\mu_\lambda).
\]
Then, for any two evidence levels $t_1,t_2$ for which the finite part is defined,
\begin{equation}
  \label{eq:weak_scheme_independence}
  \mathfrak{J}_\lambda(t_2)-\mathfrak{J}_\lambda(t_1)
  =
  \frac{(t_2-m_{\lambda,k})^2-(t_1-m_{\lambda,k})^2}
       {2\Sigma_{\lambda,kk}}.
\end{equation}
In the calibrated scaling $\Sigma_{\lambda,kk}=\varepsilon_0/(2\lambda^2)$, if
$b_i=\lambda(m_{\lambda,k}-t_i)$, this becomes
\[
  \mathfrak{J}_\lambda(t_2)-\mathfrak{J}_\lambda(t_1)
  =
  \frac{b_2^2-b_1^2}{\varepsilon_0}.
\]
\begin{proof}
Let $\Delta_t\coloneqq t-m_{\lambda,k}$ and
\[
  \rho_\lambda
  \coloneqq
  \frac{\|\Sigma_{\lambda,Ik}\|_{H_{C_\lambda}}^2}{\Sigma_{\lambda,kk}}
  =
  1-\frac{S_{k,\lambda}}{\Sigma_{\lambda,kk}}.
\]
The Gaussian entropy computation in \Cref{app:gaussian_lift} gives
\[
  \KL(\mu_{\lambda,I}^t\parallel\mu_{\lambda,I})
  =
  \frac{1}{2}\frac{\Delta_t^2}{\Sigma_{\lambda,kk}}\rho_\lambda
  -\frac{1}{2}\log(1-\rho_\lambda)-\frac{1}{2}\rho_\lambda .
\]
The conditional-residual term in \eqref{eq:unified_action} is
\[
  \frac{1}{2S_{k,\lambda}}
  \E_{\mu_\lambda^t}\!\left[
    (X_k-\mu_{k|I}(X_I))^2
  \right]
  =
  \frac{1}{2}\rho_\lambda
  +
  \frac{1}{2}\frac{\Delta_t^2}{\Sigma_{\lambda,kk}}(1-\rho_\lambda).
\]
Adding these two identities cancels the $\rho_\lambda/2$ terms and yields
\[
  \mathfrak{J}_\lambda(t)
  =
  \frac{(t-m_{\lambda,k})^2}{2\Sigma_{\lambda,kk}}
  -\frac{1}{2}\log(1-\rho_\lambda).
\]
The logarithmic term depends on the reference covariance and the conditioning
scale, but not on the evidence value $t$. Hence it cancels in the difference
between $t_2$ and $t_1$, proving \eqref{eq:weak_scheme_independence}. The same
argument shows weak independence from the mollifying kernel: replacing the
Gaussian Dirac mollifier in \Cref{prop:mollification_limit} by any centered
approximate identity with finite second moment changes the extracted finite
part only by an additive term depending on the kernel, $\lambda$, and
$S_{k,\lambda}$, not on $t$. Therefore all evidence differences
$\Delta\mathfrak{J}_\lambda$ are kernel-independent.
\end{proof}

\subsection{Gaussian Main Theorem}
\label{sec:gaussian:theorem}

We now establish the main regularization and cancellation theorems for the finite-part functional.

\begin{theorem}[Gaussian Finite-Part Regularization]
  \label{thm:finite_part_regularization}
  Under Assumption~\ref{assum:gaussian_clamp}, there exists $\lambda_0>0$ such that, for $\lambda\ge\lambda_0$, the renormalized finite-part functional under the conditioned law $\mu_{\mathrm{obs}}$ is well-defined and satisfies:
  \begin{equation}
    \label{eq:unified_regularization}
    \mathfrak{J}_\lambda(\mu_{\mathrm{obs}}; \mu_\lambda) = \frac{b_k^2}{\varepsilon_0} + O(\lambda^{-4}).
  \end{equation}
  In particular, the functional remains bounded at $O(1)$ and converges to:
  \begin{equation}
    \lim_{\lambda\to\infty} \mathfrak{J}_\lambda(\mu_{\mathrm{obs}}; \mu_\lambda) = \frac{b_k^2}{\varepsilon_0}.
  \end{equation}
\end{theorem}
\begin{proof}
  See \Cref{app:gaussian_lift} for the full derivation.
\end{proof}

\begin{corollary}[Exact Decoupled-Centered Cancellation]
  \label{cor:exact_pole_zero}
  Under Assumption~\ref{assum:gaussian_clamp}, if the target coordinate is dynamically decoupled from the bulk ($a_{Ik} = 0$) and the evidence is centered at the stationary mean ($x_k^* = m_{\lambda,k}$), then:
  \begin{equation}
    \mathfrak{J}_\lambda(\mu_{\mathrm{obs}}; \mu_\lambda) = 0 \quad \text{for every } \lambda > 0.
  \end{equation}
\end{corollary}
\begin{proof}
  The decoupling implies $\Sigma_{\lambda,Ik} = 0$, so $\KL(\mu_{\text{obs},I} \parallel \mu_{\lambda,I}) = 0$. The centering makes the conditional residual zero under the conditioning, so both terms in \eqref{eq:unified_action} vanish.
\end{proof}

% ============================================================
%  §5  ELLIPTIC GAUSSIAN FIELD EXAMPLE
% ============================================================
\section{Elliptic Gaussian Field Example}
\label{sec:spde}

We now record a standard elliptic Gaussian field example to show that the abstract Hilbert-space assumptions are nonempty and natural for covariance operators arising from stochastic evolution equations and elliptic Gaussian priors \cite{daprato2014stochastic,lord2014introduction,stuart2010inverse}.

\subsection{Elliptic Gaussian field on $[0,1]$}

Let
\[
  L=-\frac{d^2}{dx^2}+\kappa^2
\]
on $L^2(0,1)$ with Dirichlet boundary conditions and $\kappa>0$. Then $\mathcal{A}_0=-L$ with domain $H^2(0,1)\cap H_0^1(0,1)$ is dissipative, $L$ is positive self-adjoint and sectorial, and $e^{-tL}$ is an exponentially stable analytic contraction semigroup. The inverse covariance
\[
  Q_0=L^{-1}
\]
is compact, positive, self-adjoint, and trace class. With eigenfunctions $e_j(x)=\sqrt{2}\sin(\pi jx)$, its eigenvalues are
\[
  q_j=(\pi^2j^2+\kappa^2)^{-1}.
\]
Thus $Q_0$ is a typical infinite-dimensional covariance operator: it is compact and trace class, has no bounded inverse on the ambient space, and carries its inverse only as a Cameron--Martin quadratic form \cite{bogachev1998gaussian,simon2005trace}.

\subsection{Cameron--Martin range and covariance norm}

The Cameron--Martin space associated with $Q_0$ is
\[
  H_{Q_0}=\operatorname{Range}(Q_0^{1/2})=H_0^1(0,1),
\]
with equivalent covariance norm
\[
  \|h\|_{Q_0}^2
  =\langle h,Lh\rangle_{L^2}
  =\int_0^1 \bigl(|h'(x)|^2+\kappa^2|h(x)|^2\bigr)\,dx.
\]
This identifies the variational Schur projection in a Sobolev quadratic form, as in the standard covariance-norm representation of Gaussian Hilbert priors \cite{bogachev1998gaussian,stuart2010inverse}. The covariance inverse does not act as a bounded operator on $L^2(0,1)$; it acts only through this quadratic form on its Cameron--Martin domain.

\subsection{Spectral sensor conditioning}

Fix a target eigenmode $e_k$ and define the observed coordinate
\[
  X_k=\langle X,e_k\rangle_{L^2}.
\]
The unperturbed target variance is
\[
  \Sigma_{0,kk}=q_k=(\pi^2k^2+\kappa^2)^{-1}.
\]
A regularized restoring drift on this coordinate produces the exact normalization
\[
  \Sigma_{\lambda,kk}=\frac{\varepsilon_0}{2\lambda^2}-R_\lambda,
  \qquad R_\lambda=O(\lambda^{-6}),
\]
under the block assumptions of \Cref{sec:infdim}. Off-diagonal bounded perturbations of the drift operator generate the cross-covariance vector entering the variational Schur residual. The resulting subtraction is evaluated by the Cameron--Martin quadratic form above, not by a global inverse on $L^2(0,1)$.

This example shows that the abstract theorem applies to an infinite-dimensional Gaussian field: the residual is controlled by a Sobolev norm on $H_0^1(0,1)$ while the ambient covariance remains compact and non-invertible.

\section{Discussion}
\label{sec:discussion}

We have shown that singular Gaussian coordinate conditioning in a separable Hilbert space admits an inverse-free Schur normal form.  The main obstruction is the absence of a bounded inverse for compact covariance blocks.  The variational Cameron--Martin projection is the target-coordinate shorted-operator subtraction evaluated in the covariance-energy scale, and it yields the residual estimate needed for the hard-conditioning limit.  The scalar nature of the target coordinate is essential to the present proof.  Three extensions mark the boundary between the theorem proved here and a genuinely finite-rank conditioning theory. We formalize these extensions as exact theorems below.

\subsection{Matrix algebraic extension}
\label{sec:discussion:matrix-extension}

For an $m$-dimensional observed subspace $E=\operatorname{span}\{\phi_1,\ldots,\phi_m\}$, the scalar target block $\Sigma_{\lambda,kk}$ is replaced by the covariance matrix $\Sigma_{\lambda,EE}\in\mathbb{R}^{m\times m}$. The scalar calibration becomes a matrix Lyapunov balancing problem.

\begin{theorem}[Matrix Lyapunov Calibration]
  \label{thm:matrix_lyapunov}
  Let $E$ be an $m$-dimensional target subspace. Suppose the regularized drift and diffusion on $E$ are given by $A_{\lambda,E} = -\lambda \Lambda_E$ and $D_{\lambda,E} = \frac{1}{\lambda} \Delta_E$, where $\Lambda_E, \Delta_E \in \mathbb{R}^{m \times m}$ are strictly positive definite matrices. Let $\Sigma_{\lambda,EE}$ be the unique positive definite solution to the matrix Lyapunov equation:
  \begin{equation}
    \label{eq:discussion_matrix_lyapunov}
    A_{\lambda,E}\Sigma_{\lambda,EE}
    +\Sigma_{\lambda,EE}A_{\lambda,E}^{*}
    +D_{\lambda,E}=0.
  \end{equation}
  Then $\Sigma_{\lambda,EE}$ scales exactly as $\lambda^{-2}$. Specifically, the normalized matrix $\widetilde{\Sigma}_{EE} \coloneqq \lambda^2 \Sigma_{\lambda,EE}$ is independent of $\lambda$ and is the unique positive definite solution to the scale-free Lyapunov equation:
  \begin{equation}
    \label{eq:scale_free_lyapunov}
    \Lambda_E \widetilde{\Sigma}_{EE} + \widetilde{\Sigma}_{EE} \Lambda_E^* = \Delta_E.
  \end{equation}
\end{theorem}
\begin{proof}
  Substituting $A_{\lambda,E} = -\lambda \Lambda_E$ and $D_{\lambda,E} = \frac{1}{\lambda} \Delta_E$ into \eqref{eq:discussion_matrix_lyapunov} yields:
  \[
    -\lambda \Lambda_E \Sigma_{\lambda,EE} - \lambda \Sigma_{\lambda,EE} \Lambda_E^* + \frac{1}{\lambda} \Delta_E = 0.
  \]
  Multiplying the entire equation by $\lambda$ gives:
  \[
    \Lambda_E (\lambda^2 \Sigma_{\lambda,EE}) + (\lambda^2 \Sigma_{\lambda,EE}) \Lambda_E^* = \Delta_E.
  \]
  Defining $\widetilde{\Sigma}_{EE} \coloneqq \lambda^2 \Sigma_{\lambda,EE}$, we recover \eqref{eq:scale_free_lyapunov}. Since $\Lambda_E$ is positive definite, its spectrum lies in the open right half-plane, making $-\Lambda_E$ Hurwitz. Because $\Delta_E$ is positive definite, standard continuous Lyapunov theory guarantees that \eqref{eq:scale_free_lyapunov} admits a unique positive definite solution $\widetilde{\Sigma}_{EE}$. Therefore, the exact scaling $\Sigma_{\lambda,EE} = \lambda^{-2} \widetilde{\Sigma}_{EE}$ holds for all $\lambda > 0$, properly generalizing the scalar calibration condition $\Sigma_{\lambda,kk} = \varepsilon_0 / (2\lambda^2)$.
\end{proof}

\subsection{Operator-valued residuals}
\label{sec:discussion:operator-valued-residuals}

For an $m$-dimensional target space $E$, the Schur residual must be a positive-semidefinite operator on $E$. Because the free-block covariance $\Sigma_{\lambda,II}$ is compact and lacks a bounded global inverse, the formal finite-dimensional matrix expression $R_{\lambda,E} = \Sigma_{\lambda,EI}\Sigma_{\lambda,II}^{-1}\Sigma_{\lambda,IE}$ must be rigorously replaced by the shorted-operator construction evaluated in the Cameron--Martin energy norm.

\begin{theorem}[Operator-Valued Schur Residual]
  \label{thm:operator_residual}
  Let $\mathcal{H} = E \oplus \mathcal{H}_I$, with $E$ finite-dimensional. Let $\Sigma_\lambda$ be the non-negative trace-class covariance operator on $\mathcal{H}$ with blocks $\Sigma_{\lambda,EE}$, $\Sigma_{\lambda,EI}$, $\Sigma_{\lambda,IE}$, and $\Sigma_{\lambda,II}$. Define the shorted-operator residual $R_{\lambda,E}$ on $E$ by its quadratic forms:
  \begin{equation}
    \label{eq:discussion_operator_residual_form}
    \langle v,R_{\lambda,E}v\rangle_E
    =
    \sup_{\langle u,\Sigma_{\lambda,II}u\rangle>0}
    \frac{|\langle u,\Sigma_{\lambda,IE}v\rangle|^2}
         {\langle u,\Sigma_{\lambda,II}u\rangle},
    \qquad v\in E.
  \end{equation}
  Assume the structural range inclusion $\operatorname{Range}(\Sigma_{\lambda,IE}) \subset \operatorname{Range}(\Sigma_{\lambda,II}^{1/2}) \eqqcolon H_{C_\lambda}$. Then $R_{\lambda,E}$ is a well-defined bounded positive-semidefinite operator on $E$, uniquely characterized by:
  \begin{equation}
    R_{\lambda,E} = (\Sigma_{\lambda,II}^{-1/2}\Sigma_{\lambda,IE})^* (\Sigma_{\lambda,II}^{-1/2}\Sigma_{\lambda,IE}).
  \end{equation}
  Furthermore, if the cross-covariance satisfies the energy bound $\|\Sigma_{\lambda,IE} v\|_{H_{C_\lambda}} \le M \lambda^{-3} \|v\|_E$ for all $v \in E$ and some constant $M > 0$, then the operator norm of the residual satisfies $\|R_{\lambda,E}\|_{\mathcal{L}(E)} = O(\lambda^{-6})$.
\end{theorem}
\begin{proof}
  By the assumed range inclusion, for any $v \in E$, the vector $y_v \coloneqq \Sigma_{\lambda,IE} v$ lies in $H_{C_\lambda} = \operatorname{Range}(\Sigma_{\lambda,II}^{1/2})$. Thus, there exists a unique element $w_v \in \overline{\operatorname{Range}(\Sigma_{\lambda,II}^{1/2})}$ such that $\Sigma_{\lambda,II}^{1/2} w_v = y_v$. We denote $w_v \coloneqq \Sigma_{\lambda,II}^{-1/2} \Sigma_{\lambda,IE} v$.
  
  For any $u \in \mathcal{H}_I$ with $\langle u, \Sigma_{\lambda,II} u \rangle > 0$, the ratio in \eqref{eq:discussion_operator_residual_form} can be rewritten as:
  \[
    \frac{|\langle u, \Sigma_{\lambda,II}^{1/2} w_v \rangle|^2}{\|\Sigma_{\lambda,II}^{1/2} u\|^2} = \frac{|\langle \Sigma_{\lambda,II}^{1/2} u, w_v \rangle|^2}{\|\Sigma_{\lambda,II}^{1/2} u\|^2}.
  \]
  By the Cauchy--Schwarz inequality, this ratio is bounded above by $\|w_v\|^2$. The supremum is achieved by taking a sequence of vectors $\Sigma_{\lambda,II}^{1/2} u$ converging to $w_v$, yielding:
  \[
    \langle v, R_{\lambda,E} v \rangle_E = \|w_v\|^2 = \|\Sigma_{\lambda,II}^{-1/2} \Sigma_{\lambda,IE} v\|^2 = \langle v, (\Sigma_{\lambda,II}^{-1/2}\Sigma_{\lambda,IE})^* (\Sigma_{\lambda,II}^{-1/2}\Sigma_{\lambda,IE}) v \rangle_E.
  \]
  Since this holds for all $v \in E$, it uniquely defines $R_{\lambda,E}$ as a positive-semidefinite matrix on $E$. Finally, the assumption $\|\Sigma_{\lambda,IE} v\|_{H_{C_\lambda}} \le M \lambda^{-3} \|v\|_E$ means exactly that $\|w_v\| \le M \lambda^{-3} \|v\|_E$. Therefore, $\langle v, R_{\lambda,E} v \rangle_E \le M^2 \lambda^{-6} \|v\|_E^2$, which implies the operator norm bound $\|R_{\lambda,E}\|_{\mathcal{L}(E)} \le M^2 \lambda^{-6} = O(\lambda^{-6})$.
\end{proof}

\subsection{General observation operators and misalignment}
\label{sec:discussion:misalignment}

The scalar proof in this paper uses a block decomposition aligned with a single observed coordinate. However, typical finite-rank observation operators $L:\mathcal H\to\mathbb R^m$ select a subspace $E$ that does not commute with the free prior covariance, meaning the target-free split is not an eigenbasis split. This geometric misalignment generates cross-coupling that requires a precise resolvent analysis.

\begin{theorem}[Misaligned Covariance Coupling]
  \label{thm:misaligned_coupling}
  Let $P_E$ be an orthogonal projection onto an $m$-dimensional observed subspace $E$, and let $\mathcal{H}_I = E^\perp$. Let the regularized dynamics be decomposed such that the target block matches \Cref{thm:matrix_lyapunov} with $A_{\lambda,E} = -\lambda \Lambda_E$, the cross-noise is zero ($D_{\lambda,IE} = 0$), and the off-diagonal drift coupling is denoted by $A_{\lambda,IE}$. 
  If $A_{\lambda,IE}$ maps $E$ into the Cameron--Martin space $H_{Q_0}$ of the free background covariance $Q_0$ on $\mathcal{H}_I$, and the restricted drift $A_{II}$ generates an exponentially stable semigroup on $H_{Q_0}$, then the cross-covariance satisfies the exact identity:
  \begin{equation}
    \label{eq:misaligned_resolvent_identity}
    \Sigma_{\lambda,IE} = \int_0^\infty e^{t A_{II}} A_{\lambda,IE} \Sigma_{\lambda,EE} e^{t A_{\lambda,E}^*} dt.
  \end{equation}
  Furthermore, the misaligned cross-covariance energy scales as $\|\Sigma_{\lambda,IE} v\|_{H_{Q_0}} = O(\lambda^{-3})$ for any $v \in E$.
\end{theorem}
\begin{proof}
  The off-diagonal block of the global Lyapunov equation $\mathcal{A}^{(\lambda)} \Sigma_\lambda + \Sigma_\lambda (\mathcal{A}^{(\lambda)})^* + \mathcal{D}^{(\lambda)} = 0$ yields the Sylvester equation:
  \[
    A_{II} \Sigma_{\lambda,IE} + \Sigma_{\lambda,IE} A_{\lambda,E}^* + A_{\lambda,IE} \Sigma_{\lambda,EE} = 0.
  \]
  Since $A_{II}$ generates an exponentially stable semigroup $e^{t A_{II}}$ and $A_{\lambda,E} = -\lambda \Lambda_E$ is Hurwitz, the unique solution is given by the integral representation \eqref{eq:misaligned_resolvent_identity}.
  
  To bound the Cameron--Martin norm, fix $v \in E$. We have $e^{t A_{\lambda,E}^*} v = e^{-\lambda t \Lambda_E^*} v$. By \Cref{thm:matrix_lyapunov}, $\Sigma_{\lambda,EE} = O(\lambda^{-2})$. Since $A_{\lambda,IE}$ maps into $H_{Q_0}$ and $E$ is finite-dimensional, there exists a constant $C$ such that $\|A_{\lambda,IE} w\|_{H_{Q_0}} \le C \|w\|_E$ for all $w \in E$. Using the uniform exponential stability of $e^{t A_{II}}$ on $H_{Q_0}$, namely $\|e^{t A_{II}}\|_{\mathcal{L}(H_{Q_0})} \le M e^{-\omega t}$ for some $M \ge 1, \omega > 0$, we estimate:
  \begin{align*}
    \|\Sigma_{\lambda,IE} v\|_{H_{Q_0}} 
    &\le \int_0^\infty \|e^{t A_{II}}\|_{\mathcal{L}(H_{Q_0})} \|A_{\lambda,IE} \Sigma_{\lambda,EE} e^{-\lambda t \Lambda_E^*} v\|_{H_{Q_0}} dt \\
    &\le \int_0^\infty M e^{-\omega t} C \|\Sigma_{\lambda,EE}\|_{\mathcal{L}(E)} \|e^{-\lambda t \Lambda_E^*}\|_{\mathcal{L}(E)} \|v\|_E dt.
  \end{align*}
  Let $\gamma > 0$ be the smallest eigenvalue of the positive definite matrix $\Lambda_E^*$, so $\|e^{-\lambda t \Lambda_E^*}\|_{\mathcal{L}(E)} \le K e^{-\lambda \gamma t}$. Then:
  \begin{align*}
    \|\Sigma_{\lambda,IE} v\|_{H_{Q_0}} 
    &\le M C K \|\Sigma_{\lambda,EE}\|_{\mathcal{L}(E)} \|v\|_E \int_0^\infty e^{-(\omega + \lambda \gamma) t} dt \\
    &= M C K \|\Sigma_{\lambda,EE}\|_{\mathcal{L}(E)} \|v\|_E \frac{1}{\omega + \lambda \gamma}.
  \end{align*}
  Since $\|\Sigma_{\lambda,EE}\|_{\mathcal{L}(E)} = O(\lambda^{-2})$ and $(\omega + \lambda \gamma)^{-1} = O(\lambda^{-1})$, we conclude that $\|\Sigma_{\lambda,IE} v\|_{H_{Q_0}} = O(\lambda^{-3})$. Combined with \Cref{thm:operator_residual}, this verifies that the shorted-operator residual maintains the $O(\lambda^{-6})$ scaling even in the presence of observation misalignment.
\end{proof}

The result is therefore deliberately limited to covariance and Gaussian-conditioning questions in the rank-one aligned setting.  It does not assert a general non-Gaussian conditioning theory.  For non-Gaussian stationary processes, the projection residual $S_{k,\lambda}=\varepsilon_0/(2\lambda^2)-R_\lambda$ still has a second-moment interpretation, but Radon disintegration, finite-part functionals, and measure-class behavior require separate hypotheses \cite{bogachev2007measure,kallenberg2021foundations}.  Dynamical relaxation after conditioning, or transport distances between post-conditioning laws, belongs to a different problem.  The present contribution is the operator-theoretic construction of the singular Gaussian conditioning residual itself.

% ============================================================
%  \S6  CONCLUSION
% ============================================================
\section{Conclusion}
\label{sec:conclusion}

We established a Hilbert-space framework for singular coordinate conditioning of Gaussian measures. The central object is an inverse-free variational Schur projection on the Cameron--Martin range, equivalently the one-dimensional target-coordinate form of the shorted-operator subtraction. It gives canonical normalization, an $O(\lambda^{-6})$ residual estimate, Galerkin stability, and a Radon Gaussian conditioning statement. The elliptic field example shows that the assumptions are natural for trace-class covariance operators arising from stochastic evolution equations and elliptic Gaussian priors \cite{daprato2014stochastic,lord2014introduction,stuart2010inverse}.

\section*{Declaration of Generative AI and AI-Assisted Technologies in the Manuscript Preparation Process}

During the preparation of this work, the author(s) used Google DeepMind's Antigravity assistant in order to merge the LaTeX manuscripts, clean up cross-referencing tags, format mathematical notation, and edit the draft for style consistency. After using this tool/service, the author(s) reviewed and edited the content as needed and take(s) full responsibility for the content of the published article.

% ============================================================
%  REFERENCES
% ============================================================
\bibliographystyle{amsplain}

\documentclass[11pt]{amsart}

\usepackage{amsmath}
\usepackage{amssymb}
\usepackage{amsthm}
\usepackage{mathtools}
\usepackage{enumitem}
%\usepackage{thmtools}
\usepackage{hyperref}
\hypersetup{hypertexnames=false}
\usepackage[nameinlink,noabbrev]{cleveref}
\usepackage{booktabs}
\usepackage{graphicx}
\usepackage{xcolor}


% ---- Theorem environments ----
\theoremstyle{plain}
\newtheorem{theorem}{Theorem}
\newtheorem{proposition}[theorem]{Proposition}
\newtheorem{lemma}[theorem]{Lemma}
\newtheorem{corollary}[theorem]{Corollary}
\newtheorem{claim}[theorem]{Claim}

\theoremstyle{definition}
\newtheorem{definition}[theorem]{Definition}
\newtheorem{assumption}[theorem]{Assumption}
\newtheorem{example}[theorem]{Example}

\theoremstyle{remark}
\newtheorem{remark}[theorem]{Remark}

% ---- Macros ----
\newcommand{\cond}{\operatorname{cond}}
\newcommand{\KL}{D_{\operatorname{KL}}}
\newcommand{\E}{\mathbb{E}}
\newcommand{\Var}{\operatorname{Var}}
\newcommand{\Cov}{\operatorname{Cov}}
\newcommand{\tr}{\operatorname{Tr}}
\newcommand{\diag}{\operatorname{diag}}

\newcommand{\Range}{\operatorname{Range}}
\newcommand{\Span}{\operatorname{span}}
\newcommand{\id}{\mathrm{id}}
% \newcommand{\TODO}[1]{\textcolor{red}{[TODO: #1]}} % removed for submission

% ---- Cleveref names ----
\crefname{theorem}{theorem}{theorems}
\Crefname{theorem}{Theorem}{Theorems}
\crefname{proposition}{proposition}{propositions}
\Crefname{proposition}{Proposition}{Propositions}
\crefname{lemma}{lemma}{lemmas}
\Crefname{lemma}{Lemma}{Lemmas}
\crefname{corollary}{corollary}{corollaries}
\Crefname{corollary}{Corollary}{Corollaries}
\crefname{claim}{claim}{claims}
\Crefname{claim}{Claim}{Claims}
\crefname{definition}{definition}{definitions}
\Crefname{definition}{Definition}{Definitions}
\crefname{assumption}{assumption}{assumptions}
\Crefname{assumption}{Assumption}{Assumptions}
\crefname{example}{example}{examples}
\Crefname{example}{Example}{Examples}
\crefname{remark}{remark}{remarks}
\Crefname{remark}{Remark}{Remarks}
\crefname{equation}{Eq.}{Eqs.}
\Crefname{equation}{Equation}{Equations}
\crefname{figure}{fig.}{figs.}
\Crefname{figure}{Fig.}{Figs.}
\crefname{table}{table}{tables}
\Crefname{table}{Table}{Tables}
\crefname{section}{Sec.}{Secs.}
\Crefname{section}{Section}{Sections}

\title[Singular Gaussian Conditioning in Hilbert Space]{Singular Gaussian Conditioning in Hilbert Space: Inverse-Free Schur Projection and Variance Collapse}

\author{Lucia Shang}
\thanks{Manuscript received May 2026. Repository information omitted for submission.}

\begin{document}

% ============================================================
%  ABSTRACT
% ============================================================
\begin{abstract}
We study singular coordinate conditioning for Gaussian measures on separable Hilbert spaces. Since trace-class covariance blocks are compact, the usual Schur-complement formula cannot be implemented through a bounded inverse. We introduce an inverse-free variational Schur projection on the Cameron--Martin range and prove canonical normalization together with an $O(\lambda^{-6})$ residual estimate. The construction is stable under Galerkin truncation for exponentially stable analytic semigroup dynamics and applies to sectorial elliptic Gaussian fields.
\end{abstract}

\keywords{Singular conditioning; Gaussian measures; Cameron--Martin space; Schur complement; Hilbert-space covariance operators; trace-class covariance; Radon Gaussian conditioning; Galerkin approximation; stochastic evolution equations}

\maketitle

% ============================================================
%  \S1  INTRODUCTION
% ============================================================
\section{Introduction}
\label{sec:intro}

Singular coordinate conditioning arises when an exact constraint $X_k=x_k^*$ is reached as a zero-noise or infinite-precision limit of Gaussian conditioning problems.  Finite-dimensional Gaussian conditioning is governed by the ordinary Schur complement of a covariance matrix \cite{horn2012matrix,rasmussen2006gaussian}.  Infinite-dimensional Gaussian conditioning is subtler: Radon Gaussian measures on separable Hilbert or Banach spaces have compact covariance operators, and their Cameron--Martin spaces are covariance-norm domains rather than ambient statistical parameter spaces \cite{bogachev1998gaussian,daprato2014stochastic,steinwart2024conditioning}.  In Hilbert-space Gaussian conditioning, shorted operators give the covariance of conditioning on closed subspaces, while recent inverse-problem work emphasizes that positive-noise posterior formulas must be interpreted through approximation, disintegration, or Hilbert-space posterior equivalence rather than through a global covariance inverse \cite{owhadi2015conditioning,stuart2010inverse,chen2024gaussian,carerelie2024optimal,carerelie2025posterior}.

The obstruction studied here is the covariance-block obstruction behind that phenomenon.  In a finite-dimensional model, the conditional variance of a coordinate can be written as $A-BD^{-1}B^*$ after a block decomposition.  In a separable Hilbert space the corresponding block $D$ is typically compact and trace class, hence $D^{-1}$ is not a bounded operator on the ambient space \cite{bogachev1998gaussian,simon2005trace,daprato2014stochastic}.  Classical Schur-complement and shorted-operator theory, originating in Krein's theory of non-negative extensions and developed systematically by Anderson and Trapp, provides variational tools for positive block operators \cite{krein1947selfadjoint,anderson1975shorted,douglas1966majorization}.  However, this theory does not by itself give the scaling law for a singular Gaussian conditioning residual \cite{tretter2008spectral,contino2018schur,friedrich2017generalized,hassi2023complementation}.  This paper formulates the residual directly as a Cameron--Martin norm projection and never invokes a pseudoinverse or an ambient precision operator.

The hard-conditioning limit also reaches the Feldman--H\'{a}jek measure-class boundary.  With positive observational noise, Gaussian Bayesian inverse problems may remain in the equivalent-measure regime \cite{stuart2010inverse,carerelie2024optimal,carerelie2025posterior}.  Under an exact coordinate constraint, the limiting conditional law is supported on a hyperplane of prior measure zero, so the equivalent-measure regime is left \cite{feldman1958equivalence,hajek1958property,bogachev1998gaussian}.  A useful theory must therefore separate the diverging joint divergence caused by exact evidence from the finite covariance residual that remains visible in the block structure.

We prove that this finite residual is controlled by an inverse-free variational Schur projection.  Under the Lyapunov block assumptions stated below, the target residual has the exact form
\[
  S_{k,\lambda}=\frac{\varepsilon_0}{2\lambda^2}-R_\lambda,
  \qquad
  R_\lambda=O(\lambda^{-6}).
\]
Consequently,
\[
  S_{k,\lambda}^{-1}=\frac{2}{\varepsilon_0}\lambda^2+O(\lambda^{-2}).
\]
The proof uses covariance-norm forms, shorting, Douglas range inclusion, Lyapunov block identities, and trace-ideal convergence. These are standard inputs in operator and SPDE analysis \cite{krein1947selfadjoint,anderson1975shorted,douglas1966majorization} \cite{simon2005trace,daprato2014stochastic}. The Cameron--Martin space is used only as a Hilbert subspace equipped with its covariance norm; no differentiable statistical or geometric structure is assumed \cite{bogachev1998gaussian}.
The operator base is an exponentially stable analytic semigroup with bounded drift perturbations, so dissipative sectorial generators such as $-(-\Delta+\kappa^2)$ fall inside the abstract setup.

The paper has four main contributions:
\begin{enumerate}
  \item \textbf{Canonical normalization.} We identify the block-Lyapunov mechanism that fixes the observed coordinate variance at order $\lambda^{-2}$ with explicit coefficient $\varepsilon_0/2$.
  \item \textbf{Inverse-free Schur projection.} We identify the residual subtraction with the target-coordinate quadratic form of the shorted-operator variational formula, then prove the $O(\lambda^{-6})$ residual estimate in the same operator-theoretic regime where compact covariance blocks preclude a bounded Schur inverse.
  \item \textbf{Analytic-semigroup and Galerkin stability.} We work with an analytic semigroup drift base plus bounded perturbations, and prove that finite-dimensional truncations converge to the Hilbert-space residual in trace ideal norm, matching the discretization-independent perspective used in Hilbert-space Gaussian inverse problems \cite{stuart2010inverse,carerelie2024optimal,carerelie2025posterior}.
  \item \textbf{Radon Gaussian conditioning.} For Radon Gaussian laws, we construct the conditional predictor as a measurable Cameron--Martin projection and show that the associated renormalized conditioning functional remains bounded, using regular conditional probabilities and continuous Gaussian disintegration \cite{bogachev1998gaussian,lagatta2013continuous,steinwart2024conditioning}.
\end{enumerate}

\subsection{Positioning in the literature}
\label{sec:related}

The paper is closest to four bodies of work.  First, it uses the classical measure-class structure of Gaussian measures, especially the Cameron--Martin theorem and the Feldman--H\'{a}jek dichotomy, as presented in standard references on Gaussian measures \cite{feldman1958equivalence,hajek1958property,bogachev1998gaussian}.  Second, it is adjacent to Bayesian inverse problems and Hilbert-space Gaussian conditioning, including the shorted-operator description of Gaussian conditioning on closed subspaces, positive-noise posterior formulas, low-rank posterior approximations, and nonlinear observation maps \cite{owhadi2015conditioning,stuart2010inverse,chen2024gaussian,steinwart2024conditioning,carerelie2024optimal,carerelie2025posterior}.  Third, the variational residual is informed by Schur complements, shorted operators, Douglas range inclusion, and modern complementability theory for block operators \cite{krein1947selfadjoint,anderson1975shorted,douglas1966majorization,tretter2008spectral,contino2018schur,friedrich2017generalized,hassi2023complementation}.  Fourth, the motivating examples come from stochastic evolution equations and elliptic Gaussian fields, where dissipative sectorial generators produce compact trace-class stationary covariances that are naturally approximated by Galerkin truncation \cite{daprato2014stochastic,lord2014introduction,simon2005trace}.

\subsection{Organization of the paper}
\label{sec:organization}

\Cref{sec:setup} defines the operator model and establishes the canonical normalization. \Cref{sec:infdim} establishes the Hilbert-space variational Schur residual. \Cref{sec:gaussian_zero} gives the Radon Gaussian conditional extension and the renormalized conditioning functional. \Cref{sec:spde} gives an elliptic Gaussian-field example. \Cref{sec:discussion} records limitations and extensions, and \Cref{sec:conclusion} concludes. Proofs are given in \Cref{app:block_cumulant_analysis,app:gaussian_lift}; numerical diagnostics and Galerkin verification are given in \Cref{app:experiments}.

\section{Operator Model and Canonical Normalization}
\label{sec:setup}

\subsection{Operator Dynamics and Regularized Conditioning}
\label{sec:setup:dynamics}

We model the continuous-time dynamics as stochastic evolution on a separable Hilbert space $\mathcal{H}$ \cite{daprato2014stochastic}:
\begin{equation}
  \label{eq:stochastic_evolution}
  dX_t = \mathcal{A} X_t \, dt + \mathcal{D}^{1/2} dW_t
\end{equation}
where $\mathcal{A}$ is the drift operator and $\mathcal{D}$ is the trace-class diffusion operator. We decompose the Hilbert space as $\mathcal{H} = \mathcal{H}_k \oplus \mathcal{H}_I$, where $\mathcal{H}_k \coloneqq \mathbb{R}\phi_k$ is the one-dimensional target coordinate subspace, and $\mathcal{H}_I \coloneqq \mathcal{H}_k^\perp$ is the free bulk subspace. Let $\Pi_k \coloneqq \phi_k \otimes \phi_k$ and $\Pi_I \coloneqq I - \Pi_k$ be the corresponding orthogonal projections.

To approximate the exact point conditioning $X_k = x_k^*$, we set the incoming target-coordinate blocks to zero and apply a regularization parameter $\lambda > 0$.

\begin{definition}[Calibrated Approximating Dynamics Family]
  \label{def:soft_do}
  The regularized conditioning operator of strength $\lambda > 0$ at coordinate $X_k$ acts on the operator dynamics \eqref{eq:stochastic_evolution} by modifying both the drift and the diffusion:
  \begin{equation}
    \label{eq:soft_do_sde}
    \mathcal{A}^{(\lambda)} \coloneqq \mathcal{A}_{\cond} - \lambda \Pi_k, \qquad
    \mathcal{D}^{(\lambda)} \coloneqq \mathcal{D}_{\cond} + d_\lambda \Pi_k
  \end{equation}
  where $\mathcal{A}_{\cond}$ is the block drift operator with incoming target-coordinate blocks set to zero ($\Pi_k \mathcal{A}_{\cond} = 0$), $\mathcal{D}_{\cond}$ has target-coordinate cross-noise set to zero, and $d_\lambda > 0$ is the local diffusion variance of the target coordinate at regularization strength $\lambda$. In the SDE, this corresponds to:
  \begin{equation}
    \label{eq:soft_do_sde_general}
    dX_{k,t} = -\lambda (X_{k,t} - x_k^*) dt + \sqrt{d_\lambda}\, dW_{k,t}.
  \end{equation}
\end{definition}

\subsection{Canonical Normalization}
\label{sec:setup:calibration}

The block-decoupled structure under the regularized dynamics yields:
\begin{equation}
  \mathcal{A}^{(\lambda)}=
  \begin{pmatrix}-\lambda&0\\a_{Ik}&\mathcal{A}_{II}\end{pmatrix},\qquad
  \mathcal{D}^{(\lambda)}=
  \begin{pmatrix}d_\lambda&0\\0&D_{II}\end{pmatrix},
\end{equation}
where $a_{Ik} \coloneqq \mathcal{A}_{Ik} \phi_k \in \mathcal{H}_I$ represents the off-diagonal drift coupling.
Let $\Sigma_\lambda$ denote the stationary covariance operator under regularization strength $\lambda$, solving the Lyapunov equation:
\begin{equation}
  \mathcal{A}^{(\lambda)}\Sigma_\lambda + \Sigma_\lambda(\mathcal{A}^{(\lambda)})^* + \mathcal{D}^{(\lambda)} = 0.
\end{equation}
The target block $kk$ decouples completely from the free coordinates, yielding the exact scalar equation:
\begin{equation}
  -2\lambda\Sigma_{\lambda,kk} + d_\lambda = 0 \quad \Longrightarrow \quad \Sigma_{\lambda,kk} = \frac{d_\lambda}{2\lambda}.
\end{equation}

\begin{proposition}[Canonical Scaling]
  \label{prop:exact_calibration}
  To maintain a target noise scale $\varepsilon_0 > 0$ independent of the auxiliary parameter $\lambda$ under the singular limit, we define the canonical scaling condition:
  \begin{equation}
    \lambda^2 \Sigma_{\lambda,kk} \equiv \frac{\varepsilon_0}{2} \quad \text{for all } \lambda > 0,
  \end{equation}
  which uniquely requires:
  \begin{equation}
    d_\lambda = \frac{\varepsilon_0}{\lambda} \quad \text{for every } \lambda > 0.
  \end{equation}
  Substitution of this target diffusion scaling yields the calibrated approximating dynamics:
  \begin{equation}
    dX_{k,t} = -\lambda (X_{k,t} - x_k^*) dt + \sqrt{\varepsilon_0 / \lambda}\, dW_{k,t}.
  \end{equation}
\end{proposition}

All subsequent results are stated for this calibrated family.

\begin{definition}[Schur LMMSE Invariant]
  \label{def:schur_invariant}
  Let $S_{k,\lambda} \coloneqq \|X_k - P_{\mathcal{M}_I} X_k\|_{L^2}^2$ represent the linear minimum mean squared error (LMMSE) residual of the target coordinate $X_k$ on the closed linear span $\mathcal{M}_I$ of the free coordinates $X_I$. The \emph{Schur LMMSE invariant} $m_2^{\mathrm{Schur}}$ is defined as:
  \begin{equation}
    \label{eq:m2_schur_definition}
    m_2^{\mathrm{Schur}} \coloneqq \lim_{\lambda\to\infty} \lambda^2 S_{k,\lambda}.
  \end{equation}
\end{definition}

\subsection{Operator Setup and Assumptions}
\label{sec:setup:assumptions}

We state the operator-theoretic assumptions that make the limit passage independent of the Galerkin dimension.

\begin{definition}[Drift Operator with Analytic Base]
  \label{def:drift_operator}
  A drift operator on a separable Hilbert space $\mathcal{H}$ is an operator of the form $\mathcal{A} = \mathcal{A}_0 + \mathcal{K}$, where:
  \begin{enumerate}
    \item[(i)] $\mathcal{A}_0: \operatorname{dom}(\mathcal{A}_0) \subset \mathcal{H} \to \mathcal{H}$ is closed, densely defined, and dissipative, and it generates a strongly continuous analytic contraction semigroup $(e^{t\mathcal{A}_0})_{t \ge 0}$ on $\mathcal{H}$ satisfying $\|e^{t\mathcal{A}_0}\| \le e^{-\alpha t}$ for some spectral gap $\alpha > 0$; equivalently, the base operator $-\mathcal{A}_0$ is sectorial of angle $<\pi/2$ in this exponentially stable contraction setting.
    \item[(ii)] $\mathcal{K}: \mathcal{H} \to \mathcal{H}$ is a bounded linear operator.
    \item[(iii)] The full operator $\mathcal{A} = \mathcal{A}_0 + \mathcal{K}$, with $\operatorname{dom}(\mathcal{A})=\operatorname{dom}(\mathcal{A}_0)$, generates a strongly continuous analytic semigroup $(e^{t\mathcal{A}})_{t \ge 0}$ by the bounded perturbation theorem, and this semigroup remains exponentially stable: $\|e^{t\mathcal{A}}\| \le M e^{-\beta t}$ for some $M \ge 1$, $\beta > 0$.
  \end{enumerate}
  In finite dimensions ($\mathcal{H} = \mathbb{R}^n$), this reduces to a decomposition of a Hurwitz drift matrix into an exponentially stable analytic base plus a bounded perturbation.
\end{definition}

\begin{definition}[Trace-Class Noise and Covariance]
  \label{def:trace_class}
  Let $\mathcal{S}_1(\mathcal{H})$ denote the Banach space of trace-class (nuclear) operators on $\mathcal{H}$ equipped with the trace norm $\|T\|_{\mathcal{S}_1} \coloneqq \tr(|T|) = \sum_{i=1}^\infty \langle \phi_i, (T^*T)^{1/2} \phi_i \rangle < \infty$.
  \begin{enumerate}
    \item[(i)] The diffusion operator $\mathcal{D}: \mathcal{H} \to \mathcal{H}$ is positive, self-adjoint, and belongs to $\mathcal{S}_1(\mathcal{H})$.
    \item[(ii)] Assuming the pair $(\mathcal{A}, \mathcal{D}^{1/2})$ is controllable, the stationary covariance operator $\Sigma: \mathcal{H} \to \mathcal{H}$ is the unique strictly positive, self-adjoint, trace-class solution to the operator Lyapunov equation:
    \begin{equation}
      \label{eq:op_lyapunov}
      \mathcal{A}\Sigma + \Sigma\mathcal{A}^* + \mathcal{D} = 0
    \end{equation}
    which is given by the convergent Bochner integral
    \[
      \Sigma = \int_0^\infty e^{t\mathcal{A}} \mathcal{D} e^{t\mathcal{A}^*} dt.
    \]
  \end{enumerate}
\end{definition}

\begin{assumption}[Core Analytical Assumptions for Singular Hyperplane Conditioning]
  This work relies on two core analytical assumptions regarding the target-coordinate block decoupling and the outgoing coupling:
  \begin{enumerate}
    \item[(i)] \textbf{Linear Second-Moment Channel:} \label{assum:linear_second_moment} A centered linear stationary process on $\mathcal{H}=\mathcal{H}_k\oplus \mathcal{H}_I$ has finite second moments and covariance solving \Cref{eq:op_lyapunov}. The block-decoupling constraint zeroes all incoming target drift and cross-noise, so that $\mathcal{A}_{kI}=\mathcal{D}_{kI}=\mathcal{D}_{Ik}=0$ and $\mathcal{A}_{kk}=-\lambda$. The target diffusion $\mathcal{D}_{kk}=\varepsilon_0/\lambda$ is the unique scaling determined by the canonical normalization (Proposition~\ref{prop:exact_calibration}). The free noise is independent of the local target noise. (Gaussianity is not assumed).
    \item[(ii)] \textbf{Cameron--Martin Stability of Off-Diagonal Drift Coupling:} \label{assum:cm_stability} Let $Q_0\coloneqq\Sigma_{II}^{(0)}$ denote the free-block covariance produced by the free noise under block decoupling, and let $H_{Q_0}\coloneqq\operatorname{Range}(Q_0^{1/2})$ with the covariance norm $|u|_{Q_0}\coloneqq\|Q_0^{-1/2}u\|_{\mathcal{H}_I}$. The outgoing vector $a_{Ik}\coloneqq\mathcal{A}_{Ik}\phi_k$ belongs to $H_{Q_0}$, and $e^{t\mathcal{A}_{II}}$ restricts to a strongly continuous exponentially stable semigroup on $H_{Q_0}$:
    \begin{equation}
      \|e^{t\mathcal{A}_{II}}\|_{\mathcal{L}(H_{Q_0})} \le M e^{-\omega t},\qquad t\ge0,
    \end{equation}
    for constants $M\ge1$ and $\omega>0$ independent of $\lambda$.
  \end{enumerate}
  \vspace{1mm}
  \noindent\textit{Intuitive Explanation:} Part (i) states that the regularized dynamics act as a strong localized drift penalization on the target coordinate $X_k$ while isolating it from incoming influences, giving an operator implementation of the block-decoupling constraint. Part (ii) requires that the off-diagonal drift coupling $a_{Ik}$ from the target coordinate has finite covariance norm relative to the background fluctuations of the free coordinates. This ensures that target fluctuations propagating downstream do not overload the variance of the infinite-dimensional system.
\end{assumption}

\begin{remark}[Finite-Norm Outgoing Covariance Coupling and Finite-Dimensional Calibration]
  \label{rem:cm_information_flow}
  The membership $a_{Ik}\in H_{Q_0}$ is an admissibility condition,
  not a consequence of the free-block Lyapunov equation alone: an arbitrary
  outgoing vector need not lie in
  $\operatorname{Range}((\Sigma_{II}^{(0)})^{1/2})$.  It says that the
  deterministic channel carrying target fluctuations into the free-block has
  finite covariance norm relative to the free-noise background.
  
  To understand this condition intuitively, consider a finite-dimensional setting where the free-block state space is $\mathcal{H}_I = \mathbb{R}^{d-1}$ and the free stationary covariance $Q_0 \coloneqq \Sigma_{II}^{(0)}$ is strictly positive-definite. Since $Q_0 \succ 0$, its square root $Q_0^{1/2}$ is invertible, and its range is the entire space $\mathbb{R}^{d-1}$. In this case, the Cameron--Martin space is the full space $H_{Q_0} = \operatorname{Range}(Q_0^{1/2}) = \mathbb{R}^{d-1}$, and the Cameron--Martin norm $|u|_{Q_0} = \|Q_0^{-1/2} u\|_2$ is equivalent to the standard Euclidean norm. Consequently, the admissibility condition $a_{Ik} \in H_{Q_0}$ is trivially satisfied for any outgoing vector $a_{Ik} \in \mathbb{R}^{d-1}$.
  
  However, if some direction in the free coordinates is noise-free under block decoupling (meaning $Q_0$ is positive semidefinite but singular), $H_{Q_0}$ becomes a proper subspace of $\mathbb{R}^{d-1}$. In this case, the condition $a_{Ik} \in H_{Q_0}$ requires that the leakage vector $a_{Ik}$ has no component in the null space of $Q_0$. Equivalently, it prevents the target coordinate from leaking fluctuations into a downstream coordinate that has zero background noise, which would otherwise result in an infinite relative covariance norm.
  
  In infinite dimensions, $\Sigma_{II}^{(0)}$ is trace-class and compact, so its eigenvalues decay to zero, making $\operatorname{Range}((\Sigma_{II}^{(0)})^{1/2})$ a dense but strictly proper subspace of $\mathcal{H}_I$. The condition $a_{Ik} \in H_{Q_0}$ requires the off-diagonal drift coupling coefficients to decay sufficiently fast relative to the background noise eigenvalues to keep the relative covariance norm finite.
  
  Under this assumption, invariance and exponential stability on $H_{Q_0}$ imply:
  \begin{equation}
  \begin{aligned}
    (\lambda I-\mathcal{A}_{II})^{-1}a_{Ik}
       &\in H_{Q_0},\\
    |(\lambda I-\mathcal{A}_{II})^{-1}a_{Ik}|_{Q_0}
       &\le \frac{M}{\lambda+\omega}|a_{Ik}|_{Q_0}.
  \end{aligned}
  \end{equation}
\end{remark}

\begin{remark}[Constant Dependency in LMMSE Expansion]
  Consequently, the exact leakage cross-covariance
  $\Sigma_{\lambda,Ik}$ has finite Cameron--Martin norm and is
  $O(\lambda^{-3})$ in that norm.  This is a condition on the deterministic
  response in the covariance norm; it does not assert that infinite-dimensional
  Gaussian sample paths lie in $H_{Q_0}$ almost surely.
\end{remark}

% ============================================================
%  §3  HILBERT-SPACE VARIATIONAL SCHUR RESIDUAL
% ============================================================
\section{Hilbert-Space Variational Schur Residual}
\label{sec:infdim}

\subsection{Variational Schur Subtraction and Norm Bounds}
\label{sec:infdim:variational}

In finite dimensions, the Schur complement of the bulk covariance block $\Sigma_{II}$ is defined as $S_k \coloneqq \Sigma_{kk} - \Sigma_{kI} \Sigma_{II}^{-1} \Sigma_{Ik}$. Under stable and controllable dynamics, the bulk covariance $\Sigma_{II}$ is positive-definite and invertible, so the subtraction is well-defined. However, in infinite dimensions, the stationary covariance operator is trace-class and compact, meaning its restriction $\Sigma_{\lambda,II}$ lacks a bounded global inverse. The appropriate operator-theoretic replacement is the shorted operator: for a non-negative block operator, the Schur complement is recovered as the quadratic form of the short to the target subspace, with the inverse term replaced by a variational supremum \cite{krein1947selfadjoint,anderson1975shorted}. To lift the Schur complement to infinite dimensions in the present singular-conditioning problem, we formulate the projection residual directly using covariance-norm forms.

Let $Q_0 \coloneqq \Sigma_{II}^{(0)}$ denote the fixed background bulk covariance under full decoupling ($\lambda = \infty$). We define the core Cameron--Martin space $H_{Q_0} \coloneqq \operatorname{Range}(Q_0^{1/2})$ and the core covariance norm:
\begin{equation}
  |u|_{Q_0} \coloneqq \|Q_0^{-1/2} u\| \quad \text{for all } u \in H_{Q_0}.
\end{equation}
By Assumption~\ref{assum:cm_stability}, the off-diagonal drift coupling vector satisfies $a_{Ik} \in H_{Q_0}$, and the restricted drift generator $\mathcal{A}_{II}$ generates an exponentially stable semigroup on $H_{Q_0}$.

Under this assumption, the exact cross-covariance resolvent formula (derived in \Cref{app:block_cumulant_analysis}) satisfies:
\begin{equation}
  \Sigma_{\lambda,Ik} = \frac{\varepsilon_0}{2\lambda^2} (\lambda I - \mathcal{A}_{II})^{-1} a_{Ik}.
\end{equation}
This resolvent response yields the Cameron--Martin norm bound on the cross-covariance block:
\begin{equation}
  \label{eq:cross_norm_bound}
  |\Sigma_{\lambda,Ik}|_{Q_0} = O(\lambda^{-3}) \quad \text{as } \lambda \to \infty.
\end{equation}

To define the infinite-dimensional projection residual, we introduce the variational Schur subtraction $R_\lambda$:
\begin{equation}
  \label{eq:variational_schur_subtraction}
  R_\lambda \coloneqq \sup_{\langle u, \Sigma_{\lambda,II} u \rangle > 0} \frac{|\langle u, \Sigma_{\lambda,Ik} \rangle|^2}{\langle u, \Sigma_{\lambda,II} u \rangle}.
\end{equation}
Equivalently, if $\Sigma_\lambda$ is viewed as a non-negative block operator on $\mathbb{C}\phi_k\oplus\mathcal{H}_I$, then \eqref{eq:variational_schur_subtraction} is the one-dimensional Anderson--Trapp shorted-operator subtraction, namely
\begin{equation}
  \label{eq:shorted_operator_scalar}
  \big\langle \phi_k,\mathcal{S}_{\mathbb{C}\phi_k}(\Sigma_\lambda)\phi_k\big\rangle
  =\Sigma_{\lambda,kk}-R_\lambda.
\end{equation}
The restriction to vectors with $\langle u,\Sigma_{\lambda,II}u\rangle>0$ simply removes null directions of the denominator; for a non-negative block operator, such directions do not change the supremum. Thus $R_\lambda$ is not an ambient inverse computation but the covariance-norm contribution removed by shorting to the target coordinate \cite{anderson1975shorted,owhadi2015conditioning}.

Since the target noise channel and bulk noise channel are independent, the bulk state decomposes into independent background and target-driven components, meaning the bulk covariance dominates the background covariance ($\Sigma_{\lambda,II} = Q_0 + \operatorname{Cov}(Y_I) \succeq Q_0$). Douglas range inclusion then guarantees that the variational subtraction is bounded:
\begin{equation}
  0 \le R_\lambda \le |\Sigma_{\lambda,Ik}|_{Q_0}^2 = O(\lambda^{-6}).
\end{equation}

\subsection{Hilbert-Space Schur Residual Scaling}
\label{sec:infdim:scaling}

We now establish the main scaling theorem for the $L^2$-projection residual.

\begin{theorem}[Hilbert Schur Residual Invariant]
  \label{thm:inf_cancellation}
  Under Assumptions~\ref{assum:linear_second_moment} and \ref{assum:cm_stability}, the $L^2$-projection residual satisfies:
  \begin{equation}
    \label{eq:main_schur_sharp}
    S_{k,\lambda} \coloneqq \|X_k - P_{\mathcal{M}_I} X_k\|_{L^2}^2 = \frac{\varepsilon_0}{2\lambda^2} - R_\lambda,
  \end{equation}
  where the remainder satisfies $0 \le R_\lambda = O(\lambda^{-6})$. In particular, the Schur LMMSE invariant is:
  \begin{equation}
    \label{eq:m2_schur_value}
    m_2^{\mathrm{Schur}} = \lim_{\lambda\to\infty} \lambda^2 S_{k,\lambda} = \frac{\varepsilon_0}{2},
  \end{equation}
  and the inverse projection residual scaling satisfies:
  \begin{equation}
    \label{eq:schur_scaling_inverse}
    S_{k,\lambda}^{-1} = \frac{2}{\varepsilon_0}\lambda^2 + O(\lambda^{-2}).
  \end{equation}
  This result depends only on covariance and orthogonal projection; it holds for any linear stationary model satisfying these assumptions, without Gaussianity.
\end{theorem}
\begin{proof}
  See \Cref{app:block_cumulant_analysis} for the full derivation.
\end{proof}

\subsection{Convergence of Finite-Dimensional Truncations}
\label{sec:proofs:convergence}

Let $\Pi_n: \mathcal{H} \to \mathcal{H}_n \coloneqq \operatorname{span}\{\phi_1, \ldots, \phi_n\}$ be the $n$-dimensional Galerkin projection. Write $\mathcal{A}_n^{(\lambda)} \coloneqq \Pi_n \mathcal{A}^{(\lambda)} \Pi_n$, $\mathcal{D}_n^{(\lambda)} \coloneqq \Pi_n \mathcal{D}^{(\lambda)} \Pi_n$, and let $\Sigma_{n,\lambda}$ solve the truncated Lyapunov equation $\mathcal{A}_n^{(\lambda)} \Sigma_{n,\lambda} + \Sigma_{n,\lambda} (\mathcal{A}_n^{(\lambda)})^* + \mathcal{D}_n^{(\lambda)} = 0$.

\begin{proposition}[Galerkin Consistency of the Schur Invariant]
  \label{prop:galerkin_consistency}
  Under Assumption~\ref{assum:linear_second_moment}, let $\Pi_n:\mathcal{H}\to\mathcal{H}_n$ be finite-rank orthogonal Galerkin projections preserving the target-bulk block decomposition:
  \begin{equation}
    \Pi_n\phi_k=\phi_k,\qquad \Pi_n\mathcal{H}_I\subset\mathcal{H}_I,\qquad \Pi_n\to I\ \text{strongly on }\mathcal{H}.
  \end{equation}
  Assume the truncated regularized dynamics satisfy the uniform exponential stability and trace-class conditions stated in Proposition~\ref{prop:dynamic_galerkin} of \Cref{app:block_cumulant_analysis}.  Then the truncated covariance operators converge in trace norm:
  \begin{equation}
    \|\Sigma_{n,\lambda}-\Sigma_\lambda\|_{\mathcal{S}_1}\to0\quad\text{as }n\to\infty,
  \end{equation}
  and the truncated Schur invariant is exact at every level:
  \begin{equation}
    m_2^{\mathrm{Schur},(n)} \coloneqq \lim_{\lambda\to\infty} \lambda^2 S_{k,\lambda}^{(n)} = \frac{\varepsilon_0}{2} \quad \text{for all } n \ge k.
  \end{equation}
\end{proposition}
\begin{proof}
  See \Cref{app:block_cumulant_analysis} for the dynamic and spectral support.
\end{proof}

% ============================================================
%  §4  GAUSSIAN CONDITIONAL EXTENSION AND FINITE-PART FUNCTIONAL
% ============================================================
\section{Gaussian Conditional Extension and Finite-Part Functional}
\label{sec:gaussian_zero}

This section gives the Gaussian extension of the base theorem, using Radon Gaussian laws and disintegration to formulate conditioning under exact evidence and define the renormalized finite-part functional.

\subsection{Radon Gaussian Setup and Disintegration}
\label{sec:gaussian:setup}

\begin{assumption}[Hilbert Gaussian Conditioning and Affine Mean Channel]
  \label{assum:gaussian_clamp}
  In addition to Assumptions~\ref{assum:linear_second_moment} and \ref{assum:cm_stability}, let $\mu_\lambda = \mathcal{N}(m_\lambda, \Sigma_\lambda)$ be a Gaussian Radon law on $\mathcal{H}$ with injective trace-class covariance. Let $\mu_{\mathrm{obs}}$ be the canonical Gaussian disintegration member on the evidence hyperplane $X_k = x_k^*$. Suppose that a fixed deterministic forcing $b \in \mathcal{H}$ gives the stationary mean equation:
  \begin{equation}
    \mathcal{A}_{\cond} m_\lambda + b - \lambda \Pi_k (m_\lambda - x_k^* \phi_k) = 0, \quad \Pi_k \mathcal{A}_{\cond} = 0.
  \end{equation}
  We define the target deterministic offset constant:
  \begin{equation}
    c_\lambda \coloneqq \lambda(m_{\lambda,k} - x_k^*) = b_k,
  \end{equation}
  where $m_{\lambda,k} \coloneqq \langle m_\lambda, \phi_k \rangle$ and $b_k \coloneqq \langle b, \phi_k \rangle$.
\end{assumption}

Under the Gaussian Radon law, the coordinate projection $\pi_k: \mathcal{H} \to \mathbb{R}$ is continuous, so the target variable $X_k$ is a well-defined Borel random variable. The Polish structure of the Hilbert space guarantees the existence of a regular conditional probability distribution. 
By selecting the weakly continuous version of the conditional law at the single point $t = x_k^*$ (as detailed in \Cref{app:gaussian_lift}), we obtain the unique pointwise disintegration member $\mu_{\mathrm{obs}}$. Under this measure, the target coordinate satisfies the exact target zeros:
\begin{equation}
  \label{eq:exact_target_zeros}
  m_{\text{obs},k} - x_k^* = 0, \quad \Var_{\text{obs}}(X_k) = 0 \quad (\text{exact}).
\end{equation}

\subsection{Global-Inverse-Free Measurable Predictor}
\label{sec:gaussian:predictor}

Let $C_\lambda \coloneqq \Sigma_{\lambda,II}$ denote the Gaussian bulk covariance operator on $\mathcal{H}_I$. Because $C_\lambda$ is trace-class and compact, it lacks a bounded global inverse. We define the bulk Cameron--Martin space $H_{C_\lambda} \coloneqq \operatorname{Range}(C_\lambda^{1/2})$ and the corresponding covariance norm.

The leakage cross-covariance vector satisfies $\Sigma_{\lambda,Ik} \in H_{C_\lambda}$ with norm $\|\Sigma_{\lambda,Ik}\|_{H_{C_\lambda}} = O(\lambda^{-3})$. This range inclusion allows us to define the conditional predictor $\mu_{k|I}(X_I) \coloneqq \E[X_k \mid X_I]$ as a measurable Wiener functional:
\begin{equation}
  \mu_{k|I}(X_I) = m_{\lambda,k} + W_{\Sigma_{\lambda,Ik}}(X_I - m_{\lambda,I}),
\end{equation}
which is the unique LMMSE predictor in $L^2$. The estimates in \Cref{app:gaussian_lift} imply that, for all sufficiently large $\lambda$, the bulk marginals $\mu_{\text{obs},I}$ and $\mu_{\lambda,I}$ are equivalent and have finite bulk relative entropy:
\begin{equation}
  \KL(\mu_{\text{obs},I} \parallel \mu_{\lambda,I}) = O(\lambda^{-4}) \to 0 \quad \text{as } \lambda \to \infty.
\end{equation}

\subsection{Gaussian Renormalized Finite-Part Functional}
\label{sec:gaussian:functional}

Although the joint KL divergence diverges under the exact conditioning, we can define a renormalized finite-part action.

\begin{definition}[Gaussian Renormalized Finite-Part Functional]
  \label{def:unified_action}
  Whenever $S_{k,\lambda}>0$ and $\mu_{\text{obs},I}$ is equivalent to $\mu_{\lambda,I}$, we define the Gaussian renormalized finite-part functional of the conditioned law:
  \begin{equation}
    \label{eq:unified_action}
    \begin{aligned}
      \mathfrak{J}_\lambda(\mu_{\mathrm{obs}}; \mu_\lambda) \coloneqq &\KL(\mu_{\mathrm{obs},I} \parallel \mu_{\lambda,I}) \\
      &+ \frac{1}{2} S_{k,\lambda}^{-1} \E_{\mu_{\mathrm{obs}}}\left[ \left( X_k - \mu_{k|I}(X_I) \right)^2 \right].
    \end{aligned}
  \end{equation}
  This represents the finite part of the joint KL divergence under the conditional normalization scheme, where the Dirac and Gaussian normalization singularities are subtracted.
\end{definition}

\paragraph{Weak scheme-independence.}
For an evidence level $t\in\mathbb{R}$, write $\mu_\lambda^t$ for the weakly continuous conditional law from \Cref{app:gaussian_lift} and set
\[
  \mathfrak{J}_\lambda(t)\coloneqq
  \mathfrak{J}_\lambda(\mu_\lambda^t;\mu_\lambda).
\]
Then, for any two evidence levels $t_1,t_2$ for which the finite part is defined,
\begin{equation}
  \label{eq:weak_scheme_independence}
  \mathfrak{J}_\lambda(t_2)-\mathfrak{J}_\lambda(t_1)
  =
  \frac{(t_2-m_{\lambda,k})^2-(t_1-m_{\lambda,k})^2}
       {2\Sigma_{\lambda,kk}}.
\end{equation}
In the calibrated scaling $\Sigma_{\lambda,kk}=\varepsilon_0/(2\lambda^2)$, if
$b_i=\lambda(m_{\lambda,k}-t_i)$, this becomes
\[
  \mathfrak{J}_\lambda(t_2)-\mathfrak{J}_\lambda(t_1)
  =
  \frac{b_2^2-b_1^2}{\varepsilon_0}.
\]
\begin{proof}
Let $\Delta_t\coloneqq t-m_{\lambda,k}$ and
\[
  \rho_\lambda
  \coloneqq
  \frac{\|\Sigma_{\lambda,Ik}\|_{H_{C_\lambda}}^2}{\Sigma_{\lambda,kk}}
  =
  1-\frac{S_{k,\lambda}}{\Sigma_{\lambda,kk}}.
\]
The Gaussian entropy computation in \Cref{app:gaussian_lift} gives
\[
  \KL(\mu_{\lambda,I}^t\parallel\mu_{\lambda,I})
  =
  \frac{1}{2}\frac{\Delta_t^2}{\Sigma_{\lambda,kk}}\rho_\lambda
  -\frac{1}{2}\log(1-\rho_\lambda)-\frac{1}{2}\rho_\lambda .
\]
The conditional-residual term in \eqref{eq:unified_action} is
\[
  \frac{1}{2S_{k,\lambda}}
  \E_{\mu_\lambda^t}\!\left[
    (X_k-\mu_{k|I}(X_I))^2
  \right]
  =
  \frac{1}{2}\rho_\lambda
  +
  \frac{1}{2}\frac{\Delta_t^2}{\Sigma_{\lambda,kk}}(1-\rho_\lambda).
\]
Adding these two identities cancels the $\rho_\lambda/2$ terms and yields
\[
  \mathfrak{J}_\lambda(t)
  =
  \frac{(t-m_{\lambda,k})^2}{2\Sigma_{\lambda,kk}}
  -\frac{1}{2}\log(1-\rho_\lambda).
\]
The logarithmic term depends on the reference covariance and the conditioning
scale, but not on the evidence value $t$. Hence it cancels in the difference
between $t_2$ and $t_1$, proving \eqref{eq:weak_scheme_independence}. The same
argument shows weak independence from the mollifying kernel: replacing the
Gaussian Dirac mollifier in \Cref{prop:mollification_limit} by any centered
approximate identity with finite second moment changes the extracted finite
part only by an additive term depending on the kernel, $\lambda$, and
$S_{k,\lambda}$, not on $t$. Therefore all evidence differences
$\Delta\mathfrak{J}_\lambda$ are kernel-independent.
\end{proof}

\subsection{Gaussian Main Theorem}
\label{sec:gaussian:theorem}

We now establish the main regularization and cancellation theorems for the finite-part functional.

\begin{theorem}[Gaussian Finite-Part Regularization]
  \label{thm:finite_part_regularization}
  Under Assumption~\ref{assum:gaussian_clamp}, there exists $\lambda_0>0$ such that, for $\lambda\ge\lambda_0$, the renormalized finite-part functional under the conditioned law $\mu_{\mathrm{obs}}$ is well-defined and satisfies:
  \begin{equation}
    \label{eq:unified_regularization}
    \mathfrak{J}_\lambda(\mu_{\mathrm{obs}}; \mu_\lambda) = \frac{b_k^2}{\varepsilon_0} + O(\lambda^{-4}).
  \end{equation}
  In particular, the functional remains bounded at $O(1)$ and converges to:
  \begin{equation}
    \lim_{\lambda\to\infty} \mathfrak{J}_\lambda(\mu_{\mathrm{obs}}; \mu_\lambda) = \frac{b_k^2}{\varepsilon_0}.
  \end{equation}
\end{theorem}
\begin{proof}
  See \Cref{app:gaussian_lift} for the full derivation.
\end{proof}

\begin{corollary}[Exact Decoupled-Centered Cancellation]
  \label{cor:exact_pole_zero}
  Under Assumption~\ref{assum:gaussian_clamp}, if the target coordinate is dynamically decoupled from the bulk ($a_{Ik} = 0$) and the evidence is centered at the stationary mean ($x_k^* = m_{\lambda,k}$), then:
  \begin{equation}
    \mathfrak{J}_\lambda(\mu_{\mathrm{obs}}; \mu_\lambda) = 0 \quad \text{for every } \lambda > 0.
  \end{equation}
\end{corollary}
\begin{proof}
  The decoupling implies $\Sigma_{\lambda,Ik} = 0$, so $\KL(\mu_{\text{obs},I} \parallel \mu_{\lambda,I}) = 0$. The centering makes the conditional residual zero under the conditioning, so both terms in \eqref{eq:unified_action} vanish.
\end{proof}

% ============================================================
%  §5  ELLIPTIC GAUSSIAN FIELD EXAMPLE
% ============================================================
\section{Elliptic Gaussian Field Example}
\label{sec:spde}

We now record a standard elliptic Gaussian field example to show that the abstract Hilbert-space assumptions are nonempty and natural for covariance operators arising from stochastic evolution equations and elliptic Gaussian priors \cite{daprato2014stochastic,lord2014introduction,stuart2010inverse}.

\subsection{Elliptic Gaussian field on $[0,1]$}

Let
\[
  L=-\frac{d^2}{dx^2}+\kappa^2
\]
on $L^2(0,1)$ with Dirichlet boundary conditions and $\kappa>0$. Then $\mathcal{A}_0=-L$ with domain $H^2(0,1)\cap H_0^1(0,1)$ is dissipative, $L$ is positive self-adjoint and sectorial, and $e^{-tL}$ is an exponentially stable analytic contraction semigroup. The inverse covariance
\[
  Q_0=L^{-1}
\]
is compact, positive, self-adjoint, and trace class. With eigenfunctions $e_j(x)=\sqrt{2}\sin(\pi jx)$, its eigenvalues are
\[
  q_j=(\pi^2j^2+\kappa^2)^{-1}.
\]
Thus $Q_0$ is a typical infinite-dimensional covariance operator: it is compact and trace class, has no bounded inverse on the ambient space, and carries its inverse only as a Cameron--Martin quadratic form \cite{bogachev1998gaussian,simon2005trace}.

\subsection{Cameron--Martin range and covariance norm}

The Cameron--Martin space associated with $Q_0$ is
\[
  H_{Q_0}=\operatorname{Range}(Q_0^{1/2})=H_0^1(0,1),
\]
with equivalent covariance norm
\[
  \|h\|_{Q_0}^2
  =\langle h,Lh\rangle_{L^2}
  =\int_0^1 \bigl(|h'(x)|^2+\kappa^2|h(x)|^2\bigr)\,dx.
\]
This identifies the variational Schur projection in a Sobolev quadratic form, as in the standard covariance-norm representation of Gaussian Hilbert priors \cite{bogachev1998gaussian,stuart2010inverse}. The covariance inverse does not act as a bounded operator on $L^2(0,1)$; it acts only through this quadratic form on its Cameron--Martin domain.

\subsection{Spectral sensor conditioning}

Fix a target eigenmode $e_k$ and define the observed coordinate
\[
  X_k=\langle X,e_k\rangle_{L^2}.
\]
The unperturbed target variance is
\[
  \Sigma_{0,kk}=q_k=(\pi^2k^2+\kappa^2)^{-1}.
\]
A regularized restoring drift on this coordinate produces the exact normalization
\[
  \Sigma_{\lambda,kk}=\frac{\varepsilon_0}{2\lambda^2}-R_\lambda,
  \qquad R_\lambda=O(\lambda^{-6}),
\]
under the block assumptions of \Cref{sec:infdim}. Off-diagonal bounded perturbations of the drift operator generate the cross-covariance vector entering the variational Schur residual. The resulting subtraction is evaluated by the Cameron--Martin quadratic form above, not by a global inverse on $L^2(0,1)$.

This example shows that the abstract theorem applies to an infinite-dimensional Gaussian field: the residual is controlled by a Sobolev norm on $H_0^1(0,1)$ while the ambient covariance remains compact and non-invertible.

\section{Discussion}
\label{sec:discussion}

We have shown that singular Gaussian coordinate conditioning in a separable Hilbert space admits an inverse-free Schur normal form.  The main obstruction is the absence of a bounded inverse for compact covariance blocks.  The variational Cameron--Martin projection is the target-coordinate shorted-operator subtraction evaluated in the covariance-energy scale, and it yields the residual estimate needed for the hard-conditioning limit.  The scalar nature of the target coordinate is essential to the present proof.  Three extensions mark the boundary between the theorem proved here and a genuinely finite-rank conditioning theory. We formalize these extensions as exact theorems below.

\subsection{Matrix algebraic extension}
\label{sec:discussion:matrix-extension}

For an $m$-dimensional observed subspace $E=\operatorname{span}\{\phi_1,\ldots,\phi_m\}$, the scalar target block $\Sigma_{\lambda,kk}$ is replaced by the covariance matrix $\Sigma_{\lambda,EE}\in\mathbb{R}^{m\times m}$. The scalar calibration becomes a matrix Lyapunov balancing problem.

\begin{theorem}[Matrix Lyapunov Calibration]
  \label{thm:matrix_lyapunov}
  Let $E$ be an $m$-dimensional target subspace. Suppose the regularized drift and diffusion on $E$ are given by $A_{\lambda,E} = -\lambda \Lambda_E$ and $D_{\lambda,E} = \frac{1}{\lambda} \Delta_E$, where $\Lambda_E, \Delta_E \in \mathbb{R}^{m \times m}$ are strictly positive definite matrices. Let $\Sigma_{\lambda,EE}$ be the unique positive definite solution to the matrix Lyapunov equation:
  \begin{equation}
    \label{eq:discussion_matrix_lyapunov}
    A_{\lambda,E}\Sigma_{\lambda,EE}
    +\Sigma_{\lambda,EE}A_{\lambda,E}^{*}
    +D_{\lambda,E}=0.
  \end{equation}
  Then $\Sigma_{\lambda,EE}$ scales exactly as $\lambda^{-2}$. Specifically, the normalized matrix $\widetilde{\Sigma}_{EE} \coloneqq \lambda^2 \Sigma_{\lambda,EE}$ is independent of $\lambda$ and is the unique positive definite solution to the scale-free Lyapunov equation:
  \begin{equation}
    \label{eq:scale_free_lyapunov}
    \Lambda_E \widetilde{\Sigma}_{EE} + \widetilde{\Sigma}_{EE} \Lambda_E^* = \Delta_E.
  \end{equation}
\end{theorem}
\begin{proof}
  Substituting $A_{\lambda,E} = -\lambda \Lambda_E$ and $D_{\lambda,E} = \frac{1}{\lambda} \Delta_E$ into \eqref{eq:discussion_matrix_lyapunov} yields:
  \[
    -\lambda \Lambda_E \Sigma_{\lambda,EE} - \lambda \Sigma_{\lambda,EE} \Lambda_E^* + \frac{1}{\lambda} \Delta_E = 0.
  \]
  Multiplying the entire equation by $\lambda$ gives:
  \[
    \Lambda_E (\lambda^2 \Sigma_{\lambda,EE}) + (\lambda^2 \Sigma_{\lambda,EE}) \Lambda_E^* = \Delta_E.
  \]
  Defining $\widetilde{\Sigma}_{EE} \coloneqq \lambda^2 \Sigma_{\lambda,EE}$, we recover \eqref{eq:scale_free_lyapunov}. Since $\Lambda_E$ is positive definite, its spectrum lies in the open right half-plane, making $-\Lambda_E$ Hurwitz. Because $\Delta_E$ is positive definite, standard continuous Lyapunov theory guarantees that \eqref{eq:scale_free_lyapunov} admits a unique positive definite solution $\widetilde{\Sigma}_{EE}$. Therefore, the exact scaling $\Sigma_{\lambda,EE} = \lambda^{-2} \widetilde{\Sigma}_{EE}$ holds for all $\lambda > 0$, properly generalizing the scalar calibration condition $\Sigma_{\lambda,kk} = \varepsilon_0 / (2\lambda^2)$.
\end{proof}

\subsection{Operator-valued residuals}
\label{sec:discussion:operator-valued-residuals}

For an $m$-dimensional target space $E$, the Schur residual must be a positive-semidefinite operator on $E$. Because the free-block covariance $\Sigma_{\lambda,II}$ is compact and lacks a bounded global inverse, the formal finite-dimensional matrix expression $R_{\lambda,E} = \Sigma_{\lambda,EI}\Sigma_{\lambda,II}^{-1}\Sigma_{\lambda,IE}$ must be rigorously replaced by the shorted-operator construction evaluated in the Cameron--Martin energy norm.

\begin{theorem}[Operator-Valued Schur Residual]
  \label{thm:operator_residual}
  Let $\mathcal{H} = E \oplus \mathcal{H}_I$, with $E$ finite-dimensional. Let $\Sigma_\lambda$ be the non-negative trace-class covariance operator on $\mathcal{H}$ with blocks $\Sigma_{\lambda,EE}$, $\Sigma_{\lambda,EI}$, $\Sigma_{\lambda,IE}$, and $\Sigma_{\lambda,II}$. Define the shorted-operator residual $R_{\lambda,E}$ on $E$ by its quadratic forms:
  \begin{equation}
    \label{eq:discussion_operator_residual_form}
    \langle v,R_{\lambda,E}v\rangle_E
    =
    \sup_{\langle u,\Sigma_{\lambda,II}u\rangle>0}
    \frac{|\langle u,\Sigma_{\lambda,IE}v\rangle|^2}
         {\langle u,\Sigma_{\lambda,II}u\rangle},
    \qquad v\in E.
  \end{equation}
  Assume the structural range inclusion $\operatorname{Range}(\Sigma_{\lambda,IE}) \subset \operatorname{Range}(\Sigma_{\lambda,II}^{1/2}) \eqqcolon H_{C_\lambda}$. Then $R_{\lambda,E}$ is a well-defined bounded positive-semidefinite operator on $E$, uniquely characterized by:
  \begin{equation}
    R_{\lambda,E} = (\Sigma_{\lambda,II}^{-1/2}\Sigma_{\lambda,IE})^* (\Sigma_{\lambda,II}^{-1/2}\Sigma_{\lambda,IE}).
  \end{equation}
  Furthermore, if the cross-covariance satisfies the energy bound $\|\Sigma_{\lambda,IE} v\|_{H_{C_\lambda}} \le M \lambda^{-3} \|v\|_E$ for all $v \in E$ and some constant $M > 0$, then the operator norm of the residual satisfies $\|R_{\lambda,E}\|_{\mathcal{L}(E)} = O(\lambda^{-6})$.
\end{theorem}
\begin{proof}
  By the assumed range inclusion, for any $v \in E$, the vector $y_v \coloneqq \Sigma_{\lambda,IE} v$ lies in $H_{C_\lambda} = \operatorname{Range}(\Sigma_{\lambda,II}^{1/2})$. Thus, there exists a unique element $w_v \in \overline{\operatorname{Range}(\Sigma_{\lambda,II}^{1/2})}$ such that $\Sigma_{\lambda,II}^{1/2} w_v = y_v$. We denote $w_v \coloneqq \Sigma_{\lambda,II}^{-1/2} \Sigma_{\lambda,IE} v$.
  
  For any $u \in \mathcal{H}_I$ with $\langle u, \Sigma_{\lambda,II} u \rangle > 0$, the ratio in \eqref{eq:discussion_operator_residual_form} can be rewritten as:
  \[
    \frac{|\langle u, \Sigma_{\lambda,II}^{1/2} w_v \rangle|^2}{\|\Sigma_{\lambda,II}^{1/2} u\|^2} = \frac{|\langle \Sigma_{\lambda,II}^{1/2} u, w_v \rangle|^2}{\|\Sigma_{\lambda,II}^{1/2} u\|^2}.
  \]
  By the Cauchy--Schwarz inequality, this ratio is bounded above by $\|w_v\|^2$. The supremum is achieved by taking a sequence of vectors $\Sigma_{\lambda,II}^{1/2} u$ converging to $w_v$, yielding:
  \[
    \langle v, R_{\lambda,E} v \rangle_E = \|w_v\|^2 = \|\Sigma_{\lambda,II}^{-1/2} \Sigma_{\lambda,IE} v\|^2 = \langle v, (\Sigma_{\lambda,II}^{-1/2}\Sigma_{\lambda,IE})^* (\Sigma_{\lambda,II}^{-1/2}\Sigma_{\lambda,IE}) v \rangle_E.
  \]
  Since this holds for all $v \in E$, it uniquely defines $R_{\lambda,E}$ as a positive-semidefinite matrix on $E$. Finally, the assumption $\|\Sigma_{\lambda,IE} v\|_{H_{C_\lambda}} \le M \lambda^{-3} \|v\|_E$ means exactly that $\|w_v\| \le M \lambda^{-3} \|v\|_E$. Therefore, $\langle v, R_{\lambda,E} v \rangle_E \le M^2 \lambda^{-6} \|v\|_E^2$, which implies the operator norm bound $\|R_{\lambda,E}\|_{\mathcal{L}(E)} \le M^2 \lambda^{-6} = O(\lambda^{-6})$.
\end{proof}

\subsection{General observation operators and misalignment}
\label{sec:discussion:misalignment}

The scalar proof in this paper uses a block decomposition aligned with a single observed coordinate. However, typical finite-rank observation operators $L:\mathcal H\to\mathbb R^m$ select a subspace $E$ that does not commute with the free prior covariance, meaning the target-free split is not an eigenbasis split. This geometric misalignment generates cross-coupling that requires a precise resolvent analysis.

\begin{theorem}[Misaligned Covariance Coupling]
  \label{thm:misaligned_coupling}
  Let $P_E$ be an orthogonal projection onto an $m$-dimensional observed subspace $E$, and let $\mathcal{H}_I = E^\perp$. Let the regularized dynamics be decomposed such that the target block matches \Cref{thm:matrix_lyapunov} with $A_{\lambda,E} = -\lambda \Lambda_E$, the cross-noise is zero ($D_{\lambda,IE} = 0$), and the off-diagonal drift coupling is denoted by $A_{\lambda,IE}$. 
  If $A_{\lambda,IE}$ maps $E$ into the Cameron--Martin space $H_{Q_0}$ of the free background covariance $Q_0$ on $\mathcal{H}_I$, and the restricted drift $A_{II}$ generates an exponentially stable semigroup on $H_{Q_0}$, then the cross-covariance satisfies the exact identity:
  \begin{equation}
    \label{eq:misaligned_resolvent_identity}
    \Sigma_{\lambda,IE} = \int_0^\infty e^{t A_{II}} A_{\lambda,IE} \Sigma_{\lambda,EE} e^{t A_{\lambda,E}^*} dt.
  \end{equation}
  Furthermore, the misaligned cross-covariance energy scales as $\|\Sigma_{\lambda,IE} v\|_{H_{Q_0}} = O(\lambda^{-3})$ for any $v \in E$.
\end{theorem}
\begin{proof}
  The off-diagonal block of the global Lyapunov equation $\mathcal{A}^{(\lambda)} \Sigma_\lambda + \Sigma_\lambda (\mathcal{A}^{(\lambda)})^* + \mathcal{D}^{(\lambda)} = 0$ yields the Sylvester equation:
  \[
    A_{II} \Sigma_{\lambda,IE} + \Sigma_{\lambda,IE} A_{\lambda,E}^* + A_{\lambda,IE} \Sigma_{\lambda,EE} = 0.
  \]
  Since $A_{II}$ generates an exponentially stable semigroup $e^{t A_{II}}$ and $A_{\lambda,E} = -\lambda \Lambda_E$ is Hurwitz, the unique solution is given by the integral representation \eqref{eq:misaligned_resolvent_identity}.
  
  To bound the Cameron--Martin norm, fix $v \in E$. We have $e^{t A_{\lambda,E}^*} v = e^{-\lambda t \Lambda_E^*} v$. By \Cref{thm:matrix_lyapunov}, $\Sigma_{\lambda,EE} = O(\lambda^{-2})$. Since $A_{\lambda,IE}$ maps into $H_{Q_0}$ and $E$ is finite-dimensional, there exists a constant $C$ such that $\|A_{\lambda,IE} w\|_{H_{Q_0}} \le C \|w\|_E$ for all $w \in E$. Using the uniform exponential stability of $e^{t A_{II}}$ on $H_{Q_0}$, namely $\|e^{t A_{II}}\|_{\mathcal{L}(H_{Q_0})} \le M e^{-\omega t}$ for some $M \ge 1, \omega > 0$, we estimate:
  \begin{align*}
    \|\Sigma_{\lambda,IE} v\|_{H_{Q_0}} 
    &\le \int_0^\infty \|e^{t A_{II}}\|_{\mathcal{L}(H_{Q_0})} \|A_{\lambda,IE} \Sigma_{\lambda,EE} e^{-\lambda t \Lambda_E^*} v\|_{H_{Q_0}} dt \\
    &\le \int_0^\infty M e^{-\omega t} C \|\Sigma_{\lambda,EE}\|_{\mathcal{L}(E)} \|e^{-\lambda t \Lambda_E^*}\|_{\mathcal{L}(E)} \|v\|_E dt.
  \end{align*}
  Let $\gamma > 0$ be the smallest eigenvalue of the positive definite matrix $\Lambda_E^*$, so $\|e^{-\lambda t \Lambda_E^*}\|_{\mathcal{L}(E)} \le K e^{-\lambda \gamma t}$. Then:
  \begin{align*}
    \|\Sigma_{\lambda,IE} v\|_{H_{Q_0}} 
    &\le M C K \|\Sigma_{\lambda,EE}\|_{\mathcal{L}(E)} \|v\|_E \int_0^\infty e^{-(\omega + \lambda \gamma) t} dt \\
    &= M C K \|\Sigma_{\lambda,EE}\|_{\mathcal{L}(E)} \|v\|_E \frac{1}{\omega + \lambda \gamma}.
  \end{align*}
  Since $\|\Sigma_{\lambda,EE}\|_{\mathcal{L}(E)} = O(\lambda^{-2})$ and $(\omega + \lambda \gamma)^{-1} = O(\lambda^{-1})$, we conclude that $\|\Sigma_{\lambda,IE} v\|_{H_{Q_0}} = O(\lambda^{-3})$. Combined with \Cref{thm:operator_residual}, this verifies that the shorted-operator residual maintains the $O(\lambda^{-6})$ scaling even in the presence of observation misalignment.
\end{proof}

The result is therefore deliberately limited to covariance and Gaussian-conditioning questions in the rank-one aligned setting.  It does not assert a general non-Gaussian conditioning theory.  For non-Gaussian stationary processes, the projection residual $S_{k,\lambda}=\varepsilon_0/(2\lambda^2)-R_\lambda$ still has a second-moment interpretation, but Radon disintegration, finite-part functionals, and measure-class behavior require separate hypotheses \cite{bogachev2007measure,kallenberg2021foundations}.  Dynamical relaxation after conditioning, or transport distances between post-conditioning laws, belongs to a different problem.  The present contribution is the operator-theoretic construction of the singular Gaussian conditioning residual itself.

% ============================================================
%  \S6  CONCLUSION
% ============================================================
\section{Conclusion}
\label{sec:conclusion}

We established a Hilbert-space framework for singular coordinate conditioning of Gaussian measures. The central object is an inverse-free variational Schur projection on the Cameron--Martin range, equivalently the one-dimensional target-coordinate form of the shorted-operator subtraction. It gives canonical normalization, an $O(\lambda^{-6})$ residual estimate, Galerkin stability, and a Radon Gaussian conditioning statement. The elliptic field example shows that the assumptions are natural for trace-class covariance operators arising from stochastic evolution equations and elliptic Gaussian priors \cite{daprato2014stochastic,lord2014introduction,stuart2010inverse}.

\section*{Declaration of Generative AI and AI-Assisted Technologies in the Manuscript Preparation Process}

During the preparation of this work, the author(s) used Google DeepMind's Antigravity assistant in order to merge the LaTeX manuscripts, clean up cross-referencing tags, format mathematical notation, and edit the draft for style consistency. After using this tool/service, the author(s) reviewed and edited the content as needed and take(s) full responsibility for the content of the published article.

% ============================================================
%  REFERENCES
% ============================================================
\bibliographystyle{amsplain}

\documentclass[11pt]{amsart}

\usepackage{amsmath}
\usepackage{amssymb}
\usepackage{amsthm}
\usepackage{mathtools}
\usepackage{enumitem}
%\usepackage{thmtools}
\usepackage{hyperref}
\hypersetup{hypertexnames=false}
\usepackage[nameinlink,noabbrev]{cleveref}
\usepackage{booktabs}
\usepackage{graphicx}
\usepackage{xcolor}


% ---- Theorem environments ----
\theoremstyle{plain}
\newtheorem{theorem}{Theorem}
\newtheorem{proposition}[theorem]{Proposition}
\newtheorem{lemma}[theorem]{Lemma}
\newtheorem{corollary}[theorem]{Corollary}
\newtheorem{claim}[theorem]{Claim}

\theoremstyle{definition}
\newtheorem{definition}[theorem]{Definition}
\newtheorem{assumption}[theorem]{Assumption}
\newtheorem{example}[theorem]{Example}

\theoremstyle{remark}
\newtheorem{remark}[theorem]{Remark}

% ---- Macros ----
\newcommand{\cond}{\operatorname{cond}}
\newcommand{\KL}{D_{\operatorname{KL}}}
\newcommand{\E}{\mathbb{E}}
\newcommand{\Var}{\operatorname{Var}}
\newcommand{\Cov}{\operatorname{Cov}}
\newcommand{\tr}{\operatorname{Tr}}
\newcommand{\diag}{\operatorname{diag}}

\newcommand{\Range}{\operatorname{Range}}
\newcommand{\Span}{\operatorname{span}}
\newcommand{\id}{\mathrm{id}}
% \newcommand{\TODO}[1]{\textcolor{red}{[TODO: #1]}} % removed for submission

% ---- Cleveref names ----
\crefname{theorem}{theorem}{theorems}
\Crefname{theorem}{Theorem}{Theorems}
\crefname{proposition}{proposition}{propositions}
\Crefname{proposition}{Proposition}{Propositions}
\crefname{lemma}{lemma}{lemmas}
\Crefname{lemma}{Lemma}{Lemmas}
\crefname{corollary}{corollary}{corollaries}
\Crefname{corollary}{Corollary}{Corollaries}
\crefname{claim}{claim}{claims}
\Crefname{claim}{Claim}{Claims}
\crefname{definition}{definition}{definitions}
\Crefname{definition}{Definition}{Definitions}
\crefname{assumption}{assumption}{assumptions}
\Crefname{assumption}{Assumption}{Assumptions}
\crefname{example}{example}{examples}
\Crefname{example}{Example}{Examples}
\crefname{remark}{remark}{remarks}
\Crefname{remark}{Remark}{Remarks}
\crefname{equation}{Eq.}{Eqs.}
\Crefname{equation}{Equation}{Equations}
\crefname{figure}{fig.}{figs.}
\Crefname{figure}{Fig.}{Figs.}
\crefname{table}{table}{tables}
\Crefname{table}{Table}{Tables}
\crefname{section}{Sec.}{Secs.}
\Crefname{section}{Section}{Sections}

\title[Singular Gaussian Conditioning in Hilbert Space]{Singular Gaussian Conditioning in Hilbert Space: Inverse-Free Schur Projection and Variance Collapse}

\author{Lucia Shang}
\thanks{Manuscript received May 2026. Repository information omitted for submission.}

\begin{document}

% ============================================================
%  ABSTRACT
% ============================================================
\begin{abstract}
We study singular coordinate conditioning for Gaussian measures on separable Hilbert spaces. Since trace-class covariance blocks are compact, the usual Schur-complement formula cannot be implemented through a bounded inverse. We introduce an inverse-free variational Schur projection on the Cameron--Martin range and prove canonical normalization together with an $O(\lambda^{-6})$ residual estimate. The construction is stable under Galerkin truncation for exponentially stable analytic semigroup dynamics and applies to sectorial elliptic Gaussian fields.
\end{abstract}

\keywords{Singular conditioning; Gaussian measures; Cameron--Martin space; Schur complement; Hilbert-space covariance operators; trace-class covariance; Radon Gaussian conditioning; Galerkin approximation; stochastic evolution equations}

\maketitle

% ============================================================
%  \S1  INTRODUCTION
% ============================================================
\section{Introduction}
\label{sec:intro}

Singular coordinate conditioning arises when an exact constraint $X_k=x_k^*$ is reached as a zero-noise or infinite-precision limit of Gaussian conditioning problems.  Finite-dimensional Gaussian conditioning is governed by the ordinary Schur complement of a covariance matrix \cite{horn2012matrix,rasmussen2006gaussian}.  Infinite-dimensional Gaussian conditioning is subtler: Radon Gaussian measures on separable Hilbert or Banach spaces have compact covariance operators, and their Cameron--Martin spaces are covariance-norm domains rather than ambient statistical parameter spaces \cite{bogachev1998gaussian,daprato2014stochastic,steinwart2024conditioning}.  In Hilbert-space Gaussian conditioning, shorted operators give the covariance of conditioning on closed subspaces, while recent inverse-problem work emphasizes that positive-noise posterior formulas must be interpreted through approximation, disintegration, or Hilbert-space posterior equivalence rather than through a global covariance inverse \cite{owhadi2015conditioning,stuart2010inverse,chen2024gaussian,carerelie2024optimal,carerelie2025posterior}.

The obstruction studied here is the covariance-block obstruction behind that phenomenon.  In a finite-dimensional model, the conditional variance of a coordinate can be written as $A-BD^{-1}B^*$ after a block decomposition.  In a separable Hilbert space the corresponding block $D$ is typically compact and trace class, hence $D^{-1}$ is not a bounded operator on the ambient space \cite{bogachev1998gaussian,simon2005trace,daprato2014stochastic}.  Classical Schur-complement and shorted-operator theory, originating in Krein's theory of non-negative extensions and developed systematically by Anderson and Trapp, provides variational tools for positive block operators \cite{krein1947selfadjoint,anderson1975shorted,douglas1966majorization}.  However, this theory does not by itself give the scaling law for a singular Gaussian conditioning residual \cite{tretter2008spectral,contino2018schur,friedrich2017generalized,hassi2023complementation}.  This paper formulates the residual directly as a Cameron--Martin norm projection and never invokes a pseudoinverse or an ambient precision operator.

The hard-conditioning limit also reaches the Feldman--H\'{a}jek measure-class boundary.  With positive observational noise, Gaussian Bayesian inverse problems may remain in the equivalent-measure regime \cite{stuart2010inverse,carerelie2024optimal,carerelie2025posterior}.  Under an exact coordinate constraint, the limiting conditional law is supported on a hyperplane of prior measure zero, so the equivalent-measure regime is left \cite{feldman1958equivalence,hajek1958property,bogachev1998gaussian}.  A useful theory must therefore separate the diverging joint divergence caused by exact evidence from the finite covariance residual that remains visible in the block structure.

We prove that this finite residual is controlled by an inverse-free variational Schur projection.  Under the Lyapunov block assumptions stated below, the target residual has the exact form
\[
  S_{k,\lambda}=\frac{\varepsilon_0}{2\lambda^2}-R_\lambda,
  \qquad
  R_\lambda=O(\lambda^{-6}).
\]
Consequently,
\[
  S_{k,\lambda}^{-1}=\frac{2}{\varepsilon_0}\lambda^2+O(\lambda^{-2}).
\]
The proof uses covariance-norm forms, shorting, Douglas range inclusion, Lyapunov block identities, and trace-ideal convergence. These are standard inputs in operator and SPDE analysis \cite{krein1947selfadjoint,anderson1975shorted,douglas1966majorization} \cite{simon2005trace,daprato2014stochastic}. The Cameron--Martin space is used only as a Hilbert subspace equipped with its covariance norm; no differentiable statistical or geometric structure is assumed \cite{bogachev1998gaussian}.
The operator base is an exponentially stable analytic semigroup with bounded drift perturbations, so dissipative sectorial generators such as $-(-\Delta+\kappa^2)$ fall inside the abstract setup.

The paper has four main contributions:
\begin{enumerate}
  \item \textbf{Canonical normalization.} We identify the block-Lyapunov mechanism that fixes the observed coordinate variance at order $\lambda^{-2}$ with explicit coefficient $\varepsilon_0/2$.
  \item \textbf{Inverse-free Schur projection.} We identify the residual subtraction with the target-coordinate quadratic form of the shorted-operator variational formula, then prove the $O(\lambda^{-6})$ residual estimate in the same operator-theoretic regime where compact covariance blocks preclude a bounded Schur inverse.
  \item \textbf{Analytic-semigroup and Galerkin stability.} We work with an analytic semigroup drift base plus bounded perturbations, and prove that finite-dimensional truncations converge to the Hilbert-space residual in trace ideal norm, matching the discretization-independent perspective used in Hilbert-space Gaussian inverse problems \cite{stuart2010inverse,carerelie2024optimal,carerelie2025posterior}.
  \item \textbf{Radon Gaussian conditioning.} For Radon Gaussian laws, we construct the conditional predictor as a measurable Cameron--Martin projection and show that the associated renormalized conditioning functional remains bounded, using regular conditional probabilities and continuous Gaussian disintegration \cite{bogachev1998gaussian,lagatta2013continuous,steinwart2024conditioning}.
\end{enumerate}

\subsection{Positioning in the literature}
\label{sec:related}

The paper is closest to four bodies of work.  First, it uses the classical measure-class structure of Gaussian measures, especially the Cameron--Martin theorem and the Feldman--H\'{a}jek dichotomy, as presented in standard references on Gaussian measures \cite{feldman1958equivalence,hajek1958property,bogachev1998gaussian}.  Second, it is adjacent to Bayesian inverse problems and Hilbert-space Gaussian conditioning, including the shorted-operator description of Gaussian conditioning on closed subspaces, positive-noise posterior formulas, low-rank posterior approximations, and nonlinear observation maps \cite{owhadi2015conditioning,stuart2010inverse,chen2024gaussian,steinwart2024conditioning,carerelie2024optimal,carerelie2025posterior}.  Third, the variational residual is informed by Schur complements, shorted operators, Douglas range inclusion, and modern complementability theory for block operators \cite{krein1947selfadjoint,anderson1975shorted,douglas1966majorization,tretter2008spectral,contino2018schur,friedrich2017generalized,hassi2023complementation}.  Fourth, the motivating examples come from stochastic evolution equations and elliptic Gaussian fields, where dissipative sectorial generators produce compact trace-class stationary covariances that are naturally approximated by Galerkin truncation \cite{daprato2014stochastic,lord2014introduction,simon2005trace}.

\subsection{Organization of the paper}
\label{sec:organization}

\Cref{sec:setup} defines the operator model and establishes the canonical normalization. \Cref{sec:infdim} establishes the Hilbert-space variational Schur residual. \Cref{sec:gaussian_zero} gives the Radon Gaussian conditional extension and the renormalized conditioning functional. \Cref{sec:spde} gives an elliptic Gaussian-field example. \Cref{sec:discussion} records limitations and extensions, and \Cref{sec:conclusion} concludes. Proofs are given in \Cref{app:block_cumulant_analysis,app:gaussian_lift}; numerical diagnostics and Galerkin verification are given in \Cref{app:experiments}.

\section{Operator Model and Canonical Normalization}
\label{sec:setup}

\subsection{Operator Dynamics and Regularized Conditioning}
\label{sec:setup:dynamics}

We model the continuous-time dynamics as stochastic evolution on a separable Hilbert space $\mathcal{H}$ \cite{daprato2014stochastic}:
\begin{equation}
  \label{eq:stochastic_evolution}
  dX_t = \mathcal{A} X_t \, dt + \mathcal{D}^{1/2} dW_t
\end{equation}
where $\mathcal{A}$ is the drift operator and $\mathcal{D}$ is the trace-class diffusion operator. We decompose the Hilbert space as $\mathcal{H} = \mathcal{H}_k \oplus \mathcal{H}_I$, where $\mathcal{H}_k \coloneqq \mathbb{R}\phi_k$ is the one-dimensional target coordinate subspace, and $\mathcal{H}_I \coloneqq \mathcal{H}_k^\perp$ is the free bulk subspace. Let $\Pi_k \coloneqq \phi_k \otimes \phi_k$ and $\Pi_I \coloneqq I - \Pi_k$ be the corresponding orthogonal projections.

To approximate the exact point conditioning $X_k = x_k^*$, we set the incoming target-coordinate blocks to zero and apply a regularization parameter $\lambda > 0$.

\begin{definition}[Calibrated Approximating Dynamics Family]
  \label{def:soft_do}
  The regularized conditioning operator of strength $\lambda > 0$ at coordinate $X_k$ acts on the operator dynamics \eqref{eq:stochastic_evolution} by modifying both the drift and the diffusion:
  \begin{equation}
    \label{eq:soft_do_sde}
    \mathcal{A}^{(\lambda)} \coloneqq \mathcal{A}_{\cond} - \lambda \Pi_k, \qquad
    \mathcal{D}^{(\lambda)} \coloneqq \mathcal{D}_{\cond} + d_\lambda \Pi_k
  \end{equation}
  where $\mathcal{A}_{\cond}$ is the block drift operator with incoming target-coordinate blocks set to zero ($\Pi_k \mathcal{A}_{\cond} = 0$), $\mathcal{D}_{\cond}$ has target-coordinate cross-noise set to zero, and $d_\lambda > 0$ is the local diffusion variance of the target coordinate at regularization strength $\lambda$. In the SDE, this corresponds to:
  \begin{equation}
    \label{eq:soft_do_sde_general}
    dX_{k,t} = -\lambda (X_{k,t} - x_k^*) dt + \sqrt{d_\lambda}\, dW_{k,t}.
  \end{equation}
\end{definition}

\subsection{Canonical Normalization}
\label{sec:setup:calibration}

The block-decoupled structure under the regularized dynamics yields:
\begin{equation}
  \mathcal{A}^{(\lambda)}=
  \begin{pmatrix}-\lambda&0\\a_{Ik}&\mathcal{A}_{II}\end{pmatrix},\qquad
  \mathcal{D}^{(\lambda)}=
  \begin{pmatrix}d_\lambda&0\\0&D_{II}\end{pmatrix},
\end{equation}
where $a_{Ik} \coloneqq \mathcal{A}_{Ik} \phi_k \in \mathcal{H}_I$ represents the off-diagonal drift coupling.
Let $\Sigma_\lambda$ denote the stationary covariance operator under regularization strength $\lambda$, solving the Lyapunov equation:
\begin{equation}
  \mathcal{A}^{(\lambda)}\Sigma_\lambda + \Sigma_\lambda(\mathcal{A}^{(\lambda)})^* + \mathcal{D}^{(\lambda)} = 0.
\end{equation}
The target block $kk$ decouples completely from the free coordinates, yielding the exact scalar equation:
\begin{equation}
  -2\lambda\Sigma_{\lambda,kk} + d_\lambda = 0 \quad \Longrightarrow \quad \Sigma_{\lambda,kk} = \frac{d_\lambda}{2\lambda}.
\end{equation}

\begin{proposition}[Canonical Scaling]
  \label{prop:exact_calibration}
  To maintain a target noise scale $\varepsilon_0 > 0$ independent of the auxiliary parameter $\lambda$ under the singular limit, we define the canonical scaling condition:
  \begin{equation}
    \lambda^2 \Sigma_{\lambda,kk} \equiv \frac{\varepsilon_0}{2} \quad \text{for all } \lambda > 0,
  \end{equation}
  which uniquely requires:
  \begin{equation}
    d_\lambda = \frac{\varepsilon_0}{\lambda} \quad \text{for every } \lambda > 0.
  \end{equation}
  Substitution of this target diffusion scaling yields the calibrated approximating dynamics:
  \begin{equation}
    dX_{k,t} = -\lambda (X_{k,t} - x_k^*) dt + \sqrt{\varepsilon_0 / \lambda}\, dW_{k,t}.
  \end{equation}
\end{proposition}

All subsequent results are stated for this calibrated family.

\begin{definition}[Schur LMMSE Invariant]
  \label{def:schur_invariant}
  Let $S_{k,\lambda} \coloneqq \|X_k - P_{\mathcal{M}_I} X_k\|_{L^2}^2$ represent the linear minimum mean squared error (LMMSE) residual of the target coordinate $X_k$ on the closed linear span $\mathcal{M}_I$ of the free coordinates $X_I$. The \emph{Schur LMMSE invariant} $m_2^{\mathrm{Schur}}$ is defined as:
  \begin{equation}
    \label{eq:m2_schur_definition}
    m_2^{\mathrm{Schur}} \coloneqq \lim_{\lambda\to\infty} \lambda^2 S_{k,\lambda}.
  \end{equation}
\end{definition}

\subsection{Operator Setup and Assumptions}
\label{sec:setup:assumptions}

We state the operator-theoretic assumptions that make the limit passage independent of the Galerkin dimension.

\begin{definition}[Drift Operator with Analytic Base]
  \label{def:drift_operator}
  A drift operator on a separable Hilbert space $\mathcal{H}$ is an operator of the form $\mathcal{A} = \mathcal{A}_0 + \mathcal{K}$, where:
  \begin{enumerate}
    \item[(i)] $\mathcal{A}_0: \operatorname{dom}(\mathcal{A}_0) \subset \mathcal{H} \to \mathcal{H}$ is closed, densely defined, and dissipative, and it generates a strongly continuous analytic contraction semigroup $(e^{t\mathcal{A}_0})_{t \ge 0}$ on $\mathcal{H}$ satisfying $\|e^{t\mathcal{A}_0}\| \le e^{-\alpha t}$ for some spectral gap $\alpha > 0$; equivalently, the base operator $-\mathcal{A}_0$ is sectorial of angle $<\pi/2$ in this exponentially stable contraction setting.
    \item[(ii)] $\mathcal{K}: \mathcal{H} \to \mathcal{H}$ is a bounded linear operator.
    \item[(iii)] The full operator $\mathcal{A} = \mathcal{A}_0 + \mathcal{K}$, with $\operatorname{dom}(\mathcal{A})=\operatorname{dom}(\mathcal{A}_0)$, generates a strongly continuous analytic semigroup $(e^{t\mathcal{A}})_{t \ge 0}$ by the bounded perturbation theorem, and this semigroup remains exponentially stable: $\|e^{t\mathcal{A}}\| \le M e^{-\beta t}$ for some $M \ge 1$, $\beta > 0$.
  \end{enumerate}
  In finite dimensions ($\mathcal{H} = \mathbb{R}^n$), this reduces to a decomposition of a Hurwitz drift matrix into an exponentially stable analytic base plus a bounded perturbation.
\end{definition}

\begin{definition}[Trace-Class Noise and Covariance]
  \label{def:trace_class}
  Let $\mathcal{S}_1(\mathcal{H})$ denote the Banach space of trace-class (nuclear) operators on $\mathcal{H}$ equipped with the trace norm $\|T\|_{\mathcal{S}_1} \coloneqq \tr(|T|) = \sum_{i=1}^\infty \langle \phi_i, (T^*T)^{1/2} \phi_i \rangle < \infty$.
  \begin{enumerate}
    \item[(i)] The diffusion operator $\mathcal{D}: \mathcal{H} \to \mathcal{H}$ is positive, self-adjoint, and belongs to $\mathcal{S}_1(\mathcal{H})$.
    \item[(ii)] Assuming the pair $(\mathcal{A}, \mathcal{D}^{1/2})$ is controllable, the stationary covariance operator $\Sigma: \mathcal{H} \to \mathcal{H}$ is the unique strictly positive, self-adjoint, trace-class solution to the operator Lyapunov equation:
    \begin{equation}
      \label{eq:op_lyapunov}
      \mathcal{A}\Sigma + \Sigma\mathcal{A}^* + \mathcal{D} = 0
    \end{equation}
    which is given by the convergent Bochner integral
    \[
      \Sigma = \int_0^\infty e^{t\mathcal{A}} \mathcal{D} e^{t\mathcal{A}^*} dt.
    \]
  \end{enumerate}
\end{definition}

\begin{assumption}[Core Analytical Assumptions for Singular Hyperplane Conditioning]
  This work relies on two core analytical assumptions regarding the target-coordinate block decoupling and the outgoing coupling:
  \begin{enumerate}
    \item[(i)] \textbf{Linear Second-Moment Channel:} \label{assum:linear_second_moment} A centered linear stationary process on $\mathcal{H}=\mathcal{H}_k\oplus \mathcal{H}_I$ has finite second moments and covariance solving \Cref{eq:op_lyapunov}. The block-decoupling constraint zeroes all incoming target drift and cross-noise, so that $\mathcal{A}_{kI}=\mathcal{D}_{kI}=\mathcal{D}_{Ik}=0$ and $\mathcal{A}_{kk}=-\lambda$. The target diffusion $\mathcal{D}_{kk}=\varepsilon_0/\lambda$ is the unique scaling determined by the canonical normalization (Proposition~\ref{prop:exact_calibration}). The free noise is independent of the local target noise. (Gaussianity is not assumed).
    \item[(ii)] \textbf{Cameron--Martin Stability of Off-Diagonal Drift Coupling:} \label{assum:cm_stability} Let $Q_0\coloneqq\Sigma_{II}^{(0)}$ denote the free-block covariance produced by the free noise under block decoupling, and let $H_{Q_0}\coloneqq\operatorname{Range}(Q_0^{1/2})$ with the covariance norm $|u|_{Q_0}\coloneqq\|Q_0^{-1/2}u\|_{\mathcal{H}_I}$. The outgoing vector $a_{Ik}\coloneqq\mathcal{A}_{Ik}\phi_k$ belongs to $H_{Q_0}$, and $e^{t\mathcal{A}_{II}}$ restricts to a strongly continuous exponentially stable semigroup on $H_{Q_0}$:
    \begin{equation}
      \|e^{t\mathcal{A}_{II}}\|_{\mathcal{L}(H_{Q_0})} \le M e^{-\omega t},\qquad t\ge0,
    \end{equation}
    for constants $M\ge1$ and $\omega>0$ independent of $\lambda$.
  \end{enumerate}
  \vspace{1mm}
  \noindent\textit{Intuitive Explanation:} Part (i) states that the regularized dynamics act as a strong localized drift penalization on the target coordinate $X_k$ while isolating it from incoming influences, giving an operator implementation of the block-decoupling constraint. Part (ii) requires that the off-diagonal drift coupling $a_{Ik}$ from the target coordinate has finite covariance norm relative to the background fluctuations of the free coordinates. This ensures that target fluctuations propagating downstream do not overload the variance of the infinite-dimensional system.
\end{assumption}

\begin{remark}[Finite-Norm Outgoing Covariance Coupling and Finite-Dimensional Calibration]
  \label{rem:cm_information_flow}
  The membership $a_{Ik}\in H_{Q_0}$ is an admissibility condition,
  not a consequence of the free-block Lyapunov equation alone: an arbitrary
  outgoing vector need not lie in
  $\operatorname{Range}((\Sigma_{II}^{(0)})^{1/2})$.  It says that the
  deterministic channel carrying target fluctuations into the free-block has
  finite covariance norm relative to the free-noise background.
  
  To understand this condition intuitively, consider a finite-dimensional setting where the free-block state space is $\mathcal{H}_I = \mathbb{R}^{d-1}$ and the free stationary covariance $Q_0 \coloneqq \Sigma_{II}^{(0)}$ is strictly positive-definite. Since $Q_0 \succ 0$, its square root $Q_0^{1/2}$ is invertible, and its range is the entire space $\mathbb{R}^{d-1}$. In this case, the Cameron--Martin space is the full space $H_{Q_0} = \operatorname{Range}(Q_0^{1/2}) = \mathbb{R}^{d-1}$, and the Cameron--Martin norm $|u|_{Q_0} = \|Q_0^{-1/2} u\|_2$ is equivalent to the standard Euclidean norm. Consequently, the admissibility condition $a_{Ik} \in H_{Q_0}$ is trivially satisfied for any outgoing vector $a_{Ik} \in \mathbb{R}^{d-1}$.
  
  However, if some direction in the free coordinates is noise-free under block decoupling (meaning $Q_0$ is positive semidefinite but singular), $H_{Q_0}$ becomes a proper subspace of $\mathbb{R}^{d-1}$. In this case, the condition $a_{Ik} \in H_{Q_0}$ requires that the leakage vector $a_{Ik}$ has no component in the null space of $Q_0$. Equivalently, it prevents the target coordinate from leaking fluctuations into a downstream coordinate that has zero background noise, which would otherwise result in an infinite relative covariance norm.
  
  In infinite dimensions, $\Sigma_{II}^{(0)}$ is trace-class and compact, so its eigenvalues decay to zero, making $\operatorname{Range}((\Sigma_{II}^{(0)})^{1/2})$ a dense but strictly proper subspace of $\mathcal{H}_I$. The condition $a_{Ik} \in H_{Q_0}$ requires the off-diagonal drift coupling coefficients to decay sufficiently fast relative to the background noise eigenvalues to keep the relative covariance norm finite.
  
  Under this assumption, invariance and exponential stability on $H_{Q_0}$ imply:
  \begin{equation}
  \begin{aligned}
    (\lambda I-\mathcal{A}_{II})^{-1}a_{Ik}
       &\in H_{Q_0},\\
    |(\lambda I-\mathcal{A}_{II})^{-1}a_{Ik}|_{Q_0}
       &\le \frac{M}{\lambda+\omega}|a_{Ik}|_{Q_0}.
  \end{aligned}
  \end{equation}
\end{remark}

\begin{remark}[Constant Dependency in LMMSE Expansion]
  Consequently, the exact leakage cross-covariance
  $\Sigma_{\lambda,Ik}$ has finite Cameron--Martin norm and is
  $O(\lambda^{-3})$ in that norm.  This is a condition on the deterministic
  response in the covariance norm; it does not assert that infinite-dimensional
  Gaussian sample paths lie in $H_{Q_0}$ almost surely.
\end{remark}

% ============================================================
%  §3  HILBERT-SPACE VARIATIONAL SCHUR RESIDUAL
% ============================================================
\section{Hilbert-Space Variational Schur Residual}
\label{sec:infdim}

\subsection{Variational Schur Subtraction and Norm Bounds}
\label{sec:infdim:variational}

In finite dimensions, the Schur complement of the bulk covariance block $\Sigma_{II}$ is defined as $S_k \coloneqq \Sigma_{kk} - \Sigma_{kI} \Sigma_{II}^{-1} \Sigma_{Ik}$. Under stable and controllable dynamics, the bulk covariance $\Sigma_{II}$ is positive-definite and invertible, so the subtraction is well-defined. However, in infinite dimensions, the stationary covariance operator is trace-class and compact, meaning its restriction $\Sigma_{\lambda,II}$ lacks a bounded global inverse. The appropriate operator-theoretic replacement is the shorted operator: for a non-negative block operator, the Schur complement is recovered as the quadratic form of the short to the target subspace, with the inverse term replaced by a variational supremum \cite{krein1947selfadjoint,anderson1975shorted}. To lift the Schur complement to infinite dimensions in the present singular-conditioning problem, we formulate the projection residual directly using covariance-norm forms.

Let $Q_0 \coloneqq \Sigma_{II}^{(0)}$ denote the fixed background bulk covariance under full decoupling ($\lambda = \infty$). We define the core Cameron--Martin space $H_{Q_0} \coloneqq \operatorname{Range}(Q_0^{1/2})$ and the core covariance norm:
\begin{equation}
  |u|_{Q_0} \coloneqq \|Q_0^{-1/2} u\| \quad \text{for all } u \in H_{Q_0}.
\end{equation}
By Assumption~\ref{assum:cm_stability}, the off-diagonal drift coupling vector satisfies $a_{Ik} \in H_{Q_0}$, and the restricted drift generator $\mathcal{A}_{II}$ generates an exponentially stable semigroup on $H_{Q_0}$.

Under this assumption, the exact cross-covariance resolvent formula (derived in \Cref{app:block_cumulant_analysis}) satisfies:
\begin{equation}
  \Sigma_{\lambda,Ik} = \frac{\varepsilon_0}{2\lambda^2} (\lambda I - \mathcal{A}_{II})^{-1} a_{Ik}.
\end{equation}
This resolvent response yields the Cameron--Martin norm bound on the cross-covariance block:
\begin{equation}
  \label{eq:cross_norm_bound}
  |\Sigma_{\lambda,Ik}|_{Q_0} = O(\lambda^{-3}) \quad \text{as } \lambda \to \infty.
\end{equation}

To define the infinite-dimensional projection residual, we introduce the variational Schur subtraction $R_\lambda$:
\begin{equation}
  \label{eq:variational_schur_subtraction}
  R_\lambda \coloneqq \sup_{\langle u, \Sigma_{\lambda,II} u \rangle > 0} \frac{|\langle u, \Sigma_{\lambda,Ik} \rangle|^2}{\langle u, \Sigma_{\lambda,II} u \rangle}.
\end{equation}
Equivalently, if $\Sigma_\lambda$ is viewed as a non-negative block operator on $\mathbb{C}\phi_k\oplus\mathcal{H}_I$, then \eqref{eq:variational_schur_subtraction} is the one-dimensional Anderson--Trapp shorted-operator subtraction, namely
\begin{equation}
  \label{eq:shorted_operator_scalar}
  \big\langle \phi_k,\mathcal{S}_{\mathbb{C}\phi_k}(\Sigma_\lambda)\phi_k\big\rangle
  =\Sigma_{\lambda,kk}-R_\lambda.
\end{equation}
The restriction to vectors with $\langle u,\Sigma_{\lambda,II}u\rangle>0$ simply removes null directions of the denominator; for a non-negative block operator, such directions do not change the supremum. Thus $R_\lambda$ is not an ambient inverse computation but the covariance-norm contribution removed by shorting to the target coordinate \cite{anderson1975shorted,owhadi2015conditioning}.

Since the target noise channel and bulk noise channel are independent, the bulk state decomposes into independent background and target-driven components, meaning the bulk covariance dominates the background covariance ($\Sigma_{\lambda,II} = Q_0 + \operatorname{Cov}(Y_I) \succeq Q_0$). Douglas range inclusion then guarantees that the variational subtraction is bounded:
\begin{equation}
  0 \le R_\lambda \le |\Sigma_{\lambda,Ik}|_{Q_0}^2 = O(\lambda^{-6}).
\end{equation}

\subsection{Hilbert-Space Schur Residual Scaling}
\label{sec:infdim:scaling}

We now establish the main scaling theorem for the $L^2$-projection residual.

\begin{theorem}[Hilbert Schur Residual Invariant]
  \label{thm:inf_cancellation}
  Under Assumptions~\ref{assum:linear_second_moment} and \ref{assum:cm_stability}, the $L^2$-projection residual satisfies:
  \begin{equation}
    \label{eq:main_schur_sharp}
    S_{k,\lambda} \coloneqq \|X_k - P_{\mathcal{M}_I} X_k\|_{L^2}^2 = \frac{\varepsilon_0}{2\lambda^2} - R_\lambda,
  \end{equation}
  where the remainder satisfies $0 \le R_\lambda = O(\lambda^{-6})$. In particular, the Schur LMMSE invariant is:
  \begin{equation}
    \label{eq:m2_schur_value}
    m_2^{\mathrm{Schur}} = \lim_{\lambda\to\infty} \lambda^2 S_{k,\lambda} = \frac{\varepsilon_0}{2},
  \end{equation}
  and the inverse projection residual scaling satisfies:
  \begin{equation}
    \label{eq:schur_scaling_inverse}
    S_{k,\lambda}^{-1} = \frac{2}{\varepsilon_0}\lambda^2 + O(\lambda^{-2}).
  \end{equation}
  This result depends only on covariance and orthogonal projection; it holds for any linear stationary model satisfying these assumptions, without Gaussianity.
\end{theorem}
\begin{proof}
  See \Cref{app:block_cumulant_analysis} for the full derivation.
\end{proof}

\subsection{Convergence of Finite-Dimensional Truncations}
\label{sec:proofs:convergence}

Let $\Pi_n: \mathcal{H} \to \mathcal{H}_n \coloneqq \operatorname{span}\{\phi_1, \ldots, \phi_n\}$ be the $n$-dimensional Galerkin projection. Write $\mathcal{A}_n^{(\lambda)} \coloneqq \Pi_n \mathcal{A}^{(\lambda)} \Pi_n$, $\mathcal{D}_n^{(\lambda)} \coloneqq \Pi_n \mathcal{D}^{(\lambda)} \Pi_n$, and let $\Sigma_{n,\lambda}$ solve the truncated Lyapunov equation $\mathcal{A}_n^{(\lambda)} \Sigma_{n,\lambda} + \Sigma_{n,\lambda} (\mathcal{A}_n^{(\lambda)})^* + \mathcal{D}_n^{(\lambda)} = 0$.

\begin{proposition}[Galerkin Consistency of the Schur Invariant]
  \label{prop:galerkin_consistency}
  Under Assumption~\ref{assum:linear_second_moment}, let $\Pi_n:\mathcal{H}\to\mathcal{H}_n$ be finite-rank orthogonal Galerkin projections preserving the target-bulk block decomposition:
  \begin{equation}
    \Pi_n\phi_k=\phi_k,\qquad \Pi_n\mathcal{H}_I\subset\mathcal{H}_I,\qquad \Pi_n\to I\ \text{strongly on }\mathcal{H}.
  \end{equation}
  Assume the truncated regularized dynamics satisfy the uniform exponential stability and trace-class conditions stated in Proposition~\ref{prop:dynamic_galerkin} of \Cref{app:block_cumulant_analysis}.  Then the truncated covariance operators converge in trace norm:
  \begin{equation}
    \|\Sigma_{n,\lambda}-\Sigma_\lambda\|_{\mathcal{S}_1}\to0\quad\text{as }n\to\infty,
  \end{equation}
  and the truncated Schur invariant is exact at every level:
  \begin{equation}
    m_2^{\mathrm{Schur},(n)} \coloneqq \lim_{\lambda\to\infty} \lambda^2 S_{k,\lambda}^{(n)} = \frac{\varepsilon_0}{2} \quad \text{for all } n \ge k.
  \end{equation}
\end{proposition}
\begin{proof}
  See \Cref{app:block_cumulant_analysis} for the dynamic and spectral support.
\end{proof}

% ============================================================
%  §4  GAUSSIAN CONDITIONAL EXTENSION AND FINITE-PART FUNCTIONAL
% ============================================================
\section{Gaussian Conditional Extension and Finite-Part Functional}
\label{sec:gaussian_zero}

This section gives the Gaussian extension of the base theorem, using Radon Gaussian laws and disintegration to formulate conditioning under exact evidence and define the renormalized finite-part functional.

\subsection{Radon Gaussian Setup and Disintegration}
\label{sec:gaussian:setup}

\begin{assumption}[Hilbert Gaussian Conditioning and Affine Mean Channel]
  \label{assum:gaussian_clamp}
  In addition to Assumptions~\ref{assum:linear_second_moment} and \ref{assum:cm_stability}, let $\mu_\lambda = \mathcal{N}(m_\lambda, \Sigma_\lambda)$ be a Gaussian Radon law on $\mathcal{H}$ with injective trace-class covariance. Let $\mu_{\mathrm{obs}}$ be the canonical Gaussian disintegration member on the evidence hyperplane $X_k = x_k^*$. Suppose that a fixed deterministic forcing $b \in \mathcal{H}$ gives the stationary mean equation:
  \begin{equation}
    \mathcal{A}_{\cond} m_\lambda + b - \lambda \Pi_k (m_\lambda - x_k^* \phi_k) = 0, \quad \Pi_k \mathcal{A}_{\cond} = 0.
  \end{equation}
  We define the target deterministic offset constant:
  \begin{equation}
    c_\lambda \coloneqq \lambda(m_{\lambda,k} - x_k^*) = b_k,
  \end{equation}
  where $m_{\lambda,k} \coloneqq \langle m_\lambda, \phi_k \rangle$ and $b_k \coloneqq \langle b, \phi_k \rangle$.
\end{assumption}

Under the Gaussian Radon law, the coordinate projection $\pi_k: \mathcal{H} \to \mathbb{R}$ is continuous, so the target variable $X_k$ is a well-defined Borel random variable. The Polish structure of the Hilbert space guarantees the existence of a regular conditional probability distribution. 
By selecting the weakly continuous version of the conditional law at the single point $t = x_k^*$ (as detailed in \Cref{app:gaussian_lift}), we obtain the unique pointwise disintegration member $\mu_{\mathrm{obs}}$. Under this measure, the target coordinate satisfies the exact target zeros:
\begin{equation}
  \label{eq:exact_target_zeros}
  m_{\text{obs},k} - x_k^* = 0, \quad \Var_{\text{obs}}(X_k) = 0 \quad (\text{exact}).
\end{equation}

\subsection{Global-Inverse-Free Measurable Predictor}
\label{sec:gaussian:predictor}

Let $C_\lambda \coloneqq \Sigma_{\lambda,II}$ denote the Gaussian bulk covariance operator on $\mathcal{H}_I$. Because $C_\lambda$ is trace-class and compact, it lacks a bounded global inverse. We define the bulk Cameron--Martin space $H_{C_\lambda} \coloneqq \operatorname{Range}(C_\lambda^{1/2})$ and the corresponding covariance norm.

The leakage cross-covariance vector satisfies $\Sigma_{\lambda,Ik} \in H_{C_\lambda}$ with norm $\|\Sigma_{\lambda,Ik}\|_{H_{C_\lambda}} = O(\lambda^{-3})$. This range inclusion allows us to define the conditional predictor $\mu_{k|I}(X_I) \coloneqq \E[X_k \mid X_I]$ as a measurable Wiener functional:
\begin{equation}
  \mu_{k|I}(X_I) = m_{\lambda,k} + W_{\Sigma_{\lambda,Ik}}(X_I - m_{\lambda,I}),
\end{equation}
which is the unique LMMSE predictor in $L^2$. The estimates in \Cref{app:gaussian_lift} imply that, for all sufficiently large $\lambda$, the bulk marginals $\mu_{\text{obs},I}$ and $\mu_{\lambda,I}$ are equivalent and have finite bulk relative entropy:
\begin{equation}
  \KL(\mu_{\text{obs},I} \parallel \mu_{\lambda,I}) = O(\lambda^{-4}) \to 0 \quad \text{as } \lambda \to \infty.
\end{equation}

\subsection{Gaussian Renormalized Finite-Part Functional}
\label{sec:gaussian:functional}

Although the joint KL divergence diverges under the exact conditioning, we can define a renormalized finite-part action.

\begin{definition}[Gaussian Renormalized Finite-Part Functional]
  \label{def:unified_action}
  Whenever $S_{k,\lambda}>0$ and $\mu_{\text{obs},I}$ is equivalent to $\mu_{\lambda,I}$, we define the Gaussian renormalized finite-part functional of the conditioned law:
  \begin{equation}
    \label{eq:unified_action}
    \begin{aligned}
      \mathfrak{J}_\lambda(\mu_{\mathrm{obs}}; \mu_\lambda) \coloneqq &\KL(\mu_{\mathrm{obs},I} \parallel \mu_{\lambda,I}) \\
      &+ \frac{1}{2} S_{k,\lambda}^{-1} \E_{\mu_{\mathrm{obs}}}\left[ \left( X_k - \mu_{k|I}(X_I) \right)^2 \right].
    \end{aligned}
  \end{equation}
  This represents the finite part of the joint KL divergence under the conditional normalization scheme, where the Dirac and Gaussian normalization singularities are subtracted.
\end{definition}

\paragraph{Weak scheme-independence.}
For an evidence level $t\in\mathbb{R}$, write $\mu_\lambda^t$ for the weakly continuous conditional law from \Cref{app:gaussian_lift} and set
\[
  \mathfrak{J}_\lambda(t)\coloneqq
  \mathfrak{J}_\lambda(\mu_\lambda^t;\mu_\lambda).
\]
Then, for any two evidence levels $t_1,t_2$ for which the finite part is defined,
\begin{equation}
  \label{eq:weak_scheme_independence}
  \mathfrak{J}_\lambda(t_2)-\mathfrak{J}_\lambda(t_1)
  =
  \frac{(t_2-m_{\lambda,k})^2-(t_1-m_{\lambda,k})^2}
       {2\Sigma_{\lambda,kk}}.
\end{equation}
In the calibrated scaling $\Sigma_{\lambda,kk}=\varepsilon_0/(2\lambda^2)$, if
$b_i=\lambda(m_{\lambda,k}-t_i)$, this becomes
\[
  \mathfrak{J}_\lambda(t_2)-\mathfrak{J}_\lambda(t_1)
  =
  \frac{b_2^2-b_1^2}{\varepsilon_0}.
\]
\begin{proof}
Let $\Delta_t\coloneqq t-m_{\lambda,k}$ and
\[
  \rho_\lambda
  \coloneqq
  \frac{\|\Sigma_{\lambda,Ik}\|_{H_{C_\lambda}}^2}{\Sigma_{\lambda,kk}}
  =
  1-\frac{S_{k,\lambda}}{\Sigma_{\lambda,kk}}.
\]
The Gaussian entropy computation in \Cref{app:gaussian_lift} gives
\[
  \KL(\mu_{\lambda,I}^t\parallel\mu_{\lambda,I})
  =
  \frac{1}{2}\frac{\Delta_t^2}{\Sigma_{\lambda,kk}}\rho_\lambda
  -\frac{1}{2}\log(1-\rho_\lambda)-\frac{1}{2}\rho_\lambda .
\]
The conditional-residual term in \eqref{eq:unified_action} is
\[
  \frac{1}{2S_{k,\lambda}}
  \E_{\mu_\lambda^t}\!\left[
    (X_k-\mu_{k|I}(X_I))^2
  \right]
  =
  \frac{1}{2}\rho_\lambda
  +
  \frac{1}{2}\frac{\Delta_t^2}{\Sigma_{\lambda,kk}}(1-\rho_\lambda).
\]
Adding these two identities cancels the $\rho_\lambda/2$ terms and yields
\[
  \mathfrak{J}_\lambda(t)
  =
  \frac{(t-m_{\lambda,k})^2}{2\Sigma_{\lambda,kk}}
  -\frac{1}{2}\log(1-\rho_\lambda).
\]
The logarithmic term depends on the reference covariance and the conditioning
scale, but not on the evidence value $t$. Hence it cancels in the difference
between $t_2$ and $t_1$, proving \eqref{eq:weak_scheme_independence}. The same
argument shows weak independence from the mollifying kernel: replacing the
Gaussian Dirac mollifier in \Cref{prop:mollification_limit} by any centered
approximate identity with finite second moment changes the extracted finite
part only by an additive term depending on the kernel, $\lambda$, and
$S_{k,\lambda}$, not on $t$. Therefore all evidence differences
$\Delta\mathfrak{J}_\lambda$ are kernel-independent.
\end{proof}

\subsection{Gaussian Main Theorem}
\label{sec:gaussian:theorem}

We now establish the main regularization and cancellation theorems for the finite-part functional.

\begin{theorem}[Gaussian Finite-Part Regularization]
  \label{thm:finite_part_regularization}
  Under Assumption~\ref{assum:gaussian_clamp}, there exists $\lambda_0>0$ such that, for $\lambda\ge\lambda_0$, the renormalized finite-part functional under the conditioned law $\mu_{\mathrm{obs}}$ is well-defined and satisfies:
  \begin{equation}
    \label{eq:unified_regularization}
    \mathfrak{J}_\lambda(\mu_{\mathrm{obs}}; \mu_\lambda) = \frac{b_k^2}{\varepsilon_0} + O(\lambda^{-4}).
  \end{equation}
  In particular, the functional remains bounded at $O(1)$ and converges to:
  \begin{equation}
    \lim_{\lambda\to\infty} \mathfrak{J}_\lambda(\mu_{\mathrm{obs}}; \mu_\lambda) = \frac{b_k^2}{\varepsilon_0}.
  \end{equation}
\end{theorem}
\begin{proof}
  See \Cref{app:gaussian_lift} for the full derivation.
\end{proof}

\begin{corollary}[Exact Decoupled-Centered Cancellation]
  \label{cor:exact_pole_zero}
  Under Assumption~\ref{assum:gaussian_clamp}, if the target coordinate is dynamically decoupled from the bulk ($a_{Ik} = 0$) and the evidence is centered at the stationary mean ($x_k^* = m_{\lambda,k}$), then:
  \begin{equation}
    \mathfrak{J}_\lambda(\mu_{\mathrm{obs}}; \mu_\lambda) = 0 \quad \text{for every } \lambda > 0.
  \end{equation}
\end{corollary}
\begin{proof}
  The decoupling implies $\Sigma_{\lambda,Ik} = 0$, so $\KL(\mu_{\text{obs},I} \parallel \mu_{\lambda,I}) = 0$. The centering makes the conditional residual zero under the conditioning, so both terms in \eqref{eq:unified_action} vanish.
\end{proof}

% ============================================================
%  §5  ELLIPTIC GAUSSIAN FIELD EXAMPLE
% ============================================================
\section{Elliptic Gaussian Field Example}
\label{sec:spde}

We now record a standard elliptic Gaussian field example to show that the abstract Hilbert-space assumptions are nonempty and natural for covariance operators arising from stochastic evolution equations and elliptic Gaussian priors \cite{daprato2014stochastic,lord2014introduction,stuart2010inverse}.

\subsection{Elliptic Gaussian field on $[0,1]$}

Let
\[
  L=-\frac{d^2}{dx^2}+\kappa^2
\]
on $L^2(0,1)$ with Dirichlet boundary conditions and $\kappa>0$. Then $\mathcal{A}_0=-L$ with domain $H^2(0,1)\cap H_0^1(0,1)$ is dissipative, $L$ is positive self-adjoint and sectorial, and $e^{-tL}$ is an exponentially stable analytic contraction semigroup. The inverse covariance
\[
  Q_0=L^{-1}
\]
is compact, positive, self-adjoint, and trace class. With eigenfunctions $e_j(x)=\sqrt{2}\sin(\pi jx)$, its eigenvalues are
\[
  q_j=(\pi^2j^2+\kappa^2)^{-1}.
\]
Thus $Q_0$ is a typical infinite-dimensional covariance operator: it is compact and trace class, has no bounded inverse on the ambient space, and carries its inverse only as a Cameron--Martin quadratic form \cite{bogachev1998gaussian,simon2005trace}.

\subsection{Cameron--Martin range and covariance norm}

The Cameron--Martin space associated with $Q_0$ is
\[
  H_{Q_0}=\operatorname{Range}(Q_0^{1/2})=H_0^1(0,1),
\]
with equivalent covariance norm
\[
  \|h\|_{Q_0}^2
  =\langle h,Lh\rangle_{L^2}
  =\int_0^1 \bigl(|h'(x)|^2+\kappa^2|h(x)|^2\bigr)\,dx.
\]
This identifies the variational Schur projection in a Sobolev quadratic form, as in the standard covariance-norm representation of Gaussian Hilbert priors \cite{bogachev1998gaussian,stuart2010inverse}. The covariance inverse does not act as a bounded operator on $L^2(0,1)$; it acts only through this quadratic form on its Cameron--Martin domain.

\subsection{Spectral sensor conditioning}

Fix a target eigenmode $e_k$ and define the observed coordinate
\[
  X_k=\langle X,e_k\rangle_{L^2}.
\]
The unperturbed target variance is
\[
  \Sigma_{0,kk}=q_k=(\pi^2k^2+\kappa^2)^{-1}.
\]
A regularized restoring drift on this coordinate produces the exact normalization
\[
  \Sigma_{\lambda,kk}=\frac{\varepsilon_0}{2\lambda^2}-R_\lambda,
  \qquad R_\lambda=O(\lambda^{-6}),
\]
under the block assumptions of \Cref{sec:infdim}. Off-diagonal bounded perturbations of the drift operator generate the cross-covariance vector entering the variational Schur residual. The resulting subtraction is evaluated by the Cameron--Martin quadratic form above, not by a global inverse on $L^2(0,1)$.

This example shows that the abstract theorem applies to an infinite-dimensional Gaussian field: the residual is controlled by a Sobolev norm on $H_0^1(0,1)$ while the ambient covariance remains compact and non-invertible.

\section{Discussion}
\label{sec:discussion}

We have shown that singular Gaussian coordinate conditioning in a separable Hilbert space admits an inverse-free Schur normal form.  The main obstruction is the absence of a bounded inverse for compact covariance blocks.  The variational Cameron--Martin projection is the target-coordinate shorted-operator subtraction evaluated in the covariance-energy scale, and it yields the residual estimate needed for the hard-conditioning limit.  The scalar nature of the target coordinate is essential to the present proof.  Three extensions mark the boundary between the theorem proved here and a genuinely finite-rank conditioning theory. We formalize these extensions as exact theorems below.

\subsection{Matrix algebraic extension}
\label{sec:discussion:matrix-extension}

For an $m$-dimensional observed subspace $E=\operatorname{span}\{\phi_1,\ldots,\phi_m\}$, the scalar target block $\Sigma_{\lambda,kk}$ is replaced by the covariance matrix $\Sigma_{\lambda,EE}\in\mathbb{R}^{m\times m}$. The scalar calibration becomes a matrix Lyapunov balancing problem.

\begin{theorem}[Matrix Lyapunov Calibration]
  \label{thm:matrix_lyapunov}
  Let $E$ be an $m$-dimensional target subspace. Suppose the regularized drift and diffusion on $E$ are given by $A_{\lambda,E} = -\lambda \Lambda_E$ and $D_{\lambda,E} = \frac{1}{\lambda} \Delta_E$, where $\Lambda_E, \Delta_E \in \mathbb{R}^{m \times m}$ are strictly positive definite matrices. Let $\Sigma_{\lambda,EE}$ be the unique positive definite solution to the matrix Lyapunov equation:
  \begin{equation}
    \label{eq:discussion_matrix_lyapunov}
    A_{\lambda,E}\Sigma_{\lambda,EE}
    +\Sigma_{\lambda,EE}A_{\lambda,E}^{*}
    +D_{\lambda,E}=0.
  \end{equation}
  Then $\Sigma_{\lambda,EE}$ scales exactly as $\lambda^{-2}$. Specifically, the normalized matrix $\widetilde{\Sigma}_{EE} \coloneqq \lambda^2 \Sigma_{\lambda,EE}$ is independent of $\lambda$ and is the unique positive definite solution to the scale-free Lyapunov equation:
  \begin{equation}
    \label{eq:scale_free_lyapunov}
    \Lambda_E \widetilde{\Sigma}_{EE} + \widetilde{\Sigma}_{EE} \Lambda_E^* = \Delta_E.
  \end{equation}
\end{theorem}
\begin{proof}
  Substituting $A_{\lambda,E} = -\lambda \Lambda_E$ and $D_{\lambda,E} = \frac{1}{\lambda} \Delta_E$ into \eqref{eq:discussion_matrix_lyapunov} yields:
  \[
    -\lambda \Lambda_E \Sigma_{\lambda,EE} - \lambda \Sigma_{\lambda,EE} \Lambda_E^* + \frac{1}{\lambda} \Delta_E = 0.
  \]
  Multiplying the entire equation by $\lambda$ gives:
  \[
    \Lambda_E (\lambda^2 \Sigma_{\lambda,EE}) + (\lambda^2 \Sigma_{\lambda,EE}) \Lambda_E^* = \Delta_E.
  \]
  Defining $\widetilde{\Sigma}_{EE} \coloneqq \lambda^2 \Sigma_{\lambda,EE}$, we recover \eqref{eq:scale_free_lyapunov}. Since $\Lambda_E$ is positive definite, its spectrum lies in the open right half-plane, making $-\Lambda_E$ Hurwitz. Because $\Delta_E$ is positive definite, standard continuous Lyapunov theory guarantees that \eqref{eq:scale_free_lyapunov} admits a unique positive definite solution $\widetilde{\Sigma}_{EE}$. Therefore, the exact scaling $\Sigma_{\lambda,EE} = \lambda^{-2} \widetilde{\Sigma}_{EE}$ holds for all $\lambda > 0$, properly generalizing the scalar calibration condition $\Sigma_{\lambda,kk} = \varepsilon_0 / (2\lambda^2)$.
\end{proof}

\subsection{Operator-valued residuals}
\label{sec:discussion:operator-valued-residuals}

For an $m$-dimensional target space $E$, the Schur residual must be a positive-semidefinite operator on $E$. Because the free-block covariance $\Sigma_{\lambda,II}$ is compact and lacks a bounded global inverse, the formal finite-dimensional matrix expression $R_{\lambda,E} = \Sigma_{\lambda,EI}\Sigma_{\lambda,II}^{-1}\Sigma_{\lambda,IE}$ must be rigorously replaced by the shorted-operator construction evaluated in the Cameron--Martin energy norm.

\begin{theorem}[Operator-Valued Schur Residual]
  \label{thm:operator_residual}
  Let $\mathcal{H} = E \oplus \mathcal{H}_I$, with $E$ finite-dimensional. Let $\Sigma_\lambda$ be the non-negative trace-class covariance operator on $\mathcal{H}$ with blocks $\Sigma_{\lambda,EE}$, $\Sigma_{\lambda,EI}$, $\Sigma_{\lambda,IE}$, and $\Sigma_{\lambda,II}$. Define the shorted-operator residual $R_{\lambda,E}$ on $E$ by its quadratic forms:
  \begin{equation}
    \label{eq:discussion_operator_residual_form}
    \langle v,R_{\lambda,E}v\rangle_E
    =
    \sup_{\langle u,\Sigma_{\lambda,II}u\rangle>0}
    \frac{|\langle u,\Sigma_{\lambda,IE}v\rangle|^2}
         {\langle u,\Sigma_{\lambda,II}u\rangle},
    \qquad v\in E.
  \end{equation}
  Assume the structural range inclusion $\operatorname{Range}(\Sigma_{\lambda,IE}) \subset \operatorname{Range}(\Sigma_{\lambda,II}^{1/2}) \eqqcolon H_{C_\lambda}$. Then $R_{\lambda,E}$ is a well-defined bounded positive-semidefinite operator on $E$, uniquely characterized by:
  \begin{equation}
    R_{\lambda,E} = (\Sigma_{\lambda,II}^{-1/2}\Sigma_{\lambda,IE})^* (\Sigma_{\lambda,II}^{-1/2}\Sigma_{\lambda,IE}).
  \end{equation}
  Furthermore, if the cross-covariance satisfies the energy bound $\|\Sigma_{\lambda,IE} v\|_{H_{C_\lambda}} \le M \lambda^{-3} \|v\|_E$ for all $v \in E$ and some constant $M > 0$, then the operator norm of the residual satisfies $\|R_{\lambda,E}\|_{\mathcal{L}(E)} = O(\lambda^{-6})$.
\end{theorem}
\begin{proof}
  By the assumed range inclusion, for any $v \in E$, the vector $y_v \coloneqq \Sigma_{\lambda,IE} v$ lies in $H_{C_\lambda} = \operatorname{Range}(\Sigma_{\lambda,II}^{1/2})$. Thus, there exists a unique element $w_v \in \overline{\operatorname{Range}(\Sigma_{\lambda,II}^{1/2})}$ such that $\Sigma_{\lambda,II}^{1/2} w_v = y_v$. We denote $w_v \coloneqq \Sigma_{\lambda,II}^{-1/2} \Sigma_{\lambda,IE} v$.
  
  For any $u \in \mathcal{H}_I$ with $\langle u, \Sigma_{\lambda,II} u \rangle > 0$, the ratio in \eqref{eq:discussion_operator_residual_form} can be rewritten as:
  \[
    \frac{|\langle u, \Sigma_{\lambda,II}^{1/2} w_v \rangle|^2}{\|\Sigma_{\lambda,II}^{1/2} u\|^2} = \frac{|\langle \Sigma_{\lambda,II}^{1/2} u, w_v \rangle|^2}{\|\Sigma_{\lambda,II}^{1/2} u\|^2}.
  \]
  By the Cauchy--Schwarz inequality, this ratio is bounded above by $\|w_v\|^2$. The supremum is achieved by taking a sequence of vectors $\Sigma_{\lambda,II}^{1/2} u$ converging to $w_v$, yielding:
  \[
    \langle v, R_{\lambda,E} v \rangle_E = \|w_v\|^2 = \|\Sigma_{\lambda,II}^{-1/2} \Sigma_{\lambda,IE} v\|^2 = \langle v, (\Sigma_{\lambda,II}^{-1/2}\Sigma_{\lambda,IE})^* (\Sigma_{\lambda,II}^{-1/2}\Sigma_{\lambda,IE}) v \rangle_E.
  \]
  Since this holds for all $v \in E$, it uniquely defines $R_{\lambda,E}$ as a positive-semidefinite matrix on $E$. Finally, the assumption $\|\Sigma_{\lambda,IE} v\|_{H_{C_\lambda}} \le M \lambda^{-3} \|v\|_E$ means exactly that $\|w_v\| \le M \lambda^{-3} \|v\|_E$. Therefore, $\langle v, R_{\lambda,E} v \rangle_E \le M^2 \lambda^{-6} \|v\|_E^2$, which implies the operator norm bound $\|R_{\lambda,E}\|_{\mathcal{L}(E)} \le M^2 \lambda^{-6} = O(\lambda^{-6})$.
\end{proof}

\subsection{General observation operators and misalignment}
\label{sec:discussion:misalignment}

The scalar proof in this paper uses a block decomposition aligned with a single observed coordinate. However, typical finite-rank observation operators $L:\mathcal H\to\mathbb R^m$ select a subspace $E$ that does not commute with the free prior covariance, meaning the target-free split is not an eigenbasis split. This geometric misalignment generates cross-coupling that requires a precise resolvent analysis.

\begin{theorem}[Misaligned Covariance Coupling]
  \label{thm:misaligned_coupling}
  Let $P_E$ be an orthogonal projection onto an $m$-dimensional observed subspace $E$, and let $\mathcal{H}_I = E^\perp$. Let the regularized dynamics be decomposed such that the target block matches \Cref{thm:matrix_lyapunov} with $A_{\lambda,E} = -\lambda \Lambda_E$, the cross-noise is zero ($D_{\lambda,IE} = 0$), and the off-diagonal drift coupling is denoted by $A_{\lambda,IE}$. 
  If $A_{\lambda,IE}$ maps $E$ into the Cameron--Martin space $H_{Q_0}$ of the free background covariance $Q_0$ on $\mathcal{H}_I$, and the restricted drift $A_{II}$ generates an exponentially stable semigroup on $H_{Q_0}$, then the cross-covariance satisfies the exact identity:
  \begin{equation}
    \label{eq:misaligned_resolvent_identity}
    \Sigma_{\lambda,IE} = \int_0^\infty e^{t A_{II}} A_{\lambda,IE} \Sigma_{\lambda,EE} e^{t A_{\lambda,E}^*} dt.
  \end{equation}
  Furthermore, the misaligned cross-covariance energy scales as $\|\Sigma_{\lambda,IE} v\|_{H_{Q_0}} = O(\lambda^{-3})$ for any $v \in E$.
\end{theorem}
\begin{proof}
  The off-diagonal block of the global Lyapunov equation $\mathcal{A}^{(\lambda)} \Sigma_\lambda + \Sigma_\lambda (\mathcal{A}^{(\lambda)})^* + \mathcal{D}^{(\lambda)} = 0$ yields the Sylvester equation:
  \[
    A_{II} \Sigma_{\lambda,IE} + \Sigma_{\lambda,IE} A_{\lambda,E}^* + A_{\lambda,IE} \Sigma_{\lambda,EE} = 0.
  \]
  Since $A_{II}$ generates an exponentially stable semigroup $e^{t A_{II}}$ and $A_{\lambda,E} = -\lambda \Lambda_E$ is Hurwitz, the unique solution is given by the integral representation \eqref{eq:misaligned_resolvent_identity}.
  
  To bound the Cameron--Martin norm, fix $v \in E$. We have $e^{t A_{\lambda,E}^*} v = e^{-\lambda t \Lambda_E^*} v$. By \Cref{thm:matrix_lyapunov}, $\Sigma_{\lambda,EE} = O(\lambda^{-2})$. Since $A_{\lambda,IE}$ maps into $H_{Q_0}$ and $E$ is finite-dimensional, there exists a constant $C$ such that $\|A_{\lambda,IE} w\|_{H_{Q_0}} \le C \|w\|_E$ for all $w \in E$. Using the uniform exponential stability of $e^{t A_{II}}$ on $H_{Q_0}$, namely $\|e^{t A_{II}}\|_{\mathcal{L}(H_{Q_0})} \le M e^{-\omega t}$ for some $M \ge 1, \omega > 0$, we estimate:
  \begin{align*}
    \|\Sigma_{\lambda,IE} v\|_{H_{Q_0}} 
    &\le \int_0^\infty \|e^{t A_{II}}\|_{\mathcal{L}(H_{Q_0})} \|A_{\lambda,IE} \Sigma_{\lambda,EE} e^{-\lambda t \Lambda_E^*} v\|_{H_{Q_0}} dt \\
    &\le \int_0^\infty M e^{-\omega t} C \|\Sigma_{\lambda,EE}\|_{\mathcal{L}(E)} \|e^{-\lambda t \Lambda_E^*}\|_{\mathcal{L}(E)} \|v\|_E dt.
  \end{align*}
  Let $\gamma > 0$ be the smallest eigenvalue of the positive definite matrix $\Lambda_E^*$, so $\|e^{-\lambda t \Lambda_E^*}\|_{\mathcal{L}(E)} \le K e^{-\lambda \gamma t}$. Then:
  \begin{align*}
    \|\Sigma_{\lambda,IE} v\|_{H_{Q_0}} 
    &\le M C K \|\Sigma_{\lambda,EE}\|_{\mathcal{L}(E)} \|v\|_E \int_0^\infty e^{-(\omega + \lambda \gamma) t} dt \\
    &= M C K \|\Sigma_{\lambda,EE}\|_{\mathcal{L}(E)} \|v\|_E \frac{1}{\omega + \lambda \gamma}.
  \end{align*}
  Since $\|\Sigma_{\lambda,EE}\|_{\mathcal{L}(E)} = O(\lambda^{-2})$ and $(\omega + \lambda \gamma)^{-1} = O(\lambda^{-1})$, we conclude that $\|\Sigma_{\lambda,IE} v\|_{H_{Q_0}} = O(\lambda^{-3})$. Combined with \Cref{thm:operator_residual}, this verifies that the shorted-operator residual maintains the $O(\lambda^{-6})$ scaling even in the presence of observation misalignment.
\end{proof}

The result is therefore deliberately limited to covariance and Gaussian-conditioning questions in the rank-one aligned setting.  It does not assert a general non-Gaussian conditioning theory.  For non-Gaussian stationary processes, the projection residual $S_{k,\lambda}=\varepsilon_0/(2\lambda^2)-R_\lambda$ still has a second-moment interpretation, but Radon disintegration, finite-part functionals, and measure-class behavior require separate hypotheses \cite{bogachev2007measure,kallenberg2021foundations}.  Dynamical relaxation after conditioning, or transport distances between post-conditioning laws, belongs to a different problem.  The present contribution is the operator-theoretic construction of the singular Gaussian conditioning residual itself.

% ============================================================
%  \S6  CONCLUSION
% ============================================================
\section{Conclusion}
\label{sec:conclusion}

We established a Hilbert-space framework for singular coordinate conditioning of Gaussian measures. The central object is an inverse-free variational Schur projection on the Cameron--Martin range, equivalently the one-dimensional target-coordinate form of the shorted-operator subtraction. It gives canonical normalization, an $O(\lambda^{-6})$ residual estimate, Galerkin stability, and a Radon Gaussian conditioning statement. The elliptic field example shows that the assumptions are natural for trace-class covariance operators arising from stochastic evolution equations and elliptic Gaussian priors \cite{daprato2014stochastic,lord2014introduction,stuart2010inverse}.

\section*{Declaration of Generative AI and AI-Assisted Technologies in the Manuscript Preparation Process}

During the preparation of this work, the author(s) used Google DeepMind's Antigravity assistant in order to merge the LaTeX manuscripts, clean up cross-referencing tags, format mathematical notation, and edit the draft for style consistency. After using this tool/service, the author(s) reviewed and edited the content as needed and take(s) full responsibility for the content of the published article.

% ============================================================
%  REFERENCES
% ============================================================
\bibliographystyle{amsplain}
\input{paper1.bbl}

% ============================================================
%  APPENDIX
% ============================================================
\appendix

\input{paper1_appendix}
\end{document}


% ============================================================
%  APPENDIX
% ============================================================
\appendix

\section{Block Lyapunov Identities for the Calibrated Clamp}
\label{app:block_identities}

This appendix isolates the algebra used in the paper's shorting argument.  The
only inputs are the block-triangular calibrated operator family, the trace-class
stationary covariance, and the Cameron--Martin finite-energy assumption on the
outgoing drift channel.

\subsection{Block identities}

Write $\mathcal{H}=\mathcal{H}_k\oplus\mathcal{H}_I$ with
$\mathcal{H}_k=\mathbb{R}\phi_k$.  Under the calibrated clamp,
\begin{equation}
 \mathcal{A}^{(\lambda)}=
 \begin{pmatrix}-\lambda&0\\a_{Ik}&\mathcal{A}_{II}\end{pmatrix},
 \qquad
 \mathcal{D}^{(\lambda)}=
 \begin{pmatrix}\varepsilon_0/\lambda&0\\0&D_{II}\end{pmatrix}.
 \label{eq:exact_blocks}
\end{equation}
Let $\Sigma_\lambda$ be the stationary covariance solving
\[
  \mathcal{A}^{(\lambda)}\Sigma_\lambda
  +\Sigma_\lambda(\mathcal{A}^{(\lambda)})^*
  +\mathcal{D}^{(\lambda)}=0.
\]

\begin{lemma}[Block Lyapunov identities]
\label{lem:block_lyapunov_identities}
For every $\lambda>0$,
\begin{equation}
  \Sigma_{\lambda,kk}=\frac{\varepsilon_0}{2\lambda^2},
  \label{eq:Sigmakk_exact}
\end{equation}
and the target-to-bulk cross covariance is
\begin{equation}
  \Sigma_{\lambda,Ik}
  =\frac{\varepsilon_0}{2\lambda^2}
    (\lambda I-\mathcal{A}_{II})^{-1}a_{Ik}.
  \label{eq:SigmaIk_exact_vector}
\end{equation}
\end{lemma}

\begin{proof}
The $kk$ block of the Lyapunov equation is
\[
  -2\lambda\Sigma_{\lambda,kk}+\frac{\varepsilon_0}{\lambda}=0,
\]
which gives \eqref{eq:Sigmakk_exact}.  The $Ik$ block is
\[
  (\mathcal{A}_{II}-\lambda I)\Sigma_{\lambda,Ik}
  +a_{Ik}\Sigma_{\lambda,kk}=0.
\]
Substituting \eqref{eq:Sigmakk_exact} gives
\eqref{eq:SigmaIk_exact_vector}.
\end{proof}

\subsection{Cameron--Martin energy}

Let $Q_0\coloneqq\Sigma_{II}^{(0)}$ be the background bulk covariance and
$H_{Q_0}\coloneqq\Range(Q_0^{1/2})$ with energy norm
$|u|_{Q_0}=\|Q_0^{-1/2}u\|_{\mathcal{H}_I}$.  Assumption~\ref{assum:cm_stability}
states that $a_{Ik}\in H_{Q_0}$ and that the bulk resolvent is bounded on this
energy space:
\begin{equation}
  |(\lambda I-\mathcal{A}_{II})^{-1}u|_{Q_0}
  \le \frac{M}{\lambda+\omega}|u|_{Q_0}.
  \label{eq:cm_resolvent}
\end{equation}

\begin{lemma}[Finite-energy leakage]
\label{lem:finite_energy_leakage}
The cross covariance belongs to $H_{Q_0}$ and satisfies
\begin{equation}
 |\Sigma_{\lambda,Ik}|_{Q_0}
 \le
 \frac{\varepsilon_0M}{2\lambda^2(\lambda+\omega)}
 |a_{Ik}|_{Q_0}
 =O(\lambda^{-3}).
 \label{eq:cross_energy}
\end{equation}
\end{lemma}

\begin{proof}
Apply \eqref{eq:cm_resolvent} to
\eqref{eq:SigmaIk_exact_vector}.  This proves both the range statement and the
energy bound.
\end{proof}

\section{Anderson--Trapp Shorting and the Rank-One Residual}
\label{app:shorting_residual}

In finite dimensions, the Schur complement is only the bounded-inverse
shadow of Anderson--Trapp shorting.  The Hilbert space argument therefore
starts from shorting and recovers the Schur expression only when the
complementary block is boundedly invertible.

\begin{lemma}[Schur complement as the bounded-inverse shadow of shorting]
\label{lem:schur_shadow_shorting}
Let $\mathcal{E}$ and $\mathcal{F}$ be Hilbert spaces and let
\[
  T=\begin{pmatrix}A&B^*\\B&D\end{pmatrix}\ge0
\]
act on $\mathcal{E}\oplus\mathcal{F}$.  If $D$ is boundedly invertible on
$\mathcal{F}$, then the Anderson--Trapp short of $T$ to $\mathcal{E}$ has the
quadratic form
\begin{equation}
  \langle x,\mathcal{S}_{\mathcal{E}}(T)x\rangle
  =
  \langle x,Ax\rangle-\langle Bx,D^{-1}Bx\rangle.
  \label{eq:short_shadow}
\end{equation}
Thus $B^*D^{-1}B$ is exactly the variational shorting subtraction in the special
case where $D^{-1}$ is a bounded operator.
\end{lemma}

\begin{proof}
By the variational definition of the shorted operator,
\[
  \langle x,\mathcal{S}_{\mathcal{E}}(T)x\rangle
  =
  \inf_{y\in\mathcal{F}}
  \left\langle
  \begin{pmatrix}x\\y\end{pmatrix},
  T
  \begin{pmatrix}x\\y\end{pmatrix}
  \right\rangle .
\]
The displayed quadratic form equals
$\langle x,Ax\rangle+2\operatorname{Re}\langle Bx,y\rangle+\langle y,Dy\rangle$.
Since $D$ is boundedly invertible and positive, the unique minimizer is
$y=-D^{-1}Bx$.  Substitution gives \eqref{eq:short_shadow}.
\end{proof}

For the covariance block $\Sigma_\lambda$ on
$\mathbb{C}\phi_k\oplus\mathcal{H}_I$, the shorting subtraction is the
variational quantity
\begin{equation}
 R_\lambda
 =
 \sup_{\langle u,\Sigma_{\lambda,II}u\rangle>0}
 \frac{|\langle u,\Sigma_{\lambda,Ik}\rangle|^2}
      {\langle u,\Sigma_{\lambda,II}u\rangle}.
 \label{eq:schur_variational}
\end{equation}
The null directions of the denominator do not change the supremum for a
non-negative block operator.

\begin{proposition}[Rank-one shorted residual]
\label{prop:rank_one_shorted_residual}
Under Assumptions~\ref{assum:linear_second_moment} and
\ref{assum:cm_stability},
\begin{equation}
  S_{k,\lambda}
  =\Sigma_{\lambda,kk}-R_\lambda,
  \qquad
  0\le R_\lambda\le |\Sigma_{\lambda,Ik}|_{Q_0}^2=O(\lambda^{-6}).
  \label{eq:shorted_residual_bound}
\end{equation}
Consequently
\begin{equation}
  S_{k,\lambda}
  =
  \frac{\varepsilon_0}{2\lambda^2}-O(\lambda^{-6}),
  \qquad
  m_2^{\mathrm{Schur}}=\frac{\varepsilon_0}{2},
  \label{eq:Schur_asympt}
\end{equation}
and
\[
  S_{k,\lambda}^{-1}
  =
  \frac{2}{\varepsilon_0}\lambda^2+O(\lambda^{-2}).
\]
\end{proposition}

\begin{proof}
The first identity is the scalar form of Anderson--Trapp shorting applied to
the target coordinate.  Because the target noise and the free noise are
independent, the bulk covariance dominates the background covariance:
$\Sigma_{\lambda,II}=Q_0+\operatorname{Cov}(Y_I)\succeq Q_0$.  Douglas range
inclusion gives the contraction of energy norms
\[
  \|\Sigma_{\lambda,II}^{-1/2}v\|
  \le \|Q_0^{-1/2}v\|
  \quad (v\in H_{Q_0}).
\]
Using this inequality in \eqref{eq:schur_variational} and then applying
\eqref{eq:cross_energy} yields \eqref{eq:shorted_residual_bound}.  The
asymptotic formula follows from \eqref{eq:Sigmakk_exact}, and the reciprocal
expansion follows by expanding
$S_{k,\lambda}^{-1}$ around $\varepsilon_0/(2\lambda^2)$.
\end{proof}

\section{Galerkin Approximation and the Failure of Direct Inversion}
\label{app:galerkin_shorting}

Galerkin approximation is used here only to approximate trace-class covariances
and nested variational shorting problems.  It is not a license to replace the
Hilbert-space short by a global inverse of a compact covariance block.

\begin{lemma}[Trace-norm convergence of Galerkin covariances]
\label{lem:galerkin_trace_norm}
Let $\Pi_n:\mathcal{H}\to\mathcal{H}_n$ be finite-rank orthogonal projections
with $\Pi_n\phi_k=\phi_k$, $\Pi_n\mathcal{H}_I\subset\mathcal{H}_I$, and
$\Pi_n\to I$ strongly.  For fixed $\lambda>0$, set
\[
  \mathcal{A}_n^{(\lambda)}
  =\Pi_n\mathcal{A}^{(\lambda)}|_{\mathcal{H}_n},
  \qquad
  \mathcal{D}_n^{(\lambda)}
  =\Pi_n\mathcal{D}^{(\lambda)}\Pi_n .
\]
Assume that $e^{t\mathcal{A}_n^{(\lambda)}}x\to
e^{t\mathcal{A}^{(\lambda)}}x$ for every $x$, uniformly for $t$ in compact
intervals, that the semigroups are uniformly exponentially stable, satisfying
\[
  \|e^{t\mathcal{A}_n^{(\lambda)}}\|, \
  \|e^{t\mathcal{A}^{(\lambda)}}\|\le Me^{-\beta t},
\]
and that $\mathcal{D}_n^{(\lambda)}\to\mathcal{D}^{(\lambda)}$ in
$\mathcal{S}_1(\mathcal{H})$.  Then
\[
  \Sigma_{n,\lambda}
  =
  \int_0^\infty
  e^{t\mathcal{A}_n^{(\lambda)}}
  \mathcal{D}_n^{(\lambda)}
  e^{t\mathcal{A}_n^{(\lambda)*}}\,dt
  \longrightarrow
  \Sigma_\lambda
  \quad\text{in }\mathcal{S}_1(\mathcal{H}).
\]
\end{lemma}

\begin{proof}
Omitting the superscript $(\lambda)$, write the difference as a Bochner
integral and split the integrand:
\begin{multline*}
  \|e^{t\mathcal{A}_n}\mathcal{D}_ne^{t\mathcal{A}_n^*}
    -e^{t\mathcal{A}}\mathcal{D}e^{t\mathcal{A}^*}\|_{\mathcal{S}_1}
  \le
  M^2e^{-2\beta t}\|\mathcal{D}_n-\mathcal{D}\|_{\mathcal{S}_1}\\
  +
  \|e^{t\mathcal{A}_n}\mathcal{D}e^{t\mathcal{A}_n^*}
    -e^{t\mathcal{A}}\mathcal{D}e^{t\mathcal{A}^*}\|_{\mathcal{S}_1}.
\end{multline*}
The first term tends to zero.  For the second, Gr\"umm's theorem for trace
ideals implies trace-norm convergence from strong convergence of the left and
right semigroups, uniform boundedness, and $\mathcal{D}\in\mathcal{S}_1$
\cite{simon2005trace}.  The full integrand is dominated by
\[
  M^2e^{-2\beta t}
  \bigl(\|\mathcal{D}_n\|_{\mathcal{S}_1}
        +\|\mathcal{D}\|_{\mathcal{S}_1}\bigr),
\]
with a finite uniform trace-norm bound because
$\mathcal{D}_n\to\mathcal{D}$ in $\mathcal{S}_1$.  Dominated convergence for
Bochner integrals gives the claim.
\end{proof}

\begin{lemma}[Galerkin convergence of variational shorting]
\label{lem:galerkin_variational_shorting}
For every block-preserving Galerkin truncation satisfying the hypotheses of
\Cref{lem:galerkin_trace_norm} and the uniform finite-energy condition
\[
  a_{Ik}^{(n)}\in H_{Q_0^{(n)}},
  \qquad
  \sup_n |a_{Ik}^{(n)}|_{Q_0^{(n)}}<\infty,
\]
the truncated residual satisfies
\[
  S_{k,\lambda}^{(n)}
  =
  \frac{\varepsilon_0}{2\lambda^2}-R_{n,\lambda},
  \qquad
  0\le R_{n,\lambda}\le C\lambda^{-6},
\]
with $C$ independent of $n$.  Hence
$m_2^{\mathrm{Schur},(n)}=\varepsilon_0/2$ for every admissible truncation.
Moreover, for fixed $\lambda$, nested finite-dimensional test spaces in the
bulk Cameron--Martin space give monotone Rayleigh--Ritz convergence of the
variational subtraction to $R_\lambda$.
\end{lemma}

\begin{proof}
The block-preserving truncation has the same triangular form as
\eqref{eq:exact_blocks}, so the finite-dimensional Lyapunov blocks give
\[
  \Sigma_{n,\lambda,kk}=\frac{\varepsilon_0}{2\lambda^2},
  \qquad
  \Sigma_{n,\lambda,Ik}
  =
  \frac{\varepsilon_0}{2\lambda^2}
  (\lambda I-\mathcal{A}_{II}^{(n)})^{-1}a_{Ik}^{(n)}.
\]
The uniform finite-energy assumption gives
$|\Sigma_{n,\lambda,Ik}|_{Q_0^{(n)}}=O(\lambda^{-3})$ uniformly in $n$.
The same Douglas domination argument used in
\Cref{prop:rank_one_shorted_residual} then yields
$0\le R_{n,\lambda}\le C\lambda^{-6}$ and the stated value of
$m_2^{\mathrm{Schur},(n)}$.

For the variational convergence statement, let
$C_\lambda=\Sigma_{\lambda,II}$ and let $V_m$ be nested finite-dimensional
subspaces whose union is dense in $H_{C_\lambda}=\Range(C_\lambda^{1/2})$.
Then
\[
  R_\lambda^{[m]}
  =
  \sup_{\substack{u\in V_m\\ \langle u,C_\lambda u\rangle>0}}
  \frac{|\langle u,\Sigma_{\lambda,Ik}\rangle|^2}
       {\langle u,C_\lambda u\rangle}
\]
is monotone increasing in $m$ and bounded above by $R_\lambda$.  Density in the
Cameron--Martin energy norm gives $R_\lambda^{[m]}\uparrow R_\lambda$.  In an
eigenbasis of $C_\lambda$, this is the Rayleigh--Ritz identity
\[
  R_\lambda^{[m]}
  =
  \sum_{j=1}^m
  \frac{|\langle\Sigma_{\lambda,Ik},e_j\rangle|^2}{\sigma_j}
  \uparrow
  \|\Sigma_{\lambda,Ik}\|_{H_{C_\lambda}}^2.
\]
\end{proof}

The obstruction to direct inversion is not numerical bookkeeping.  In infinite
dimension the compact covariance block has no bounded inverse on the ambient
Hilbert space.  In finite-dimensional Galerkin matrices the condition number
$\kappa(D_N)$ diverges as the smallest covariance eigenvalue tends to zero, so
the formula $B_N^*D_N^{-1}B_N$ amplifies truncation and roundoff errors.  Even
with exact arithmetic, the limits do not commute: fixing $N$ and then sending
$\lambda\to\infty$ resolves a different finite matrix problem than first
passing to the trace-class covariance and then taking the calibrated shorting
limit.  The stable object is therefore the nested variational short, not direct
matrix inversion.

\section{Gaussian Radon Conditioning and the Fredholm Finite Part}
\label{app:gaussian_fredholm}

The Gaussian and Fredholm layer is the Radon representation of the shorted
residual constructed above.  It is not a second mechanism: disintegration
supplies the exact-evidence law, while the shorted covariance residual supplies
the finite conditional variance used by the Gaussian entropy calculation.

\subsection{Pointwise disintegration}

Let $\mu_\lambda=\mathcal{N}(m_\lambda,\Sigma_\lambda)$ be a Radon Gaussian
measure on $\mathcal{H}=\mathcal{H}_k\oplus\mathcal{H}_I$, and let
$\pi_k(X)=\langle X,\phi_k\rangle$.

\begin{lemma}[Continuous disintegration and exact target zeros]
\label{lem:disintegration}
For each $t\in\mathbb{R}$, there is a unique weakly continuous Gaussian
disintegration member
$\mu_\lambda^t=\mathcal{N}(m_\lambda^t,\Sigma_\lambda^t)$ supported on
$X_k=t$, where
\[
  m_\lambda^t=
  \begin{pmatrix}
    t\\
    m_{\lambda,I}
    +\Sigma_{\lambda,Ik}\Sigma_{\lambda,kk}^{-1}(t-m_{\lambda,k})
  \end{pmatrix},
\]
and
\[
  \Sigma_\lambda^t=
  \begin{pmatrix}
    0&0\\
    0&\Sigma_{\lambda,II}
      -\Sigma_{\lambda,Ik}\Sigma_{\lambda,kk}^{-1}\Sigma_{\lambda,kI}
  \end{pmatrix}.
\]
At $t=x_k^*$, the observed law $\mu_{\mathrm{obs}}=\mu_\lambda^{x_k^*}$
satisfies
\[
  \mathbb{E}_{\mathrm{obs}}[X_k]-x_k^*=0,
  \qquad
  \operatorname{Var}_{\mathrm{obs}}(X_k)=0.
\]
\end{lemma}

\begin{proof}
Regular conditional probabilities exist because $\mathcal{H}$ is Polish.  The
displayed family is the weakly continuous Gaussian disintegration along the
continuous scalar observation map, and this continuous version is unique
\cite{bogachev1998gaussian,lagatta2013continuous}.  Evaluating at $x_k^*$ fixes
the target coordinate almost surely.
\end{proof}

\subsection{Gaussian energy projection and entropy}

Let $C_\lambda=\Sigma_{\lambda,II}$.  Since $C_\lambda$ is trace-class, its
inverse is not a bounded ambient operator; the relevant object is the
Cameron--Martin space $H_{C_\lambda}=\Range(C_\lambda^{1/2})$.

\begin{lemma}[Measurable Cameron--Martin energy projection]
\label{lem:conditional_predictor}
The cross covariance belongs to $H_{C_\lambda}$ and
\[
  \|\Sigma_{\lambda,Ik}\|_{H_{C_\lambda}}
  \le |\Sigma_{\lambda,Ik}|_{Q_0}=O(\lambda^{-3}).
\]
The conditional mean projection is the measurable Wiener functional
\[
  \mu_{k|I}(X_I)
  =
  m_{\lambda,k}
  +W_{\Sigma_{\lambda,Ik}}(X_I-m_{\lambda,I}),
\]
and, for all sufficiently large $\lambda$,
\[
  \operatorname{Var}(X_k\mid X_I)
  =
  \Sigma_{\lambda,kk}
  -\|\Sigma_{\lambda,Ik}\|_{H_{C_\lambda}}^2
  =S_{k,\lambda}>0.
\]
\end{lemma}

\begin{proof}
The domination $\Sigma_{\lambda,II}\succeq Q_0$ gives
$H_{Q_0}\subseteq H_{C_\lambda}$ by Douglas range inclusion
\cite{douglas1966majorization}.  Combining this inclusion with
\eqref{eq:cross_energy} gives the norm bound.  The range condition is exactly
the condition for the measurable Wiener functional associated with
$\Sigma_{\lambda,Ik}$ to exist under the bulk Gaussian law
\cite{bogachev1998gaussian}.  Positivity follows from
\Cref{prop:rank_one_shorted_residual}.
\end{proof}

\begin{lemma}[Feldman--Hajek equivalence and trace-anomaly decay]
\label{lem:feldman_hajek}
There is $\lambda_0>0$ such that for $\lambda\ge\lambda_0$ the bulk marginals
$\mu_{\mathrm{obs},I}$ and $\mu_{\lambda,I}$ are equivalent and
\[
  2\KL(\mu_{\mathrm{obs},I}\parallel\mu_{\lambda,I})
  =
  \frac{(x_k^*-m_{\lambda,k})^2}{\Sigma_{\lambda,kk}}\rho_\lambda
  -\log(1-\rho_\lambda)-\rho_\lambda,
\]
where
\[
  \rho_\lambda
  =
  \frac{\|\Sigma_{\lambda,Ik}\|_{H_{C_\lambda}}^2}
       {\Sigma_{\lambda,kk}}
  =
  1-\frac{S_{k,\lambda}}{\Sigma_{\lambda,kk}}
  =O(\lambda^{-4}).
\]
In particular,
$\KL(\mu_{\mathrm{obs},I}\parallel\mu_{\lambda,I})=O(\lambda^{-4})$.
\end{lemma}

\begin{proof}
The covariance change under disintegration is the rank-one operator
$\Sigma_{\lambda,Ik}\Sigma_{\lambda,kk}^{-1}\Sigma_{\lambda,kI}$.  By
\Cref{lem:conditional_predictor}, its relative covariance perturbation has the
single non-unit eigenvalue $1-\rho_\lambda$.  For large $\lambda$,
$\rho_\lambda<1$, so Feldman--Hajek equivalence applies
\cite{feldman1958equivalence,bogachev1998gaussian}.  The Gaussian relative
entropy formula gives the displayed expression.  Assumption~\ref{assum:gaussian_clamp}
has $m_{\lambda,k}-x_k^*=O(\lambda^{-1})$, while
$\Sigma_{\lambda,kk}=O(\lambda^{-2})$ and $\rho_\lambda=O(\lambda^{-4})$; the
logarithmic remainder is $O(\rho_\lambda^2)$.  Hence the entropy is
$O(\lambda^{-4})$.
\end{proof}

\subsection{Mollified divergence and finite part}

For $\delta>0$, set
\[
  \nu_{\lambda,\delta}
  =
  \mu_{\mathrm{obs},I}\otimes\mathcal{N}(x_k^*,\delta^2)
\]
under the product identification
$\mathcal{H}=\mathcal{H}_I\times\mathcal{H}_k$.

\begin{proposition}[Mollification limit]
\label{prop:mollification_limit}
For $\lambda\ge\lambda_0$ as in \Cref{lem:feldman_hajek}, define
\[
  C_{\lambda,\delta}
  =
  \frac12\left(\log\frac{S_{k,\lambda}}{\delta^2}-1\right).
\]
Then
\[
  \lim_{\delta\downarrow0}
  \left[
    \KL(\nu_{\lambda,\delta}\parallel\mu_\lambda)
    -C_{\lambda,\delta}
  \right]
  =
  \mathfrak{J}_\lambda(\mu_{\mathrm{obs}};\mu_\lambda).
\]
\end{proposition}

\begin{proof}
Let $r_\lambda(X)=X_k-\mu_{k|I}(X_I)$.  The entropy chain rule and the scalar
Gaussian formula give
\[
  \KL(\nu_{\lambda,\delta}\parallel\mu_\lambda)
  =
  \KL(\mu_{\mathrm{obs},I}\parallel\mu_{\lambda,I})
  +C_{\lambda,\delta}
  +\frac{\delta^2}{2S_{k,\lambda}}
  +\frac{1}{2S_{k,\lambda}}
   \mathbb{E}_{\mu_{\mathrm{obs}}}[r_\lambda(X)^2].
    \]
    The $\delta^2$ term vanishes as $\delta\downarrow0$, and the remaining terms
    are exactly Definition~\ref{def:unified_action}.
\end{proof}

\begin{lemma}[Weak scheme-independence]
\label{lem:weak_scheme_independence}
For an evidence level $t\in\mathbb{R}$, let $\mu_\lambda^t$ be the weakly continuous
Gaussian disintegration from \Cref{lem:disintegration} and define
\[
\mathfrak{J}_\lambda(t)\coloneqq
\mathfrak{J}_\lambda(\mu_\lambda^t;\mu_\lambda).
\]
For any two evidence levels $t_1,t_2$ for which the finite part is defined,
\begin{equation}
    \mathfrak{J}_\lambda(t_2)-\mathfrak{J}_\lambda(t_1)
    =
    \frac{(t_2-m_{\lambda,k})^2-(t_1-m_{\lambda,k})^2}
    {2\Sigma_{\lambda,kk}}.
\end{equation}
\begin{proof}
Let $\Delta_t\coloneqq t-m_{\lambda,k}$ and
\[
    \rho_\lambda
    \coloneqq
    \frac{\|\Sigma_{\lambda,Ik}\|_{H_{C_\lambda}}^2}{\Sigma_{\lambda,kk}}
    =
    1-\frac{S_{k,\lambda}}{\Sigma_{\lambda,kk}}.
\]
The Gaussian entropy computation in \Cref{lem:feldman_hajek} gives
\[
    \KL(\mu_{\lambda,I}^t\parallel\mu_{\lambda,I})
    =
    \frac{1}{2}\frac{\Delta_t^2}{\Sigma_{\lambda,kk}}\rho_\lambda
    -\frac{1}{2}\log(1-\rho_\lambda)-\frac{1}{2}\rho_\lambda .
\]
The conditional-residual term in \eqref{eq:unified_action} is
\[
    \frac{1}{2S_{k,\lambda}}
    \E_{\mu_\lambda^t}\!\left[
    (X_k-\mu_{k|I}(X_I))^2
    \right]
    =
    \frac{1}{2}\rho_\lambda
    +
    \frac{1}{2}\frac{\Delta_t^2}{\Sigma_{\lambda,kk}}(1-\rho_\lambda).
\]
Adding these two identities cancels the $\rho_\lambda/2$ terms and yields
\[
    \mathfrak{J}_\lambda(t)
    =
    \frac{(t-m_{\lambda,k})^2}{2\Sigma_{\lambda,kk}}
    -\frac{1}{2}\log(1-\rho_\lambda).
\]
The logarithmic term depends on $C_\lambda$ and $\lambda$, but not on $t$. Hence it
 cancels in the difference between $t_2$ and $t_1$, proving
 \eqref{eq:weak_scheme_independence}. Replacing the Gaussian Dirac mollifier in
 \Cref{prop:mollification_limit} by any centered approximate identity with finite
 second moment changes the extracted finite part only by an additive term depending on
 the kernel, $\lambda$, and $S_{k,\lambda}$, not on $t$. Therefore all evidence
 differences $\Delta\mathfrak{J}_\lambda$ are kernel-independent.
\end{proof}
\end{lemma}

\begin{proof}[Proof of Theorem~\ref{thm:finite_part_regularization}]
Put $\Delta_\lambda=x_k^*-m_{\lambda,k}$ and
$Z_\lambda=W_{\Sigma_{\lambda,Ik}}(X_I-m_{\lambda,I})$.  The Gaussian
conditional formulas give
\[
  \mathbb{E}_{\mathrm{obs}}[Z_\lambda]
  =\rho_\lambda\Delta_\lambda,
  \qquad
  \operatorname{Var}_{\mathrm{obs}}(Z_\lambda)
  =
  \Sigma_{\lambda,kk}\rho_\lambda(1-\rho_\lambda).
\]
Since $X_k=x_k^*$ under $\mu_{\mathrm{obs}}$ and
$S_{k,\lambda}=\Sigma_{\lambda,kk}(1-\rho_\lambda)$,
\[
  \frac{1}{2S_{k,\lambda}}
  \mathbb{E}_{\mu_{\mathrm{obs}}}
  \bigl[(X_k-\mu_{k|I}(X_I))^2\bigr]
  =
  \frac{\rho_\lambda}{2}
  +
  \frac{\Delta_\lambda^2}{2\Sigma_{\lambda,kk}}(1-\rho_\lambda).
\]
The calibrated identities yield the relation
$\Delta_\lambda^2 / (2\Sigma_{\lambda,kk}) = b_k^2 / \varepsilon_0$
as well as $\rho_\lambda = O(\lambda^{-4})$.  Adding
\Cref{lem:feldman_hajek} yields
\[
  \mathfrak{J}_\lambda(\mu_{\mathrm{obs}};\mu_\lambda)
  =
  \frac{b_k^2}{\varepsilon_0}+O(\lambda^{-4}).
\]
\end{proof}

If $a_{Ik}=0$ and $x_k^*=m_{\lambda,k}$, then
$\Sigma_{\lambda,Ik}=0$, the bulk marginals coincide, and the conditional
residual term in Definition~\ref{def:unified_action} is zero.  Hence
$\mathfrak{J}_\lambda(\mu_{\mathrm{obs}};\mu_\lambda)=0$ exactly for every
$\lambda>0$.

\section{Admissibility Boundary and Elliptic Benchmark}
\label{app:admissibility_boundary}

The finite-energy assumption is a boundary condition for the theory, not a
technical artifact.  The off-diagonal channel must land in the Cameron--Martin
range of the background covariance.  If it does not, the variational subtraction
need not be finite and the $O(\lambda^{-6})$ rank-one residual estimate has no
meaning.

The elliptic example of \Cref{sec:spde} makes this boundary concrete.  For
$L=-d^2/dx^2+\kappa^2$ on $L^2(0,1)$ with Dirichlet boundary conditions and
$Q_0=L^{-1}$, the Cameron--Martin space is
\[
  H_{Q_0}=H_0^1(0,1),
  \qquad
  \|h\|_{Q_0}^2
  =
  \int_0^1\bigl(|h'(x)|^2+\kappa^2|h(x)|^2\bigr)\,dx .
\]
Smooth spatially averaged observations are admissible when their representing
vectors lie in $H_0^1(0,1)$; for example, any fixed
$g\in C_c^\infty(0,1)$ has finite energy.  A point-evaluation channel is the
opposite limit.  If $g_\varepsilon$ is a standard compactly supported mollifier
concentrating at $x_0\in(0,1)$, then
\[
  \|g_\varepsilon\|_{Q_0}^2
  =
  \int_0^1\bigl(|g_\varepsilon'(x)|^2+\kappa^2|g_\varepsilon(x)|^2\bigr)\,dx
  \longrightarrow \infty
  \qquad(\varepsilon\downarrow0).
\]
Thus each fixed smoothed sensor can satisfy the finite-energy condition, but
the Dirac point-evaluation limit does not.  The shorting theorem is therefore a
theory of admissible Cameron--Martin channels, not of unsmoothed point sensors
outside the covariance-energy domain.

\end{document}


% ============================================================
%  APPENDIX
% ============================================================
\appendix

\section{Block Lyapunov Identities for the Calibrated Clamp}
\label{app:block_identities}

This appendix isolates the algebra used in the paper's shorting argument.  The
only inputs are the block-triangular calibrated operator family, the trace-class
stationary covariance, and the Cameron--Martin finite-energy assumption on the
outgoing drift channel.

\subsection{Block identities}

Write $\mathcal{H}=\mathcal{H}_k\oplus\mathcal{H}_I$ with
$\mathcal{H}_k=\mathbb{R}\phi_k$.  Under the calibrated clamp,
\begin{equation}
 \mathcal{A}^{(\lambda)}=
 \begin{pmatrix}-\lambda&0\\a_{Ik}&\mathcal{A}_{II}\end{pmatrix},
 \qquad
 \mathcal{D}^{(\lambda)}=
 \begin{pmatrix}\varepsilon_0/\lambda&0\\0&D_{II}\end{pmatrix}.
 \label{eq:exact_blocks}
\end{equation}
Let $\Sigma_\lambda$ be the stationary covariance solving
\[
  \mathcal{A}^{(\lambda)}\Sigma_\lambda
  +\Sigma_\lambda(\mathcal{A}^{(\lambda)})^*
  +\mathcal{D}^{(\lambda)}=0.
\]

\begin{lemma}[Block Lyapunov identities]
\label{lem:block_lyapunov_identities}
For every $\lambda>0$,
\begin{equation}
  \Sigma_{\lambda,kk}=\frac{\varepsilon_0}{2\lambda^2},
  \label{eq:Sigmakk_exact}
\end{equation}
and the target-to-bulk cross covariance is
\begin{equation}
  \Sigma_{\lambda,Ik}
  =\frac{\varepsilon_0}{2\lambda^2}
    (\lambda I-\mathcal{A}_{II})^{-1}a_{Ik}.
  \label{eq:SigmaIk_exact_vector}
\end{equation}
\end{lemma}

\begin{proof}
The $kk$ block of the Lyapunov equation is
\[
  -2\lambda\Sigma_{\lambda,kk}+\frac{\varepsilon_0}{\lambda}=0,
\]
which gives \eqref{eq:Sigmakk_exact}.  The $Ik$ block is
\[
  (\mathcal{A}_{II}-\lambda I)\Sigma_{\lambda,Ik}
  +a_{Ik}\Sigma_{\lambda,kk}=0.
\]
Substituting \eqref{eq:Sigmakk_exact} gives
\eqref{eq:SigmaIk_exact_vector}.
\end{proof}

\subsection{Cameron--Martin energy}

Let $Q_0\coloneqq\Sigma_{II}^{(0)}$ be the background bulk covariance and
$H_{Q_0}\coloneqq\Range(Q_0^{1/2})$ with energy norm
$|u|_{Q_0}=\|Q_0^{-1/2}u\|_{\mathcal{H}_I}$.  Assumption~\ref{assum:cm_stability}
states that $a_{Ik}\in H_{Q_0}$ and that the bulk resolvent is bounded on this
energy space:
\begin{equation}
  |(\lambda I-\mathcal{A}_{II})^{-1}u|_{Q_0}
  \le \frac{M}{\lambda+\omega}|u|_{Q_0}.
  \label{eq:cm_resolvent}
\end{equation}

\begin{lemma}[Finite-energy leakage]
\label{lem:finite_energy_leakage}
The cross covariance belongs to $H_{Q_0}$ and satisfies
\begin{equation}
 |\Sigma_{\lambda,Ik}|_{Q_0}
 \le
 \frac{\varepsilon_0M}{2\lambda^2(\lambda+\omega)}
 |a_{Ik}|_{Q_0}
 =O(\lambda^{-3}).
 \label{eq:cross_energy}
\end{equation}
\end{lemma}

\begin{proof}
Apply \eqref{eq:cm_resolvent} to
\eqref{eq:SigmaIk_exact_vector}.  This proves both the range statement and the
energy bound.
\end{proof}

\section{Anderson--Trapp Shorting and the Rank-One Residual}
\label{app:shorting_residual}

In finite dimensions, the Schur complement is only the bounded-inverse
shadow of Anderson--Trapp shorting.  The Hilbert space argument therefore
starts from shorting and recovers the Schur expression only when the
complementary block is boundedly invertible.

\begin{lemma}[Schur complement as the bounded-inverse shadow of shorting]
\label{lem:schur_shadow_shorting}
Let $\mathcal{E}$ and $\mathcal{F}$ be Hilbert spaces and let
\[
  T=\begin{pmatrix}A&B^*\\B&D\end{pmatrix}\ge0
\]
act on $\mathcal{E}\oplus\mathcal{F}$.  If $D$ is boundedly invertible on
$\mathcal{F}$, then the Anderson--Trapp short of $T$ to $\mathcal{E}$ has the
quadratic form
\begin{equation}
  \langle x,\mathcal{S}_{\mathcal{E}}(T)x\rangle
  =
  \langle x,Ax\rangle-\langle Bx,D^{-1}Bx\rangle.
  \label{eq:short_shadow}
\end{equation}
Thus $B^*D^{-1}B$ is exactly the variational shorting subtraction in the special
case where $D^{-1}$ is a bounded operator.
\end{lemma}

\begin{proof}
By the variational definition of the shorted operator,
\[
  \langle x,\mathcal{S}_{\mathcal{E}}(T)x\rangle
  =
  \inf_{y\in\mathcal{F}}
  \left\langle
  \begin{pmatrix}x\\y\end{pmatrix},
  T
  \begin{pmatrix}x\\y\end{pmatrix}
  \right\rangle .
\]
The displayed quadratic form equals
$\langle x,Ax\rangle+2\operatorname{Re}\langle Bx,y\rangle+\langle y,Dy\rangle$.
Since $D$ is boundedly invertible and positive, the unique minimizer is
$y=-D^{-1}Bx$.  Substitution gives \eqref{eq:short_shadow}.
\end{proof}

For the covariance block $\Sigma_\lambda$ on
$\mathbb{C}\phi_k\oplus\mathcal{H}_I$, the shorting subtraction is the
variational quantity
\begin{equation}
 R_\lambda
 =
 \sup_{\langle u,\Sigma_{\lambda,II}u\rangle>0}
 \frac{|\langle u,\Sigma_{\lambda,Ik}\rangle|^2}
      {\langle u,\Sigma_{\lambda,II}u\rangle}.
 \label{eq:schur_variational}
\end{equation}
The null directions of the denominator do not change the supremum for a
non-negative block operator.

\begin{proposition}[Rank-one shorted residual]
\label{prop:rank_one_shorted_residual}
Under Assumptions~\ref{assum:linear_second_moment} and
\ref{assum:cm_stability},
\begin{equation}
  S_{k,\lambda}
  =\Sigma_{\lambda,kk}-R_\lambda,
  \qquad
  0\le R_\lambda\le |\Sigma_{\lambda,Ik}|_{Q_0}^2=O(\lambda^{-6}).
  \label{eq:shorted_residual_bound}
\end{equation}
Consequently
\begin{equation}
  S_{k,\lambda}
  =
  \frac{\varepsilon_0}{2\lambda^2}-O(\lambda^{-6}),
  \qquad
  m_2^{\mathrm{Schur}}=\frac{\varepsilon_0}{2},
  \label{eq:Schur_asympt}
\end{equation}
and
\[
  S_{k,\lambda}^{-1}
  =
  \frac{2}{\varepsilon_0}\lambda^2+O(\lambda^{-2}).
\]
\end{proposition}

\begin{proof}
The first identity is the scalar form of Anderson--Trapp shorting applied to
the target coordinate.  Because the target noise and the free noise are
independent, the bulk covariance dominates the background covariance:
$\Sigma_{\lambda,II}=Q_0+\operatorname{Cov}(Y_I)\succeq Q_0$.  Douglas range
inclusion gives the contraction of energy norms
\[
  \|\Sigma_{\lambda,II}^{-1/2}v\|
  \le \|Q_0^{-1/2}v\|
  \quad (v\in H_{Q_0}).
\]
Using this inequality in \eqref{eq:schur_variational} and then applying
\eqref{eq:cross_energy} yields \eqref{eq:shorted_residual_bound}.  The
asymptotic formula follows from \eqref{eq:Sigmakk_exact}, and the reciprocal
expansion follows by expanding
$S_{k,\lambda}^{-1}$ around $\varepsilon_0/(2\lambda^2)$.
\end{proof}

\section{Galerkin Approximation and the Failure of Direct Inversion}
\label{app:galerkin_shorting}

Galerkin approximation is used here only to approximate trace-class covariances
and nested variational shorting problems.  It is not a license to replace the
Hilbert-space short by a global inverse of a compact covariance block.

\begin{lemma}[Trace-norm convergence of Galerkin covariances]
\label{lem:galerkin_trace_norm}
Let $\Pi_n:\mathcal{H}\to\mathcal{H}_n$ be finite-rank orthogonal projections
with $\Pi_n\phi_k=\phi_k$, $\Pi_n\mathcal{H}_I\subset\mathcal{H}_I$, and
$\Pi_n\to I$ strongly.  For fixed $\lambda>0$, set
\[
  \mathcal{A}_n^{(\lambda)}
  =\Pi_n\mathcal{A}^{(\lambda)}|_{\mathcal{H}_n},
  \qquad
  \mathcal{D}_n^{(\lambda)}
  =\Pi_n\mathcal{D}^{(\lambda)}\Pi_n .
\]
Assume that $e^{t\mathcal{A}_n^{(\lambda)}}x\to
e^{t\mathcal{A}^{(\lambda)}}x$ for every $x$, uniformly for $t$ in compact
intervals, that the semigroups are uniformly exponentially stable, satisfying
\[
  \|e^{t\mathcal{A}_n^{(\lambda)}}\|, \
  \|e^{t\mathcal{A}^{(\lambda)}}\|\le Me^{-\beta t},
\]
and that $\mathcal{D}_n^{(\lambda)}\to\mathcal{D}^{(\lambda)}$ in
$\mathcal{S}_1(\mathcal{H})$.  Then
\[
  \Sigma_{n,\lambda}
  =
  \int_0^\infty
  e^{t\mathcal{A}_n^{(\lambda)}}
  \mathcal{D}_n^{(\lambda)}
  e^{t\mathcal{A}_n^{(\lambda)*}}\,dt
  \longrightarrow
  \Sigma_\lambda
  \quad\text{in }\mathcal{S}_1(\mathcal{H}).
\]
\end{lemma}

\begin{proof}
Omitting the superscript $(\lambda)$, write the difference as a Bochner
integral and split the integrand:
\begin{multline*}
  \|e^{t\mathcal{A}_n}\mathcal{D}_ne^{t\mathcal{A}_n^*}
    -e^{t\mathcal{A}}\mathcal{D}e^{t\mathcal{A}^*}\|_{\mathcal{S}_1}
  \le
  M^2e^{-2\beta t}\|\mathcal{D}_n-\mathcal{D}\|_{\mathcal{S}_1}\\
  +
  \|e^{t\mathcal{A}_n}\mathcal{D}e^{t\mathcal{A}_n^*}
    -e^{t\mathcal{A}}\mathcal{D}e^{t\mathcal{A}^*}\|_{\mathcal{S}_1}.
\end{multline*}
The first term tends to zero.  For the second, Gr\"umm's theorem for trace
ideals implies trace-norm convergence from strong convergence of the left and
right semigroups, uniform boundedness, and $\mathcal{D}\in\mathcal{S}_1$
\cite{simon2005trace}.  The full integrand is dominated by
\[
  M^2e^{-2\beta t}
  \bigl(\|\mathcal{D}_n\|_{\mathcal{S}_1}
        +\|\mathcal{D}\|_{\mathcal{S}_1}\bigr),
\]
with a finite uniform trace-norm bound because
$\mathcal{D}_n\to\mathcal{D}$ in $\mathcal{S}_1$.  Dominated convergence for
Bochner integrals gives the claim.
\end{proof}

\begin{lemma}[Galerkin convergence of variational shorting]
\label{lem:galerkin_variational_shorting}
For every block-preserving Galerkin truncation satisfying the hypotheses of
\Cref{lem:galerkin_trace_norm} and the uniform finite-energy condition
\[
  a_{Ik}^{(n)}\in H_{Q_0^{(n)}},
  \qquad
  \sup_n |a_{Ik}^{(n)}|_{Q_0^{(n)}}<\infty,
\]
the truncated residual satisfies
\[
  S_{k,\lambda}^{(n)}
  =
  \frac{\varepsilon_0}{2\lambda^2}-R_{n,\lambda},
  \qquad
  0\le R_{n,\lambda}\le C\lambda^{-6},
\]
with $C$ independent of $n$.  Hence
$m_2^{\mathrm{Schur},(n)}=\varepsilon_0/2$ for every admissible truncation.
Moreover, for fixed $\lambda$, nested finite-dimensional test spaces in the
bulk Cameron--Martin space give monotone Rayleigh--Ritz convergence of the
variational subtraction to $R_\lambda$.
\end{lemma}

\begin{proof}
The block-preserving truncation has the same triangular form as
\eqref{eq:exact_blocks}, so the finite-dimensional Lyapunov blocks give
\[
  \Sigma_{n,\lambda,kk}=\frac{\varepsilon_0}{2\lambda^2},
  \qquad
  \Sigma_{n,\lambda,Ik}
  =
  \frac{\varepsilon_0}{2\lambda^2}
  (\lambda I-\mathcal{A}_{II}^{(n)})^{-1}a_{Ik}^{(n)}.
\]
The uniform finite-energy assumption gives
$|\Sigma_{n,\lambda,Ik}|_{Q_0^{(n)}}=O(\lambda^{-3})$ uniformly in $n$.
The same Douglas domination argument used in
\Cref{prop:rank_one_shorted_residual} then yields
$0\le R_{n,\lambda}\le C\lambda^{-6}$ and the stated value of
$m_2^{\mathrm{Schur},(n)}$.

For the variational convergence statement, let
$C_\lambda=\Sigma_{\lambda,II}$ and let $V_m$ be nested finite-dimensional
subspaces whose union is dense in $H_{C_\lambda}=\Range(C_\lambda^{1/2})$.
Then
\[
  R_\lambda^{[m]}
  =
  \sup_{\substack{u\in V_m\\ \langle u,C_\lambda u\rangle>0}}
  \frac{|\langle u,\Sigma_{\lambda,Ik}\rangle|^2}
       {\langle u,C_\lambda u\rangle}
\]
is monotone increasing in $m$ and bounded above by $R_\lambda$.  Density in the
Cameron--Martin energy norm gives $R_\lambda^{[m]}\uparrow R_\lambda$.  In an
eigenbasis of $C_\lambda$, this is the Rayleigh--Ritz identity
\[
  R_\lambda^{[m]}
  =
  \sum_{j=1}^m
  \frac{|\langle\Sigma_{\lambda,Ik},e_j\rangle|^2}{\sigma_j}
  \uparrow
  \|\Sigma_{\lambda,Ik}\|_{H_{C_\lambda}}^2.
\]
\end{proof}

The obstruction to direct inversion is not numerical bookkeeping.  In infinite
dimension the compact covariance block has no bounded inverse on the ambient
Hilbert space.  In finite-dimensional Galerkin matrices the condition number
$\kappa(D_N)$ diverges as the smallest covariance eigenvalue tends to zero, so
the formula $B_N^*D_N^{-1}B_N$ amplifies truncation and roundoff errors.  Even
with exact arithmetic, the limits do not commute: fixing $N$ and then sending
$\lambda\to\infty$ resolves a different finite matrix problem than first
passing to the trace-class covariance and then taking the calibrated shorting
limit.  The stable object is therefore the nested variational short, not direct
matrix inversion.

\section{Gaussian Radon Conditioning and the Fredholm Finite Part}
\label{app:gaussian_fredholm}

The Gaussian and Fredholm layer is the Radon representation of the shorted
residual constructed above.  It is not a second mechanism: disintegration
supplies the exact-evidence law, while the shorted covariance residual supplies
the finite conditional variance used by the Gaussian entropy calculation.

\subsection{Pointwise disintegration}

Let $\mu_\lambda=\mathcal{N}(m_\lambda,\Sigma_\lambda)$ be a Radon Gaussian
measure on $\mathcal{H}=\mathcal{H}_k\oplus\mathcal{H}_I$, and let
$\pi_k(X)=\langle X,\phi_k\rangle$.

\begin{lemma}[Continuous disintegration and exact target zeros]
\label{lem:disintegration}
For each $t\in\mathbb{R}$, there is a unique weakly continuous Gaussian
disintegration member
$\mu_\lambda^t=\mathcal{N}(m_\lambda^t,\Sigma_\lambda^t)$ supported on
$X_k=t$, where
\[
  m_\lambda^t=
  \begin{pmatrix}
    t\\
    m_{\lambda,I}
    +\Sigma_{\lambda,Ik}\Sigma_{\lambda,kk}^{-1}(t-m_{\lambda,k})
  \end{pmatrix},
\]
and
\[
  \Sigma_\lambda^t=
  \begin{pmatrix}
    0&0\\
    0&\Sigma_{\lambda,II}
      -\Sigma_{\lambda,Ik}\Sigma_{\lambda,kk}^{-1}\Sigma_{\lambda,kI}
  \end{pmatrix}.
\]
At $t=x_k^*$, the observed law $\mu_{\mathrm{obs}}=\mu_\lambda^{x_k^*}$
satisfies
\[
  \mathbb{E}_{\mathrm{obs}}[X_k]-x_k^*=0,
  \qquad
  \operatorname{Var}_{\mathrm{obs}}(X_k)=0.
\]
\end{lemma}

\begin{proof}
Regular conditional probabilities exist because $\mathcal{H}$ is Polish.  The
displayed family is the weakly continuous Gaussian disintegration along the
continuous scalar observation map, and this continuous version is unique
\cite{bogachev1998gaussian,lagatta2013continuous}.  Evaluating at $x_k^*$ fixes
the target coordinate almost surely.
\end{proof}

\subsection{Gaussian energy projection and entropy}

Let $C_\lambda=\Sigma_{\lambda,II}$.  Since $C_\lambda$ is trace-class, its
inverse is not a bounded ambient operator; the relevant object is the
Cameron--Martin space $H_{C_\lambda}=\Range(C_\lambda^{1/2})$.

\begin{lemma}[Measurable Cameron--Martin energy projection]
\label{lem:conditional_predictor}
The cross covariance belongs to $H_{C_\lambda}$ and
\[
  \|\Sigma_{\lambda,Ik}\|_{H_{C_\lambda}}
  \le |\Sigma_{\lambda,Ik}|_{Q_0}=O(\lambda^{-3}).
\]
The conditional mean projection is the measurable Wiener functional
\[
  \mu_{k|I}(X_I)
  =
  m_{\lambda,k}
  +W_{\Sigma_{\lambda,Ik}}(X_I-m_{\lambda,I}),
\]
and, for all sufficiently large $\lambda$,
\[
  \operatorname{Var}(X_k\mid X_I)
  =
  \Sigma_{\lambda,kk}
  -\|\Sigma_{\lambda,Ik}\|_{H_{C_\lambda}}^2
  =S_{k,\lambda}>0.
\]
\end{lemma}

\begin{proof}
The domination $\Sigma_{\lambda,II}\succeq Q_0$ gives
$H_{Q_0}\subseteq H_{C_\lambda}$ by Douglas range inclusion
\cite{douglas1966majorization}.  Combining this inclusion with
\eqref{eq:cross_energy} gives the norm bound.  The range condition is exactly
the condition for the measurable Wiener functional associated with
$\Sigma_{\lambda,Ik}$ to exist under the bulk Gaussian law
\cite{bogachev1998gaussian}.  Positivity follows from
\Cref{prop:rank_one_shorted_residual}.
\end{proof}

\begin{lemma}[Feldman--Hajek equivalence and trace-anomaly decay]
\label{lem:feldman_hajek}
There is $\lambda_0>0$ such that for $\lambda\ge\lambda_0$ the bulk marginals
$\mu_{\mathrm{obs},I}$ and $\mu_{\lambda,I}$ are equivalent and
\[
  2\KL(\mu_{\mathrm{obs},I}\parallel\mu_{\lambda,I})
  =
  \frac{(x_k^*-m_{\lambda,k})^2}{\Sigma_{\lambda,kk}}\rho_\lambda
  -\log(1-\rho_\lambda)-\rho_\lambda,
\]
where
\[
  \rho_\lambda
  =
  \frac{\|\Sigma_{\lambda,Ik}\|_{H_{C_\lambda}}^2}
       {\Sigma_{\lambda,kk}}
  =
  1-\frac{S_{k,\lambda}}{\Sigma_{\lambda,kk}}
  =O(\lambda^{-4}).
\]
In particular,
$\KL(\mu_{\mathrm{obs},I}\parallel\mu_{\lambda,I})=O(\lambda^{-4})$.
\end{lemma}

\begin{proof}
The covariance change under disintegration is the rank-one operator
$\Sigma_{\lambda,Ik}\Sigma_{\lambda,kk}^{-1}\Sigma_{\lambda,kI}$.  By
\Cref{lem:conditional_predictor}, its relative covariance perturbation has the
single non-unit eigenvalue $1-\rho_\lambda$.  For large $\lambda$,
$\rho_\lambda<1$, so Feldman--Hajek equivalence applies
\cite{feldman1958equivalence,bogachev1998gaussian}.  The Gaussian relative
entropy formula gives the displayed expression.  Assumption~\ref{assum:gaussian_clamp}
has $m_{\lambda,k}-x_k^*=O(\lambda^{-1})$, while
$\Sigma_{\lambda,kk}=O(\lambda^{-2})$ and $\rho_\lambda=O(\lambda^{-4})$; the
logarithmic remainder is $O(\rho_\lambda^2)$.  Hence the entropy is
$O(\lambda^{-4})$.
\end{proof}

\subsection{Mollified divergence and finite part}

For $\delta>0$, set
\[
  \nu_{\lambda,\delta}
  =
  \mu_{\mathrm{obs},I}\otimes\mathcal{N}(x_k^*,\delta^2)
\]
under the product identification
$\mathcal{H}=\mathcal{H}_I\times\mathcal{H}_k$.

\begin{proposition}[Mollification limit]
\label{prop:mollification_limit}
For $\lambda\ge\lambda_0$ as in \Cref{lem:feldman_hajek}, define
\[
  C_{\lambda,\delta}
  =
  \frac12\left(\log\frac{S_{k,\lambda}}{\delta^2}-1\right).
\]
Then
\[
  \lim_{\delta\downarrow0}
  \left[
    \KL(\nu_{\lambda,\delta}\parallel\mu_\lambda)
    -C_{\lambda,\delta}
  \right]
  =
  \mathfrak{J}_\lambda(\mu_{\mathrm{obs}};\mu_\lambda).
\]
\end{proposition}

\begin{proof}
Let $r_\lambda(X)=X_k-\mu_{k|I}(X_I)$.  The entropy chain rule and the scalar
Gaussian formula give
\[
  \KL(\nu_{\lambda,\delta}\parallel\mu_\lambda)
  =
  \KL(\mu_{\mathrm{obs},I}\parallel\mu_{\lambda,I})
  +C_{\lambda,\delta}
  +\frac{\delta^2}{2S_{k,\lambda}}
  +\frac{1}{2S_{k,\lambda}}
   \mathbb{E}_{\mu_{\mathrm{obs}}}[r_\lambda(X)^2].
    \]
    The $\delta^2$ term vanishes as $\delta\downarrow0$, and the remaining terms
    are exactly Definition~\ref{def:unified_action}.
\end{proof}

\begin{lemma}[Weak scheme-independence]
\label{lem:weak_scheme_independence}
For an evidence level $t\in\mathbb{R}$, let $\mu_\lambda^t$ be the weakly continuous
Gaussian disintegration from \Cref{lem:disintegration} and define
\[
\mathfrak{J}_\lambda(t)\coloneqq
\mathfrak{J}_\lambda(\mu_\lambda^t;\mu_\lambda).
\]
For any two evidence levels $t_1,t_2$ for which the finite part is defined,
\begin{equation}
    \mathfrak{J}_\lambda(t_2)-\mathfrak{J}_\lambda(t_1)
    =
    \frac{(t_2-m_{\lambda,k})^2-(t_1-m_{\lambda,k})^2}
    {2\Sigma_{\lambda,kk}}.
\end{equation}
\begin{proof}
Let $\Delta_t\coloneqq t-m_{\lambda,k}$ and
\[
    \rho_\lambda
    \coloneqq
    \frac{\|\Sigma_{\lambda,Ik}\|_{H_{C_\lambda}}^2}{\Sigma_{\lambda,kk}}
    =
    1-\frac{S_{k,\lambda}}{\Sigma_{\lambda,kk}}.
\]
The Gaussian entropy computation in \Cref{lem:feldman_hajek} gives
\[
    \KL(\mu_{\lambda,I}^t\parallel\mu_{\lambda,I})
    =
    \frac{1}{2}\frac{\Delta_t^2}{\Sigma_{\lambda,kk}}\rho_\lambda
    -\frac{1}{2}\log(1-\rho_\lambda)-\frac{1}{2}\rho_\lambda .
\]
The conditional-residual term in \eqref{eq:unified_action} is
\[
    \frac{1}{2S_{k,\lambda}}
    \E_{\mu_\lambda^t}\!\left[
    (X_k-\mu_{k|I}(X_I))^2
    \right]
    =
    \frac{1}{2}\rho_\lambda
    +
    \frac{1}{2}\frac{\Delta_t^2}{\Sigma_{\lambda,kk}}(1-\rho_\lambda).
\]
Adding these two identities cancels the $\rho_\lambda/2$ terms and yields
\[
    \mathfrak{J}_\lambda(t)
    =
    \frac{(t-m_{\lambda,k})^2}{2\Sigma_{\lambda,kk}}
    -\frac{1}{2}\log(1-\rho_\lambda).
\]
The logarithmic term depends on $C_\lambda$ and $\lambda$, but not on $t$. Hence it
 cancels in the difference between $t_2$ and $t_1$, proving
 \eqref{eq:weak_scheme_independence}. Replacing the Gaussian Dirac mollifier in
 \Cref{prop:mollification_limit} by any centered approximate identity with finite
 second moment changes the extracted finite part only by an additive term depending on
 the kernel, $\lambda$, and $S_{k,\lambda}$, not on $t$. Therefore all evidence
 differences $\Delta\mathfrak{J}_\lambda$ are kernel-independent.
\end{proof}
\end{lemma}

\begin{proof}[Proof of Theorem~\ref{thm:finite_part_regularization}]
Put $\Delta_\lambda=x_k^*-m_{\lambda,k}$ and
$Z_\lambda=W_{\Sigma_{\lambda,Ik}}(X_I-m_{\lambda,I})$.  The Gaussian
conditional formulas give
\[
  \mathbb{E}_{\mathrm{obs}}[Z_\lambda]
  =\rho_\lambda\Delta_\lambda,
  \qquad
  \operatorname{Var}_{\mathrm{obs}}(Z_\lambda)
  =
  \Sigma_{\lambda,kk}\rho_\lambda(1-\rho_\lambda).
\]
Since $X_k=x_k^*$ under $\mu_{\mathrm{obs}}$ and
$S_{k,\lambda}=\Sigma_{\lambda,kk}(1-\rho_\lambda)$,
\[
  \frac{1}{2S_{k,\lambda}}
  \mathbb{E}_{\mu_{\mathrm{obs}}}
  \bigl[(X_k-\mu_{k|I}(X_I))^2\bigr]
  =
  \frac{\rho_\lambda}{2}
  +
  \frac{\Delta_\lambda^2}{2\Sigma_{\lambda,kk}}(1-\rho_\lambda).
\]
The calibrated identities yield the relation
$\Delta_\lambda^2 / (2\Sigma_{\lambda,kk}) = b_k^2 / \varepsilon_0$
as well as $\rho_\lambda = O(\lambda^{-4})$.  Adding
\Cref{lem:feldman_hajek} yields
\[
  \mathfrak{J}_\lambda(\mu_{\mathrm{obs}};\mu_\lambda)
  =
  \frac{b_k^2}{\varepsilon_0}+O(\lambda^{-4}).
\]
\end{proof}

If $a_{Ik}=0$ and $x_k^*=m_{\lambda,k}$, then
$\Sigma_{\lambda,Ik}=0$, the bulk marginals coincide, and the conditional
residual term in Definition~\ref{def:unified_action} is zero.  Hence
$\mathfrak{J}_\lambda(\mu_{\mathrm{obs}};\mu_\lambda)=0$ exactly for every
$\lambda>0$.

\section{Admissibility Boundary and Elliptic Benchmark}
\label{app:admissibility_boundary}

The finite-energy assumption is a boundary condition for the theory, not a
technical artifact.  The off-diagonal channel must land in the Cameron--Martin
range of the background covariance.  If it does not, the variational subtraction
need not be finite and the $O(\lambda^{-6})$ rank-one residual estimate has no
meaning.

The elliptic example of \Cref{sec:spde} makes this boundary concrete.  For
$L=-d^2/dx^2+\kappa^2$ on $L^2(0,1)$ with Dirichlet boundary conditions and
$Q_0=L^{-1}$, the Cameron--Martin space is
\[
  H_{Q_0}=H_0^1(0,1),
  \qquad
  \|h\|_{Q_0}^2
  =
  \int_0^1\bigl(|h'(x)|^2+\kappa^2|h(x)|^2\bigr)\,dx .
\]
Smooth spatially averaged observations are admissible when their representing
vectors lie in $H_0^1(0,1)$; for example, any fixed
$g\in C_c^\infty(0,1)$ has finite energy.  A point-evaluation channel is the
opposite limit.  If $g_\varepsilon$ is a standard compactly supported mollifier
concentrating at $x_0\in(0,1)$, then
\[
  \|g_\varepsilon\|_{Q_0}^2
  =
  \int_0^1\bigl(|g_\varepsilon'(x)|^2+\kappa^2|g_\varepsilon(x)|^2\bigr)\,dx
  \longrightarrow \infty
  \qquad(\varepsilon\downarrow0).
\]
Thus each fixed smoothed sensor can satisfy the finite-energy condition, but
the Dirac point-evaluation limit does not.  The shorting theorem is therefore a
theory of admissible Cameron--Martin channels, not of unsmoothed point sensors
outside the covariance-energy domain.

\end{document}


% ============================================================
%  APPENDIX
% ============================================================
\appendix

\section{Block Lyapunov Identities for the Calibrated Clamp}
\label{app:block_identities}

This appendix isolates the algebra used in the paper's shorting argument.  The
only inputs are the block-triangular calibrated operator family, the trace-class
stationary covariance, and the Cameron--Martin finite-energy assumption on the
outgoing drift channel.

\subsection{Block identities}

Write $\mathcal{H}=\mathcal{H}_k\oplus\mathcal{H}_I$ with
$\mathcal{H}_k=\mathbb{R}\phi_k$.  Under the calibrated clamp,
\begin{equation}
 \mathcal{A}^{(\lambda)}=
 \begin{pmatrix}-\lambda&0\\a_{Ik}&\mathcal{A}_{II}\end{pmatrix},
 \qquad
 \mathcal{D}^{(\lambda)}=
 \begin{pmatrix}\varepsilon_0/\lambda&0\\0&D_{II}\end{pmatrix}.
 \label{eq:exact_blocks}
\end{equation}
Let $\Sigma_\lambda$ be the stationary covariance solving
\[
  \mathcal{A}^{(\lambda)}\Sigma_\lambda
  +\Sigma_\lambda(\mathcal{A}^{(\lambda)})^*
  +\mathcal{D}^{(\lambda)}=0.
\]

\begin{lemma}[Block Lyapunov identities]
\label{lem:block_lyapunov_identities}
For every $\lambda>0$,
\begin{equation}
  \Sigma_{\lambda,kk}=\frac{\varepsilon_0}{2\lambda^2},
  \label{eq:Sigmakk_exact}
\end{equation}
and the target-to-bulk cross covariance is
\begin{equation}
  \Sigma_{\lambda,Ik}
  =\frac{\varepsilon_0}{2\lambda^2}
    (\lambda I-\mathcal{A}_{II})^{-1}a_{Ik}.
  \label{eq:SigmaIk_exact_vector}
\end{equation}
\end{lemma}

\begin{proof}
The $kk$ block of the Lyapunov equation is
\[
  -2\lambda\Sigma_{\lambda,kk}+\frac{\varepsilon_0}{\lambda}=0,
\]
which gives \eqref{eq:Sigmakk_exact}.  The $Ik$ block is
\[
  (\mathcal{A}_{II}-\lambda I)\Sigma_{\lambda,Ik}
  +a_{Ik}\Sigma_{\lambda,kk}=0.
\]
Substituting \eqref{eq:Sigmakk_exact} gives
\eqref{eq:SigmaIk_exact_vector}.
\end{proof}

\subsection{Cameron--Martin energy}

Let $Q_0\coloneqq\Sigma_{II}^{(0)}$ be the background bulk covariance and
$H_{Q_0}\coloneqq\Range(Q_0^{1/2})$ with energy norm
$|u|_{Q_0}=\|Q_0^{-1/2}u\|_{\mathcal{H}_I}$.  Assumption~\ref{assum:cm_stability}
states that $a_{Ik}\in H_{Q_0}$ and that the bulk resolvent is bounded on this
energy space:
\begin{equation}
  |(\lambda I-\mathcal{A}_{II})^{-1}u|_{Q_0}
  \le \frac{M}{\lambda+\omega}|u|_{Q_0}.
  \label{eq:cm_resolvent}
\end{equation}

\begin{lemma}[Finite-energy leakage]
\label{lem:finite_energy_leakage}
The cross covariance belongs to $H_{Q_0}$ and satisfies
\begin{equation}
 |\Sigma_{\lambda,Ik}|_{Q_0}
 \le
 \frac{\varepsilon_0M}{2\lambda^2(\lambda+\omega)}
 |a_{Ik}|_{Q_0}
 =O(\lambda^{-3}).
 \label{eq:cross_energy}
\end{equation}
\end{lemma}

\begin{proof}
Apply \eqref{eq:cm_resolvent} to
\eqref{eq:SigmaIk_exact_vector}.  This proves both the range statement and the
energy bound.
\end{proof}

\section{Anderson--Trapp Shorting and the Rank-One Residual}
\label{app:shorting_residual}

In finite dimensions, the Schur complement is only the bounded-inverse
shadow of Anderson--Trapp shorting.  The Hilbert space argument therefore
starts from shorting and recovers the Schur expression only when the
complementary block is boundedly invertible.

\begin{lemma}[Schur complement as the bounded-inverse shadow of shorting]
\label{lem:schur_shadow_shorting}
Let $\mathcal{E}$ and $\mathcal{F}$ be Hilbert spaces and let
\[
  T=\begin{pmatrix}A&B^*\\B&D\end{pmatrix}\ge0
\]
act on $\mathcal{E}\oplus\mathcal{F}$.  If $D$ is boundedly invertible on
$\mathcal{F}$, then the Anderson--Trapp short of $T$ to $\mathcal{E}$ has the
quadratic form
\begin{equation}
  \langle x,\mathcal{S}_{\mathcal{E}}(T)x\rangle
  =
  \langle x,Ax\rangle-\langle Bx,D^{-1}Bx\rangle.
  \label{eq:short_shadow}
\end{equation}
Thus $B^*D^{-1}B$ is exactly the variational shorting subtraction in the special
case where $D^{-1}$ is a bounded operator.
\end{lemma}

\begin{proof}
By the variational definition of the shorted operator,
\[
  \langle x,\mathcal{S}_{\mathcal{E}}(T)x\rangle
  =
  \inf_{y\in\mathcal{F}}
  \left\langle
  \begin{pmatrix}x\\y\end{pmatrix},
  T
  \begin{pmatrix}x\\y\end{pmatrix}
  \right\rangle .
\]
The displayed quadratic form equals
$\langle x,Ax\rangle+2\operatorname{Re}\langle Bx,y\rangle+\langle y,Dy\rangle$.
Since $D$ is boundedly invertible and positive, the unique minimizer is
$y=-D^{-1}Bx$.  Substitution gives \eqref{eq:short_shadow}.
\end{proof}

For the covariance block $\Sigma_\lambda$ on
$\mathbb{C}\phi_k\oplus\mathcal{H}_I$, the shorting subtraction is the
variational quantity
\begin{equation}
 R_\lambda
 =
 \sup_{\langle u,\Sigma_{\lambda,II}u\rangle>0}
 \frac{|\langle u,\Sigma_{\lambda,Ik}\rangle|^2}
      {\langle u,\Sigma_{\lambda,II}u\rangle}.
 \label{eq:schur_variational}
\end{equation}
The null directions of the denominator do not change the supremum for a
non-negative block operator.

\begin{proposition}[Rank-one shorted residual]
\label{prop:rank_one_shorted_residual}
Under Assumptions~\ref{assum:linear_second_moment} and
\ref{assum:cm_stability},
\begin{equation}
  S_{k,\lambda}
  =\Sigma_{\lambda,kk}-R_\lambda,
  \qquad
  0\le R_\lambda\le |\Sigma_{\lambda,Ik}|_{Q_0}^2=O(\lambda^{-6}).
  \label{eq:shorted_residual_bound}
\end{equation}
Consequently
\begin{equation}
  S_{k,\lambda}
  =
  \frac{\varepsilon_0}{2\lambda^2}-O(\lambda^{-6}),
  \qquad
  m_2^{\mathrm{Schur}}=\frac{\varepsilon_0}{2},
  \label{eq:Schur_asympt}
\end{equation}
and
\[
  S_{k,\lambda}^{-1}
  =
  \frac{2}{\varepsilon_0}\lambda^2+O(\lambda^{-2}).
\]
\end{proposition}

\begin{proof}
The first identity is the scalar form of Anderson--Trapp shorting applied to
the target coordinate.  Because the target noise and the free noise are
independent, the bulk covariance dominates the background covariance:
$\Sigma_{\lambda,II}=Q_0+\operatorname{Cov}(Y_I)\succeq Q_0$.  Douglas range
inclusion gives the contraction of energy norms
\[
  \|\Sigma_{\lambda,II}^{-1/2}v\|
  \le \|Q_0^{-1/2}v\|
  \quad (v\in H_{Q_0}).
\]
Using this inequality in \eqref{eq:schur_variational} and then applying
\eqref{eq:cross_energy} yields \eqref{eq:shorted_residual_bound}.  The
asymptotic formula follows from \eqref{eq:Sigmakk_exact}, and the reciprocal
expansion follows by expanding
$S_{k,\lambda}^{-1}$ around $\varepsilon_0/(2\lambda^2)$.
\end{proof}

\section{Galerkin Approximation and the Failure of Direct Inversion}
\label{app:galerkin_shorting}

Galerkin approximation is used here only to approximate trace-class covariances
and nested variational shorting problems.  It is not a license to replace the
Hilbert-space short by a global inverse of a compact covariance block.

\begin{lemma}[Trace-norm convergence of Galerkin covariances]
\label{lem:galerkin_trace_norm}
Let $\Pi_n:\mathcal{H}\to\mathcal{H}_n$ be finite-rank orthogonal projections
with $\Pi_n\phi_k=\phi_k$, $\Pi_n\mathcal{H}_I\subset\mathcal{H}_I$, and
$\Pi_n\to I$ strongly.  For fixed $\lambda>0$, set
\[
  \mathcal{A}_n^{(\lambda)}
  =\Pi_n\mathcal{A}^{(\lambda)}|_{\mathcal{H}_n},
  \qquad
  \mathcal{D}_n^{(\lambda)}
  =\Pi_n\mathcal{D}^{(\lambda)}\Pi_n .
\]
Assume that $e^{t\mathcal{A}_n^{(\lambda)}}x\to
e^{t\mathcal{A}^{(\lambda)}}x$ for every $x$, uniformly for $t$ in compact
intervals, that the semigroups are uniformly exponentially stable, satisfying
\[
  \|e^{t\mathcal{A}_n^{(\lambda)}}\|, \
  \|e^{t\mathcal{A}^{(\lambda)}}\|\le Me^{-\beta t},
\]
and that $\mathcal{D}_n^{(\lambda)}\to\mathcal{D}^{(\lambda)}$ in
$\mathcal{S}_1(\mathcal{H})$.  Then
\[
  \Sigma_{n,\lambda}
  =
  \int_0^\infty
  e^{t\mathcal{A}_n^{(\lambda)}}
  \mathcal{D}_n^{(\lambda)}
  e^{t\mathcal{A}_n^{(\lambda)*}}\,dt
  \longrightarrow
  \Sigma_\lambda
  \quad\text{in }\mathcal{S}_1(\mathcal{H}).
\]
\end{lemma}

\begin{proof}
Omitting the superscript $(\lambda)$, write the difference as a Bochner
integral and split the integrand:
\begin{multline*}
  \|e^{t\mathcal{A}_n}\mathcal{D}_ne^{t\mathcal{A}_n^*}
    -e^{t\mathcal{A}}\mathcal{D}e^{t\mathcal{A}^*}\|_{\mathcal{S}_1}
  \le
  M^2e^{-2\beta t}\|\mathcal{D}_n-\mathcal{D}\|_{\mathcal{S}_1}\\
  +
  \|e^{t\mathcal{A}_n}\mathcal{D}e^{t\mathcal{A}_n^*}
    -e^{t\mathcal{A}}\mathcal{D}e^{t\mathcal{A}^*}\|_{\mathcal{S}_1}.
\end{multline*}
The first term tends to zero.  For the second, Gr\"umm's theorem for trace
ideals implies trace-norm convergence from strong convergence of the left and
right semigroups, uniform boundedness, and $\mathcal{D}\in\mathcal{S}_1$
\cite{simon2005trace}.  The full integrand is dominated by
\[
  M^2e^{-2\beta t}
  \bigl(\|\mathcal{D}_n\|_{\mathcal{S}_1}
        +\|\mathcal{D}\|_{\mathcal{S}_1}\bigr),
\]
with a finite uniform trace-norm bound because
$\mathcal{D}_n\to\mathcal{D}$ in $\mathcal{S}_1$.  Dominated convergence for
Bochner integrals gives the claim.
\end{proof}

\begin{lemma}[Galerkin convergence of variational shorting]
\label{lem:galerkin_variational_shorting}
For every block-preserving Galerkin truncation satisfying the hypotheses of
\Cref{lem:galerkin_trace_norm} and the uniform finite-energy condition
\[
  a_{Ik}^{(n)}\in H_{Q_0^{(n)}},
  \qquad
  \sup_n |a_{Ik}^{(n)}|_{Q_0^{(n)}}<\infty,
\]
the truncated residual satisfies
\[
  S_{k,\lambda}^{(n)}
  =
  \frac{\varepsilon_0}{2\lambda^2}-R_{n,\lambda},
  \qquad
  0\le R_{n,\lambda}\le C\lambda^{-6},
\]
with $C$ independent of $n$.  Hence
$m_2^{\mathrm{Schur},(n)}=\varepsilon_0/2$ for every admissible truncation.
Moreover, for fixed $\lambda$, nested finite-dimensional test spaces in the
bulk Cameron--Martin space give monotone Rayleigh--Ritz convergence of the
variational subtraction to $R_\lambda$.
\end{lemma}

\begin{proof}
The block-preserving truncation has the same triangular form as
\eqref{eq:exact_blocks}, so the finite-dimensional Lyapunov blocks give
\[
  \Sigma_{n,\lambda,kk}=\frac{\varepsilon_0}{2\lambda^2},
  \qquad
  \Sigma_{n,\lambda,Ik}
  =
  \frac{\varepsilon_0}{2\lambda^2}
  (\lambda I-\mathcal{A}_{II}^{(n)})^{-1}a_{Ik}^{(n)}.
\]
The uniform finite-energy assumption gives
$|\Sigma_{n,\lambda,Ik}|_{Q_0^{(n)}}=O(\lambda^{-3})$ uniformly in $n$.
The same Douglas domination argument used in
\Cref{prop:rank_one_shorted_residual} then yields
$0\le R_{n,\lambda}\le C\lambda^{-6}$ and the stated value of
$m_2^{\mathrm{Schur},(n)}$.

For the variational convergence statement, let
$C_\lambda=\Sigma_{\lambda,II}$ and let $V_m$ be nested finite-dimensional
subspaces whose union is dense in $H_{C_\lambda}=\Range(C_\lambda^{1/2})$.
Then
\[
  R_\lambda^{[m]}
  =
  \sup_{\substack{u\in V_m\\ \langle u,C_\lambda u\rangle>0}}
  \frac{|\langle u,\Sigma_{\lambda,Ik}\rangle|^2}
       {\langle u,C_\lambda u\rangle}
\]
is monotone increasing in $m$ and bounded above by $R_\lambda$.  Density in the
Cameron--Martin energy norm gives $R_\lambda^{[m]}\uparrow R_\lambda$.  In an
eigenbasis of $C_\lambda$, this is the Rayleigh--Ritz identity
\[
  R_\lambda^{[m]}
  =
  \sum_{j=1}^m
  \frac{|\langle\Sigma_{\lambda,Ik},e_j\rangle|^2}{\sigma_j}
  \uparrow
  \|\Sigma_{\lambda,Ik}\|_{H_{C_\lambda}}^2.
\]
\end{proof}

The obstruction to direct inversion is not numerical bookkeeping.  In infinite
dimension the compact covariance block has no bounded inverse on the ambient
Hilbert space.  In finite-dimensional Galerkin matrices the condition number
$\kappa(D_N)$ diverges as the smallest covariance eigenvalue tends to zero, so
the formula $B_N^*D_N^{-1}B_N$ amplifies truncation and roundoff errors.  Even
with exact arithmetic, the limits do not commute: fixing $N$ and then sending
$\lambda\to\infty$ resolves a different finite matrix problem than first
passing to the trace-class covariance and then taking the calibrated shorting
limit.  The stable object is therefore the nested variational short, not direct
matrix inversion.

\section{Gaussian Radon Conditioning and the Fredholm Finite Part}
\label{app:gaussian_fredholm}

The Gaussian and Fredholm layer is the Radon representation of the shorted
residual constructed above.  It is not a second mechanism: disintegration
supplies the exact-evidence law, while the shorted covariance residual supplies
the finite conditional variance used by the Gaussian entropy calculation.

\subsection{Pointwise disintegration}

Let $\mu_\lambda=\mathcal{N}(m_\lambda,\Sigma_\lambda)$ be a Radon Gaussian
measure on $\mathcal{H}=\mathcal{H}_k\oplus\mathcal{H}_I$, and let
$\pi_k(X)=\langle X,\phi_k\rangle$.

\begin{lemma}[Continuous disintegration and exact target zeros]
\label{lem:disintegration}
For each $t\in\mathbb{R}$, there is a unique weakly continuous Gaussian
disintegration member
$\mu_\lambda^t=\mathcal{N}(m_\lambda^t,\Sigma_\lambda^t)$ supported on
$X_k=t$, where
\[
  m_\lambda^t=
  \begin{pmatrix}
    t\\
    m_{\lambda,I}
    +\Sigma_{\lambda,Ik}\Sigma_{\lambda,kk}^{-1}(t-m_{\lambda,k})
  \end{pmatrix},
\]
and
\[
  \Sigma_\lambda^t=
  \begin{pmatrix}
    0&0\\
    0&\Sigma_{\lambda,II}
      -\Sigma_{\lambda,Ik}\Sigma_{\lambda,kk}^{-1}\Sigma_{\lambda,kI}
  \end{pmatrix}.
\]
At $t=x_k^*$, the observed law $\mu_{\mathrm{obs}}=\mu_\lambda^{x_k^*}$
satisfies
\[
  \mathbb{E}_{\mathrm{obs}}[X_k]-x_k^*=0,
  \qquad
  \operatorname{Var}_{\mathrm{obs}}(X_k)=0.
\]
\end{lemma}

\begin{proof}
Regular conditional probabilities exist because $\mathcal{H}$ is Polish.  The
displayed family is the weakly continuous Gaussian disintegration along the
continuous scalar observation map, and this continuous version is unique
\cite{bogachev1998gaussian,lagatta2013continuous}.  Evaluating at $x_k^*$ fixes
the target coordinate almost surely.
\end{proof}

\subsection{Gaussian energy projection and entropy}

Let $C_\lambda=\Sigma_{\lambda,II}$.  Since $C_\lambda$ is trace-class, its
inverse is not a bounded ambient operator; the relevant object is the
Cameron--Martin space $H_{C_\lambda}=\Range(C_\lambda^{1/2})$.

\begin{lemma}[Measurable Cameron--Martin energy projection]
\label{lem:conditional_predictor}
The cross covariance belongs to $H_{C_\lambda}$ and
\[
  \|\Sigma_{\lambda,Ik}\|_{H_{C_\lambda}}
  \le |\Sigma_{\lambda,Ik}|_{Q_0}=O(\lambda^{-3}).
\]
The conditional mean projection is the measurable Wiener functional
\[
  \mu_{k|I}(X_I)
  =
  m_{\lambda,k}
  +W_{\Sigma_{\lambda,Ik}}(X_I-m_{\lambda,I}),
\]
and, for all sufficiently large $\lambda$,
\[
  \operatorname{Var}(X_k\mid X_I)
  =
  \Sigma_{\lambda,kk}
  -\|\Sigma_{\lambda,Ik}\|_{H_{C_\lambda}}^2
  =S_{k,\lambda}>0.
\]
\end{lemma}

\begin{proof}
The domination $\Sigma_{\lambda,II}\succeq Q_0$ gives
$H_{Q_0}\subseteq H_{C_\lambda}$ by Douglas range inclusion
\cite{douglas1966majorization}.  Combining this inclusion with
\eqref{eq:cross_energy} gives the norm bound.  The range condition is exactly
the condition for the measurable Wiener functional associated with
$\Sigma_{\lambda,Ik}$ to exist under the bulk Gaussian law
\cite{bogachev1998gaussian}.  Positivity follows from
\Cref{prop:rank_one_shorted_residual}.
\end{proof}

\begin{lemma}[Feldman--Hajek equivalence and trace-anomaly decay]
\label{lem:feldman_hajek}
There is $\lambda_0>0$ such that for $\lambda\ge\lambda_0$ the bulk marginals
$\mu_{\mathrm{obs},I}$ and $\mu_{\lambda,I}$ are equivalent and
\[
  2\KL(\mu_{\mathrm{obs},I}\parallel\mu_{\lambda,I})
  =
  \frac{(x_k^*-m_{\lambda,k})^2}{\Sigma_{\lambda,kk}}\rho_\lambda
  -\log(1-\rho_\lambda)-\rho_\lambda,
\]
where
\[
  \rho_\lambda
  =
  \frac{\|\Sigma_{\lambda,Ik}\|_{H_{C_\lambda}}^2}
       {\Sigma_{\lambda,kk}}
  =
  1-\frac{S_{k,\lambda}}{\Sigma_{\lambda,kk}}
  =O(\lambda^{-4}).
\]
In particular,
$\KL(\mu_{\mathrm{obs},I}\parallel\mu_{\lambda,I})=O(\lambda^{-4})$.
\end{lemma}

\begin{proof}
The covariance change under disintegration is the rank-one operator
$\Sigma_{\lambda,Ik}\Sigma_{\lambda,kk}^{-1}\Sigma_{\lambda,kI}$.  By
\Cref{lem:conditional_predictor}, its relative covariance perturbation has the
single non-unit eigenvalue $1-\rho_\lambda$.  For large $\lambda$,
$\rho_\lambda<1$, so Feldman--Hajek equivalence applies
\cite{feldman1958equivalence,bogachev1998gaussian}.  The Gaussian relative
entropy formula gives the displayed expression.  Assumption~\ref{assum:gaussian_clamp}
has $m_{\lambda,k}-x_k^*=O(\lambda^{-1})$, while
$\Sigma_{\lambda,kk}=O(\lambda^{-2})$ and $\rho_\lambda=O(\lambda^{-4})$; the
logarithmic remainder is $O(\rho_\lambda^2)$.  Hence the entropy is
$O(\lambda^{-4})$.
\end{proof}

\subsection{Mollified divergence and finite part}

For $\delta>0$, set
\[
  \nu_{\lambda,\delta}
  =
  \mu_{\mathrm{obs},I}\otimes\mathcal{N}(x_k^*,\delta^2)
\]
under the product identification
$\mathcal{H}=\mathcal{H}_I\times\mathcal{H}_k$.

\begin{proposition}[Mollification limit]
\label{prop:mollification_limit}
For $\lambda\ge\lambda_0$ as in \Cref{lem:feldman_hajek}, define
\[
  C_{\lambda,\delta}
  =
  \frac12\left(\log\frac{S_{k,\lambda}}{\delta^2}-1\right).
\]
Then
\[
  \lim_{\delta\downarrow0}
  \left[
    \KL(\nu_{\lambda,\delta}\parallel\mu_\lambda)
    -C_{\lambda,\delta}
  \right]
  =
  \mathfrak{J}_\lambda(\mu_{\mathrm{obs}};\mu_\lambda).
\]
\end{proposition}

\begin{proof}
Let $r_\lambda(X)=X_k-\mu_{k|I}(X_I)$.  The entropy chain rule and the scalar
Gaussian formula give
\[
  \KL(\nu_{\lambda,\delta}\parallel\mu_\lambda)
  =
  \KL(\mu_{\mathrm{obs},I}\parallel\mu_{\lambda,I})
  +C_{\lambda,\delta}
  +\frac{\delta^2}{2S_{k,\lambda}}
  +\frac{1}{2S_{k,\lambda}}
   \mathbb{E}_{\mu_{\mathrm{obs}}}[r_\lambda(X)^2].
    \]
    The $\delta^2$ term vanishes as $\delta\downarrow0$, and the remaining terms
    are exactly Definition~\ref{def:unified_action}.
\end{proof}

\begin{lemma}[Weak scheme-independence]
\label{lem:weak_scheme_independence}
For an evidence level $t\in\mathbb{R}$, let $\mu_\lambda^t$ be the weakly continuous
Gaussian disintegration from \Cref{lem:disintegration} and define
\[
\mathfrak{J}_\lambda(t)\coloneqq
\mathfrak{J}_\lambda(\mu_\lambda^t;\mu_\lambda).
\]
For any two evidence levels $t_1,t_2$ for which the finite part is defined,
\begin{equation}
    \mathfrak{J}_\lambda(t_2)-\mathfrak{J}_\lambda(t_1)
    =
    \frac{(t_2-m_{\lambda,k})^2-(t_1-m_{\lambda,k})^2}
    {2\Sigma_{\lambda,kk}}.
\end{equation}
\begin{proof}
Let $\Delta_t\coloneqq t-m_{\lambda,k}$ and
\[
    \rho_\lambda
    \coloneqq
    \frac{\|\Sigma_{\lambda,Ik}\|_{H_{C_\lambda}}^2}{\Sigma_{\lambda,kk}}
    =
    1-\frac{S_{k,\lambda}}{\Sigma_{\lambda,kk}}.
\]
The Gaussian entropy computation in \Cref{lem:feldman_hajek} gives
\[
    \KL(\mu_{\lambda,I}^t\parallel\mu_{\lambda,I})
    =
    \frac{1}{2}\frac{\Delta_t^2}{\Sigma_{\lambda,kk}}\rho_\lambda
    -\frac{1}{2}\log(1-\rho_\lambda)-\frac{1}{2}\rho_\lambda .
\]
The conditional-residual term in \eqref{eq:unified_action} is
\[
    \frac{1}{2S_{k,\lambda}}
    \E_{\mu_\lambda^t}\!\left[
    (X_k-\mu_{k|I}(X_I))^2
    \right]
    =
    \frac{1}{2}\rho_\lambda
    +
    \frac{1}{2}\frac{\Delta_t^2}{\Sigma_{\lambda,kk}}(1-\rho_\lambda).
\]
Adding these two identities cancels the $\rho_\lambda/2$ terms and yields
\[
    \mathfrak{J}_\lambda(t)
    =
    \frac{(t-m_{\lambda,k})^2}{2\Sigma_{\lambda,kk}}
    -\frac{1}{2}\log(1-\rho_\lambda).
\]
The logarithmic term depends on $C_\lambda$ and $\lambda$, but not on $t$. Hence it
 cancels in the difference between $t_2$ and $t_1$, proving
 \eqref{eq:weak_scheme_independence}. Replacing the Gaussian Dirac mollifier in
 \Cref{prop:mollification_limit} by any centered approximate identity with finite
 second moment changes the extracted finite part only by an additive term depending on
 the kernel, $\lambda$, and $S_{k,\lambda}$, not on $t$. Therefore all evidence
 differences $\Delta\mathfrak{J}_\lambda$ are kernel-independent.
\end{proof}
\end{lemma}

\begin{proof}[Proof of Theorem~\ref{thm:finite_part_regularization}]
Put $\Delta_\lambda=x_k^*-m_{\lambda,k}$ and
$Z_\lambda=W_{\Sigma_{\lambda,Ik}}(X_I-m_{\lambda,I})$.  The Gaussian
conditional formulas give
\[
  \mathbb{E}_{\mathrm{obs}}[Z_\lambda]
  =\rho_\lambda\Delta_\lambda,
  \qquad
  \operatorname{Var}_{\mathrm{obs}}(Z_\lambda)
  =
  \Sigma_{\lambda,kk}\rho_\lambda(1-\rho_\lambda).
\]
Since $X_k=x_k^*$ under $\mu_{\mathrm{obs}}$ and
$S_{k,\lambda}=\Sigma_{\lambda,kk}(1-\rho_\lambda)$,
\[
  \frac{1}{2S_{k,\lambda}}
  \mathbb{E}_{\mu_{\mathrm{obs}}}
  \bigl[(X_k-\mu_{k|I}(X_I))^2\bigr]
  =
  \frac{\rho_\lambda}{2}
  +
  \frac{\Delta_\lambda^2}{2\Sigma_{\lambda,kk}}(1-\rho_\lambda).
\]
The calibrated identities yield the relation
$\Delta_\lambda^2 / (2\Sigma_{\lambda,kk}) = b_k^2 / \varepsilon_0$
as well as $\rho_\lambda = O(\lambda^{-4})$.  Adding
\Cref{lem:feldman_hajek} yields
\[
  \mathfrak{J}_\lambda(\mu_{\mathrm{obs}};\mu_\lambda)
  =
  \frac{b_k^2}{\varepsilon_0}+O(\lambda^{-4}).
\]
\end{proof}

If $a_{Ik}=0$ and $x_k^*=m_{\lambda,k}$, then
$\Sigma_{\lambda,Ik}=0$, the bulk marginals coincide, and the conditional
residual term in Definition~\ref{def:unified_action} is zero.  Hence
$\mathfrak{J}_\lambda(\mu_{\mathrm{obs}};\mu_\lambda)=0$ exactly for every
$\lambda>0$.

\section{Admissibility Boundary and Elliptic Benchmark}
\label{app:admissibility_boundary}

The finite-energy assumption is a boundary condition for the theory, not a
technical artifact.  The off-diagonal channel must land in the Cameron--Martin
range of the background covariance.  If it does not, the variational subtraction
need not be finite and the $O(\lambda^{-6})$ rank-one residual estimate has no
meaning.

The elliptic example of \Cref{sec:spde} makes this boundary concrete.  For
$L=-d^2/dx^2+\kappa^2$ on $L^2(0,1)$ with Dirichlet boundary conditions and
$Q_0=L^{-1}$, the Cameron--Martin space is
\[
  H_{Q_0}=H_0^1(0,1),
  \qquad
  \|h\|_{Q_0}^2
  =
  \int_0^1\bigl(|h'(x)|^2+\kappa^2|h(x)|^2\bigr)\,dx .
\]
Smooth spatially averaged observations are admissible when their representing
vectors lie in $H_0^1(0,1)$; for example, any fixed
$g\in C_c^\infty(0,1)$ has finite energy.  A point-evaluation channel is the
opposite limit.  If $g_\varepsilon$ is a standard compactly supported mollifier
concentrating at $x_0\in(0,1)$, then
\[
  \|g_\varepsilon\|_{Q_0}^2
  =
  \int_0^1\bigl(|g_\varepsilon'(x)|^2+\kappa^2|g_\varepsilon(x)|^2\bigr)\,dx
  \longrightarrow \infty
  \qquad(\varepsilon\downarrow0).
\]
Thus each fixed smoothed sensor can satisfy the finite-energy condition, but
the Dirac point-evaluation limit does not.  The shorting theorem is therefore a
theory of admissible Cameron--Martin channels, not of unsmoothed point sensors
outside the covariance-energy domain.

\end{document}
